\documentclass[a4paper, 12pt]{article}
\usepackage[margin=2cm]{geometry}
\usepackage{color}
\usepackage{amsmath,bm}
\usepackage{graphicx}
\usepackage{booktabs,multirow}
\usepackage{float}
\usepackage{authblk}
\usepackage{pgffor}
\usepackage[no-weekday]{eukdate}

\renewcommand{\thefigure}{S\arabic{figure}}
\renewcommand{\thetable}{S\arabic{table}}

\DeclareMathOperator{\RelRMSE}{RelRMSE}
\DeclareMathOperator{\Base}{Base}
\DeclareMathOperator{\Uni}{Uni}

\title{Multivariate reconciliation for hierarchical time series}

\author[1,2]{Ana Caroline Pinheiro}
\author[1]{Rodrigo de Souza Bulh{\~o}es}
\author[3]{Rob J. Hyndman}
\author[1,4]{Paulo Canas Rodrigues}

%% Author affiliations
\affil[1]{Department of Statistics, Federal University of Bahia, Salvador, Brazil}
\affil[2]{Secretary of Health of the State of Bahia, Salvador, Brazil}
\affil[3]{Department of Econometrics \& Business Statistics, Monash University, Melbourne, Victoria, Australia}
\affil[4]{Department of Business Management, University of Pretoria, Pretoria, South Africa}

\begin{document}

\maketitle
\newpage

%%%% APENAS MATERIAL SUPLEMENTAR(IMGS, FIGURAS, ETC) %%%%

\section*{Figures}

\foreach \i in {1,...,9}{%
  \begin{figure}[H]
    \centering
    \includegraphics[width=\textwidth]{Imagens/\i.pdf}
    \caption{Simulation of a hierarchical multivariate time series, with the same structure as Figure 1, with $m=2$, under Scenario \i.}
    \label{fig:sim\i}
  \end{figure}
}

\newpage

\section*{Tables}

\begin{table}[!htb]
\caption{$\RelRMSE_{\Base}$ for all series in the hierarchy across different forecast horizons. ARIMA, ETS, and VAR models for base forecasts. Using the shrinkage approach to estimate $\boldsymbol{W}$ under scenario 1. Values in red indicate a $\RelRMSE_{\Base}$ less than 0.}
\centering\resizebox{1\textwidth}{!}{%
\begin{tabular}{llcccccccccccc}\hline
\multirow{2}{*}{Series} & \multirow{2}{*}{Variable}
& \multicolumn{12}{c}{Forecast horizons ($h$)} \\ \cline{3-14}
& & 1 & 2 & 3 & 4 & 5 & 6 & 7 & 8 & 9 & 10 & 11 & 12 \\ \hline
ARIMA &  &  &  &  &  &  &  &  &  &  &  &  & \\
\hline
Total & 1 & 0.048 & 0.045 & 0.043 & 0.035 & 0.036 & 0.036 & 0.041 & 0.045 & 0.049 & 0.047 & 0.055 & 0.054\\
Total & 2 & 0.040 & 0.040 & 0.021 & 0.018 & 0.024 & 0.030 & 0.029 & 0.016 & 0.026 & 0.029 & 0.035 & 0.020\\
A & 1 & 0.020 & 0.025 & 0.017 & 0.009 & 0.016 & 0.019 & 0.017 & 0.022 & 0.028 & 0.018 & 0.020 & 0.036\\
A & 2 & 0.029 & 0.021 & 0.033 & 0.023 & 0.030 & 0.024 & 0.022 & 0.024 & 0.027 & 0.026 & 0.032 & 0.031\\
B & 1 & 0.030 & 0.032 & 0.023 & 0.018 & 0.016 & 0.033 & 0.031 & 0.029 & 0.029 & 0.030 & 0.024 & 0.035\\
B & 2 & 0.025 & 0.028 & 0.024 & 0.022 & 0.038 & 0.029 & 0.021 & 0.027 & 0.036 & 0.035 & 0.033 & 0.036\\
AA & 1 & 0.014 & 0.002 & 0.004 & 0.003 & 0.004 & 0.004 & 0.006 & 0.004 & 0.008 & 0.008 & 0.015 & 0.009\\
AA & 2 & 0.020 & 0.010 & 0.012 & 0.008 & 0.011 & 0.005 & 0.016 & 0.012 & 0.017 & 0.012 & 0.020 & 0.014\\
AB & 1 & 0.026 & 0.013 & 0.019 & 0.012 & 0.020 & 0.014 & 0.016 & 0.019 & 0.019 & 0.015 & 0.017 & 0.008\\
AB & 2 & 0.016 & 0.015 & 0.014 & 0.016 & 0.024 & 0.020 & 0.016 & 0.010 & 0.013 & 0.006 & 0.007 & 0.008\\
BA & 1 & 0.016 & 0.006 & 0.010 & 0.009 & 0.009 & 0.011 & 0.013 & 0.008 & 0.013 & 0.009 & 0.010 & 0.007\\
BA & 2 & 0.013 & 0.011 & 0.006 & 0.003 & 0.007 & \textcolor{red}{-0.001} & 0.010 & 0.008 & 0.010 & 0.008 & 0.003 & 0.005\\
BB & 1 & 0.010 & 0.010 & 0.008 & 0.003 & 0.014 & 0.013 & 0.022 & 0.015 & 0.024 & 0.022 & 0.023 & 0.019\\
BB & 2 & 0.014 & 0.014 & 0.017 & 0.007 & 0.008 & 0.016 & 0.015 & 0.011 & 0.012 & 0.018 & 0.016 & 0.012\\
BC & 1 & 0.024 & 0.012 & 0.018 & 0.019 & 0.029 & 0.022 & 0.019 & 0.012 & 0.020 & 0.010 & 0.020 & 0.013\\
BC & 2 & 0.009 & 0.007 & 0.004 & 0.003 & \textcolor{red}{-0.003} & \textcolor{red}{-0.000} & 0.006 & 0.007 & \textcolor{red}{-0.000} & \textcolor{red}{-0.005} & 0.008 & 0.005\\
\hline
ETS &  &  &  &  &  &  &  &  &  &  &  &  & \\
\hline
Total & 1 & 0.003 & 0.002 & 0.003 & 0.002 & 0.004 & 0.002 & 0.003 & 0.001 & 0.003 & 0.003 & 0.002 & 0.001\\
Total & 2 & 0.005 & 0.008 & 0.004 & 0.004 & 0.004 & 0.002 & 0.002 & 0.002 & 0.002 & 0.002 & 0.002 & 0.002\\
A & 1 & 0.010 & 0.004 & 0.006 & 0.003 & 0.004 & 0.002 & 0.003 & 0.002 & 0.003 & 0.002 & 0.003 & 0.002\\
A & 2 & 0.000 & \textcolor{red}{-0.001} & 0.001 & \textcolor{red}{-0.001} & 0.000 & 0.002 & 0.000 & 0.000 & \textcolor{red}{-0.000} & 0.001 & 0.000 & 0.000\\
B & 1 & 0.009 & 0.007 & 0.006 & 0.005 & 0.005 & 0.004 & 0.003 & 0.002 & 0.002 & 0.002 & 0.002 & 0.002\\
B & 2 & 0.007 & 0.005 & 0.003 & 0.000 & 0.002 & 0.004 & 0.004 & 0.004 & 0.006 & 0.006 & 0.005 & 0.004\\
AA & 1 & \textcolor{red}{-0.000} & 0.003 & 0.000 & 0.001 & 0.003 & 0.003 & \textcolor{red}{-0.000} & 0.000 & 0.002 & 0.002 & \textcolor{red}{-0.001} & 0.001\\
AA & 2 & 0.003 & 0.001 & 0.002 & 0.001 & 0.002 & 0.001 & 0.001 & 0.000 & 0.001 & 0.001 & 0.000 & 0.000\\
AB & 1 & 0.001 & 0.000 & 0.000 & 0.001 & 0.002 & 0.000 & \textcolor{red}{-0.001} & 0.000 & 0.000 & \textcolor{red}{-0.000} & \textcolor{red}{-0.001} & 0.000\\
AB & 2 & 0.007 & 0.003 & 0.000 & 0.002 & 0.001 & 0.000 & 0.001 & \textcolor{red}{-0.000} & 0.001 & \textcolor{red}{-0.000} & \textcolor{red}{-0.000} & \textcolor{red}{-0.000}\\
BA & 1 & 0.003 & 0.001 & \textcolor{red}{-0.001} & \textcolor{red}{-0.000} & \textcolor{red}{-0.001} & \textcolor{red}{-0.001} & 0.000 & 0.001 & \textcolor{red}{-0.000} & \textcolor{red}{-0.000} & 0.000 & 0.000\\
BA & 2 & 0.006 & 0.003 & 0.003 & 0.005 & 0.005 & 0.002 & 0.003 & 0.002 & 0.002 & 0.001 & 0.001 & 0.001\\
BB & 1 & 0.001 & 0.002 & 0.001 & \textcolor{red}{-0.000} & \textcolor{red}{-0.000} & \textcolor{red}{-0.000} & \textcolor{red}{-0.000} & 0.001 & \textcolor{red}{-0.000} & \textcolor{red}{-0.001} & \textcolor{red}{-0.000} & 0.000\\
BB & 2 & 0.003 & 0.000 & 0.001 & 0.001 & 0.001 & 0.002 & \textcolor{red}{-0.000} & \textcolor{red}{-0.001} & \textcolor{red}{-0.001} & \textcolor{red}{-0.002} & \textcolor{red}{-0.002} & \textcolor{red}{-0.001}\\
BC & 1 & 0.001 & \textcolor{red}{-0.000} & \textcolor{red}{-0.001} & \textcolor{red}{-0.001} & \textcolor{red}{-0.001} & \textcolor{red}{-0.001} & \textcolor{red}{-0.001} & 0.000 & 0.000 & \textcolor{red}{-0.001} & \textcolor{red}{-0.001} & \textcolor{red}{-0.001}\\
BC & 2 & \textcolor{red}{-0.000} & 0.000 & \textcolor{red}{-0.002} & 0.000 & \textcolor{red}{-0.001} & \textcolor{red}{-0.002} & \textcolor{red}{-0.002} & \textcolor{red}{-0.001} & \textcolor{red}{-0.001} & \textcolor{red}{-0.000} & \textcolor{red}{-0.000} & \textcolor{red}{-0.001}\\
\hline
VAR &  &  &  &  &  &  &  &  &  &  &  &  & \\
\hline
Total & 1 & 0.022 & 0.036 & 0.020 & 0.024 & 0.013 & 0.035 & 0.016 & 0.031 & 0.003 & 0.040 & 0.002 & 0.025\\
Total & 2 & 0.019 & 0.034 & 0.015 & 0.021 & 0.015 & 0.035 & 0.014 & 0.017 & 0.006 & 0.039 & \textcolor{red}{-0.003} & 0.023\\
A & 1 & 0.023 & 0.027 & 0.013 & 0.006 & 0.015 & 0.017 & 0.015 & 0.019 & 0.022 & 0.018 & 0.013 & 0.015\\
A & 2 & 0.027 & 0.028 & 0.028 & 0.025 & 0.028 & 0.030 & 0.031 & 0.021 & 0.020 & 0.023 & 0.027 & 0.021\\
B & 1 & 0.027 & 0.030 & 0.019 & 0.018 & 0.025 & 0.036 & 0.022 & 0.021 & 0.019 & 0.019 & 0.007 & 0.017\\
B & 2 & 0.018 & 0.026 & 0.017 & 0.015 & 0.012 & 0.024 & 0.005 & 0.013 & 0.017 & 0.021 & 0.005 & 0.013\\
AA & 1 & 0.008 & \textcolor{red}{-0.002} & 0.002 & \textcolor{red}{-0.000} & 0.010 & \textcolor{red}{-0.004} & 0.000 & 0.001 & 0.019 & 0.002 & 0.015 & 0.003\\
AA & 2 & 0.009 & \textcolor{red}{-0.007} & 0.002 & \textcolor{red}{-0.007} & 0.006 & \textcolor{red}{-0.005} & 0.014 & 0.010 & 0.021 & 0.001 & 0.015 & 0.004\\
AB & 1 & 0.008 & \textcolor{red}{-0.000} & 0.007 & 0.001 & 0.016 & 0.004 & 0.015 & \textcolor{red}{-0.000} & 0.012 & 0.000 & 0.015 & \textcolor{red}{-0.001}\\
AB & 2 & 0.010 & 0.001 & 0.004 & 0.002 & 0.017 & 0.009 & 0.013 & \textcolor{red}{-0.002} & 0.015 & 0.003 & 0.017 & 0.001\\
BA & 1 & 0.008 & \textcolor{red}{-0.008} & 0.008 & 0.004 & 0.009 & 0.002 & 0.013 & 0.001 & 0.011 & \textcolor{red}{-0.001} & 0.016 & 0.007\\
BA & 2 & 0.014 & 0.004 & 0.007 & 0.010 & 0.020 & 0.009 & 0.020 & 0.012 & 0.021 & 0.003 & 0.013 & 0.003\\
BB & 1 & 0.009 & 0.001 & 0.006 & 0.005 & 0.011 & 0.004 & 0.012 & 0.004 & 0.023 & 0.000 & 0.018 & 0.003\\
BB & 2 & 0.003 & \textcolor{red}{-0.004} & 0.009 & \textcolor{red}{-0.002} & 0.007 & \textcolor{red}{-0.003} & 0.010 & \textcolor{red}{-0.001} & 0.013 & \textcolor{red}{-0.006} & 0.010 & 0.000\\
BC & 1 & 0.006 & \textcolor{red}{-0.007} & 0.011 & 0.005 & 0.007 & 0.001 & 0.023 & 0.005 & 0.011 & 0.006 & 0.024 & 0.002\\
BC & 2 & 0.005 & 0.001 & 0.006 & \textcolor{red}{-0.001} & 0.003 & \textcolor{red}{-0.000} & 0.018 & 0.007 & 0.019 & 0.011 & 0.022 & 0.002\\
\hline
    \end{tabular}%
    }
\end{table}

\begin{table}[!htb]
\caption{$\RelRMSE_{\Base}$ for all series in the hierarchy across different forecast horizons. ARIMA, ETS, and VAR models for base forecasts. Using the shrinkage approach to estimate $\boldsymbol{W}$ under scenario 2. Values in red indicate a $\RelRMSE_{\Base}$ less than 0.}
\centering\resizebox{1\textwidth}{!}{%
\begin{tabular}{llcccccccccccc}\hline
\multirow{2}{*}{Series} & \multirow{2}{*}{Variable}
& \multicolumn{12}{c}{Forecast horizons ($h$)} \\ \cline{3-14}
& & 1 & 2 & 3 & 4 & 5 & 6 & 7 & 8 & 9 & 10 & 11 & 12 \\ \hline
ARIMA &  &  &  &  &  &  &  &  &  &  &  &  & \\
\hline
Total & 1 & 0.029 & 0.032 & 0.031 & 0.017 & 0.036 & 0.040 & 0.050 & 0.039 & 0.040 & 0.061 & 0.064 & 0.064\\
Total & 2 & 0.033 & 0.028 & 0.029 & 0.022 & 0.026 & 0.028 & 0.036 & 0.035 & 0.044 & 0.055 & 0.057 & 0.054\\
A & 1 & 0.037 & 0.036 & 0.021 & 0.018 & 0.039 & 0.036 & 0.021 & 0.037 & 0.036 & 0.031 & 0.027 & 0.026\\
A & 2 & 0.041 & 0.050 & 0.033 & 0.027 & 0.041 & 0.036 & 0.033 & 0.035 & 0.042 & 0.041 & 0.043 & 0.043\\
B & 1 & 0.034 & 0.030 & 0.015 & 0.013 & 0.036 & 0.034 & 0.032 & 0.044 & 0.048 & 0.055 & 0.046 & 0.055\\
B & 2 & 0.025 & 0.039 & 0.028 & 0.023 & 0.024 & 0.043 & 0.047 & 0.043 & 0.041 & 0.054 & 0.050 & 0.048\\
AA & 1 & 0.022 & 0.019 & 0.019 & 0.007 & 0.014 & 0.011 & 0.030 & 0.026 & 0.031 & 0.036 & 0.040 & 0.033\\
AA & 2 & 0.018 & 0.023 & 0.024 & 0.019 & 0.016 & 0.018 & 0.029 & 0.026 & 0.029 & 0.029 & 0.033 & 0.031\\
AB & 1 & 0.023 & 0.018 & 0.024 & 0.017 & 0.021 & 0.018 & 0.021 & 0.018 & 0.017 & 0.016 & 0.021 & 0.015\\
AB & 2 & 0.022 & 0.016 & 0.016 & 0.012 & 0.015 & 0.011 & 0.018 & 0.016 & 0.012 & 0.017 & 0.020 & 0.025\\
BA & 1 & 0.011 & 0.013 & 0.014 & 0.010 & 0.008 & 0.006 & 0.010 & 0.011 & 0.014 & 0.008 & 0.005 & 0.003\\
BA & 2 & 0.015 & 0.011 & 0.005 & 0.002 & 0.012 & 0.011 & 0.011 & 0.011 & 0.017 & 0.019 & 0.008 & 0.012\\
BB & 1 & 0.013 & 0.008 & 0.011 & 0.011 & 0.020 & 0.020 & 0.023 & 0.013 & 0.020 & 0.017 & 0.024 & 0.025\\
BB & 2 & 0.020 & 0.009 & 0.009 & 0.009 & 0.023 & 0.021 & 0.025 & 0.017 & 0.028 & 0.023 & 0.028 & 0.028\\
BC & 1 & 0.009 & 0.016 & 0.020 & 0.011 & 0.010 & 0.009 & 0.018 & 0.019 & 0.020 & 0.023 & 0.026 & 0.022\\
BC & 2 & 0.027 & 0.016 & 0.013 & 0.007 & 0.024 & 0.017 & 0.023 & 0.022 & 0.027 & 0.029 & 0.024 & 0.025\\
\hline
ETS &  &  &  &  &  &  &  &  &  &  &  &  & \\
\hline
Total & 1 & 0.002 & 0.005 & 0.003 & 0.001 & 0.002 & 0.002 & 0.002 & 0.001 & 0.001 & 0.002 & 0.002 & 0.002\\
Total & 2 & 0.004 & 0.003 & 0.003 & 0.001 & 0.002 & 0.002 & 0.003 & 0.000 & 0.002 & 0.001 & 0.002 & 0.001\\
A & 1 & \textcolor{red}{-0.001} & 0.001 & 0.001 & 0.000 & 0.000 & 0.002 & 0.001 & 0.000 & \textcolor{red}{-0.000} & \textcolor{red}{-0.000} & 0.000 & \textcolor{red}{-0.000}\\
A & 2 & 0.004 & 0.002 & \textcolor{red}{-0.001} & 0.003 & 0.003 & 0.001 & 0.001 & 0.001 & 0.000 & \textcolor{red}{-0.000} & \textcolor{red}{-0.000} & 0.002\\
B & 1 & 0.004 & 0.001 & 0.003 & \textcolor{red}{-0.000} & \textcolor{red}{-0.000} & 0.001 & 0.002 & 0.000 & \textcolor{red}{-0.000} & 0.001 & 0.001 & 0.001\\
B & 2 & \textcolor{red}{-0.003} & \textcolor{red}{-0.001} & \textcolor{red}{-0.002} & \textcolor{red}{-0.003} & \textcolor{red}{-0.003} & \textcolor{red}{-0.001} & \textcolor{red}{-0.001} & \textcolor{red}{-0.002} & \textcolor{red}{-0.002} & \textcolor{red}{-0.000} & \textcolor{red}{-0.001} & \textcolor{red}{-0.002}\\
AA & 1 & 0.003 & \textcolor{red}{-0.001} & \textcolor{red}{-0.000} & \textcolor{red}{-0.000} & 0.001 & \textcolor{red}{-0.001} & \textcolor{red}{-0.001} & \textcolor{red}{-0.001} & \textcolor{red}{-0.000} & \textcolor{red}{-0.001} & \textcolor{red}{-0.001} & \textcolor{red}{-0.001}\\
AA & 2 & 0.002 & 0.003 & 0.001 & \textcolor{red}{-0.000} & 0.000 & \textcolor{red}{-0.000} & \textcolor{red}{-0.000} & 0.000 & 0.000 & \textcolor{red}{-0.000} & \textcolor{red}{-0.000} & 0.000\\
AB & 1 & 0.006 & 0.002 & 0.001 & 0.002 & 0.002 & 0.000 & 0.000 & 0.001 & 0.001 & 0.000 & \textcolor{red}{-0.000} & 0.000\\
AB & 2 & 0.004 & 0.002 & 0.002 & 0.001 & 0.000 & 0.001 & 0.002 & 0.001 & 0.001 & 0.001 & 0.001 & \textcolor{red}{-0.000}\\
BA & 1 & 0.003 & \textcolor{red}{-0.001} & \textcolor{red}{-0.001} & 0.001 & 0.001 & 0.000 & 0.001 & 0.000 & 0.001 & \textcolor{red}{-0.000} & \textcolor{red}{-0.000} & 0.000\\
BA & 2 & 0.002 & \textcolor{red}{-0.001} & \textcolor{red}{-0.001} & \textcolor{red}{-0.001} & \textcolor{red}{-0.001} & \textcolor{red}{-0.001} & \textcolor{red}{-0.001} & \textcolor{red}{-0.000} & \textcolor{red}{-0.001} & \textcolor{red}{-0.002} & \textcolor{red}{-0.002} & \textcolor{red}{-0.001}\\
BB & 1 & 0.002 & 0.001 & \textcolor{red}{-0.001} & \textcolor{red}{-0.001} & \textcolor{red}{-0.000} & \textcolor{red}{-0.000} & \textcolor{red}{-0.001} & \textcolor{red}{-0.001} & 0.000 & \textcolor{red}{-0.001} & \textcolor{red}{-0.001} & \textcolor{red}{-0.001}\\
BB & 2 & 0.002 & 0.001 & 0.000 & 0.000 & \textcolor{red}{-0.000} & \textcolor{red}{-0.000} & 0.000 & 0.000 & 0.000 & \textcolor{red}{-0.000} & 0.000 & \textcolor{red}{-0.000}\\
BC & 1 & \textcolor{red}{-0.000} & 0.002 & 0.000 & \textcolor{red}{-0.000} & 0.001 & 0.001 & 0.001 & 0.000 & 0.000 & 0.001 & 0.001 & \textcolor{red}{-0.000}\\
BC & 2 & 0.005 & 0.000 & 0.003 & 0.002 & 0.003 & 0.000 & 0.000 & 0.001 & 0.001 & 0.001 & 0.001 & 0.000\\
\hline
VAR &  &  &  &  &  &  &  &  &  &  &  &  & \\
\hline
Total & 1 & 0.006 & 0.032 & 0.010 & 0.023 & 0.012 & 0.031 & 0.005 & 0.019 & \textcolor{red}{-0.014} & 0.017 & \textcolor{red}{-0.018} & 0.022\\
Total & 2 & 0.001 & 0.033 & 0.013 & 0.015 & 0.005 & 0.027 & 0.005 & 0.026 & \textcolor{red}{-0.015} & 0.018 & \textcolor{red}{-0.019} & 0.020\\
A & 1 & 0.028 & 0.034 & 0.016 & 0.017 & 0.025 & 0.021 & 0.020 & 0.012 & 0.026 & 0.018 & 0.023 & 0.015\\
A & 2 & 0.039 & 0.038 & 0.017 & 0.008 & 0.021 & 0.016 & 0.016 & 0.010 & 0.025 & 0.015 & 0.020 & 0.012\\
B & 1 & 0.030 & 0.047 & 0.030 & 0.022 & 0.029 & 0.044 & 0.026 & 0.024 & 0.021 & 0.029 & 0.012 & 0.017\\
B & 2 & 0.020 & 0.046 & 0.026 & 0.020 & 0.025 & 0.040 & 0.030 & 0.027 & 0.024 & 0.032 & 0.016 & 0.018\\
AA & 1 & 0.024 & 0.007 & 0.024 & 0.003 & 0.022 & 0.010 & 0.022 & 0.003 & 0.023 & 0.008 & 0.029 & \textcolor{red}{-0.003}\\
AA & 2 & 0.017 & 0.004 & 0.021 & 0.007 & 0.023 & 0.002 & 0.016 & \textcolor{red}{-0.002} & 0.019 & 0.004 & 0.026 & \textcolor{red}{-0.003}\\
AB & 1 & 0.015 & \textcolor{red}{-0.001} & 0.008 & 0.004 & 0.022 & 0.008 & 0.020 & 0.005 & 0.025 & \textcolor{red}{-0.002} & 0.023 & 0.005\\
AB & 2 & 0.013 & \textcolor{red}{-0.009} & 0.007 & \textcolor{red}{-0.001} & 0.016 & 0.001 & 0.020 & 0.004 & 0.021 & \textcolor{red}{-0.004} & 0.023 & 0.005\\
BA & 1 & 0.015 & 0.003 & 0.012 & \textcolor{red}{-0.004} & 0.011 & \textcolor{red}{-0.003} & 0.010 & \textcolor{red}{-0.005} & 0.019 & \textcolor{red}{-0.002} & 0.018 & \textcolor{red}{-0.001}\\
BA & 2 & 0.016 & \textcolor{red}{-0.000} & 0.007 & \textcolor{red}{-0.001} & 0.013 & \textcolor{red}{-0.001} & 0.006 & \textcolor{red}{-0.003} & 0.026 & \textcolor{red}{-0.002} & 0.018 & 0.001\\
BB & 1 & 0.013 & \textcolor{red}{-0.008} & 0.004 & \textcolor{red}{-0.005} & 0.005 & \textcolor{red}{-0.012} & 0.005 & \textcolor{red}{-0.012} & 0.015 & \textcolor{red}{-0.004} & 0.023 & \textcolor{red}{-0.001}\\
BB & 2 & 0.012 & \textcolor{red}{-0.009} & 0.005 & \textcolor{red}{-0.000} & 0.014 & \textcolor{red}{-0.010} & 0.009 & \textcolor{red}{-0.012} & 0.016 & \textcolor{red}{-0.004} & 0.024 & 0.000\\
BC & 1 & 0.007 & \textcolor{red}{-0.008} & 0.005 & \textcolor{red}{-0.002} & 0.006 & \textcolor{red}{-0.011} & 0.011 & \textcolor{red}{-0.002} & 0.021 & 0.001 & 0.026 & 0.004\\
BC & 2 & 0.012 & \textcolor{red}{-0.018} & \textcolor{red}{-0.003} & \textcolor{red}{-0.004} & 0.013 & \textcolor{red}{-0.012} & 0.011 & \textcolor{red}{-0.001} & 0.023 & \textcolor{red}{-0.001} & 0.023 & 0.002\\
\hline
    \end{tabular}%
    }
\end{table}

\begin{table}[!htb]
\caption{$\RelRMSE_{\Base}$ for all series in the hierarchy across different forecast horizons. ARIMA, ETS, and VAR models for base forecasts. Using the shrinkage approach to estimate $\boldsymbol{W}$ under scenario 3. Values in red indicate a $\RelRMSE_{\Base}$ less than 0.}
\centering\resizebox{1\textwidth}{!}{%
\begin{tabular}{llcccccccccccc}\hline
\multirow{2}{*}{Series} & \multirow{2}{*}{Variable}
& \multicolumn{12}{c}{Forecast horizons ($h$)} \\ \cline{3-14}
& & 1 & 2 & 3 & 4 & 5 & 6 & 7 & 8 & 9 & 10 & 11 & 12 \\ \hline
ARIMA &  &  &  &  &  &  &  &  &  &  &  &  & \\
\hline
Total & 1 & 0.027 & 0.042 & 0.034 & 0.036 & 0.040 & 0.035 & 0.031 & 0.038 & 0.039 & 0.052 & 0.051 & 0.054\\
Total & 2 & 0.033 & 0.036 & 0.021 & 0.020 & 0.036 & 0.029 & 0.024 & 0.014 & 0.038 & 0.052 & 0.032 & 0.030\\
A & 1 & 0.029 & 0.029 & 0.022 & 0.012 & 0.039 & 0.026 & 0.032 & 0.033 & 0.031 & 0.034 & 0.028 & 0.038\\
A & 2 & 0.032 & 0.019 & 0.023 & 0.013 & 0.029 & 0.022 & 0.028 & 0.016 & 0.030 & 0.028 & 0.032 & 0.020\\
B & 1 & 0.026 & 0.025 & 0.035 & 0.038 & 0.030 & 0.034 & 0.036 & 0.033 & 0.028 & 0.039 & 0.034 & 0.038\\
B & 2 & 0.038 & 0.041 & 0.011 & 0.019 & 0.022 & 0.030 & 0.013 & 0.024 & 0.017 & 0.031 & 0.025 & 0.030\\
AA & 1 & 0.018 & 0.015 & 0.027 & 0.014 & 0.021 & 0.007 & 0.019 & 0.009 & 0.029 & 0.012 & 0.019 & 0.008\\
AA & 2 & 0.018 & 0.027 & 0.021 & 0.015 & 0.027 & 0.025 & 0.015 & 0.016 & 0.028 & 0.027 & 0.023 & 0.015\\
AB & 1 & 0.013 & 0.009 & 0.020 & \textcolor{red}{-0.000} & 0.014 & 0.008 & 0.029 & 0.004 & 0.014 & 0.010 & 0.020 & 0.000\\
AB & 2 & 0.014 & 0.019 & 0.034 & 0.018 & 0.015 & 0.013 & 0.028 & 0.025 & 0.026 & 0.017 & 0.033 & 0.024\\
BA & 1 & 0.013 & 0.011 & 0.001 & 0.001 & 0.007 & \textcolor{red}{-0.000} & \textcolor{red}{-0.001} & 0.001 & 0.009 & 0.011 & 0.013 & 0.006\\
BA & 2 & 0.024 & 0.018 & 0.015 & 0.004 & 0.019 & 0.010 & 0.012 & 0.008 & 0.016 & 0.006 & 0.009 & 0.013\\
BB & 1 & 0.018 & 0.011 & 0.007 & 0.013 & 0.024 & 0.012 & 0.020 & 0.005 & 0.015 & 0.001 & 0.016 & 0.005\\
BB & 2 & 0.010 & 0.001 & 0.014 & 0.012 & 0.017 & 0.007 & 0.008 & 0.005 & 0.006 & \textcolor{red}{-0.005} & 0.009 & 0.009\\
BC & 1 & 0.020 & 0.005 & 0.019 & 0.008 & 0.028 & 0.013 & 0.019 & 0.022 & 0.037 & 0.014 & 0.026 & 0.022\\
BC & 2 & \textcolor{red}{-0.000} & 0.013 & 0.011 & 0.007 & 0.013 & 0.011 & 0.016 & 0.013 & 0.020 & 0.013 & 0.014 & 0.016\\
\hline
ETS &  &  &  &  &  &  &  &  &  &  &  &  & \\
\hline
Total & 1 & 0.003 & 0.011 & 0.009 & 0.009 & 0.007 & 0.006 & 0.005 & 0.008 & 0.007 & 0.005 & 0.003 & 0.006\\
Total & 2 & 0.008 & 0.010 & 0.008 & 0.008 & 0.008 & 0.007 & 0.007 & 0.010 & 0.008 & 0.007 & 0.008 & 0.006\\
A & 1 & 0.003 & 0.006 & 0.006 & 0.003 & 0.004 & 0.005 & 0.005 & 0.001 & 0.000 & 0.001 & 0.001 & 0.003\\
A & 2 & 0.006 & 0.006 & 0.004 & 0.001 & 0.002 & 0.002 & 0.001 & 0.002 & 0.004 & 0.005 & 0.004 & 0.005\\
B & 1 & 0.013 & 0.008 & 0.008 & 0.005 & 0.006 & 0.006 & 0.004 & 0.006 & 0.004 & 0.006 & 0.005 & 0.006\\
B & 2 & 0.011 & 0.010 & 0.007 & 0.001 & 0.005 & 0.007 & 0.004 & 0.004 & 0.002 & 0.005 & 0.004 & 0.002\\
AA & 1 & 0.006 & 0.006 & 0.008 & 0.003 & 0.006 & 0.005 & 0.003 & 0.004 & 0.005 & 0.001 & 0.002 & 0.002\\
AA & 2 & 0.008 & 0.005 & 0.006 & 0.005 & 0.002 & 0.004 & 0.004 & 0.004 & 0.005 & 0.004 & 0.004 & 0.003\\
AB & 1 & 0.004 & 0.001 & 0.001 & 0.005 & 0.007 & 0.005 & 0.005 & 0.005 & 0.004 & 0.005 & 0.002 & 0.002\\
AB & 2 & 0.007 & 0.006 & 0.007 & 0.004 & 0.006 & 0.004 & 0.004 & 0.004 & 0.002 & 0.000 & 0.001 & 0.001\\
BA & 1 & \textcolor{red}{-0.001} & 0.002 & 0.000 & 0.002 & 0.002 & 0.006 & 0.007 & 0.006 & 0.004 & 0.003 & 0.004 & 0.003\\
BA & 2 & 0.001 & 0.003 & 0.003 & 0.001 & 0.001 & 0.003 & 0.004 & 0.003 & 0.001 & 0.001 & 0.002 & 0.004\\
BB & 1 & 0.002 & 0.005 & 0.005 & 0.002 & 0.001 & 0.000 & 0.003 & 0.002 & 0.003 & \textcolor{red}{-0.000} & 0.000 & 0.000\\
BB & 2 & 0.002 & 0.005 & 0.005 & 0.005 & 0.004 & 0.003 & 0.007 & 0.007 & 0.004 & 0.003 & 0.001 & 0.003\\
BC & 1 & 0.007 & 0.002 & 0.005 & 0.006 & 0.005 & 0.002 & 0.004 & 0.004 & 0.006 & 0.004 & 0.005 & 0.004\\
BC & 2 & \textcolor{red}{-0.003} & \textcolor{red}{-0.001} & 0.004 & 0.003 & 0.003 & 0.001 & 0.001 & 0.002 & 0.004 & 0.004 & 0.003 & 0.003\\
\hline
VAR &  &  &  &  &  &  &  &  &  &  &  &  & \\
\hline
Total & 1 & 0.021 & 0.021 & 0.017 & 0.010 & 0.004 & 0.020 & 0.015 & 0.018 & \textcolor{red}{-0.001} & 0.019 & 0.009 & 0.009\\
Total & 2 & 0.010 & 0.018 & 0.016 & 0.023 & 0.014 & 0.012 & 0.010 & 0.011 & 0.012 & 0.011 & 0.001 & 0.016\\
A & 1 & 0.018 & 0.015 & 0.011 & \textcolor{red}{-0.001} & 0.011 & 0.011 & 0.005 & 0.005 & 0.013 & 0.013 & 0.007 & 0.011\\
A & 2 & 0.017 & 0.013 & 0.010 & 0.003 & 0.008 & 0.019 & 0.013 & 0.013 & 0.014 & 0.021 & 0.008 & 0.004\\
B & 1 & 0.019 & 0.011 & 0.011 & 0.012 & 0.010 & 0.012 & 0.018 & 0.016 & 0.013 & 0.016 & 0.011 & 0.010\\
B & 2 & 0.025 & 0.009 & 0.010 & 0.010 & 0.017 & 0.021 & 0.010 & 0.017 & 0.012 & 0.015 & 0.005 & 0.005\\
AA & 1 & 0.007 & 0.003 & \textcolor{red}{-0.005} & 0.003 & 0.015 & 0.003 & 0.003 & 0.004 & 0.004 & 0.001 & \textcolor{red}{-0.000} & \textcolor{red}{-0.000}\\
AA & 2 & 0.002 & 0.004 & 0.004 & 0.006 & 0.010 & 0.012 & 0.008 & 0.001 & 0.004 & \textcolor{red}{-0.003} & 0.008 & 0.001\\
AB & 1 & \textcolor{red}{-0.007} & \textcolor{red}{-0.009} & 0.002 & \textcolor{red}{-0.000} & \textcolor{red}{-0.003} & \textcolor{red}{-0.002} & \textcolor{red}{-0.003} & \textcolor{red}{-0.002} & 0.001 & \textcolor{red}{-0.002} & 0.007 & 0.006\\
AB & 2 & 0.001 & 0.005 & 0.008 & 0.010 & 0.015 & 0.006 & 0.008 & 0.007 & 0.008 & \textcolor{red}{-0.001} & 0.006 & 0.003\\
BA & 1 & 0.000 & 0.004 & \textcolor{red}{-0.004} & \textcolor{red}{-0.006} & \textcolor{red}{-0.002} & \textcolor{red}{-0.006} & \textcolor{red}{-0.004} & 0.000 & 0.007 & \textcolor{red}{-0.004} & 0.004 & 0.001\\
BA & 2 & \textcolor{red}{-0.008} & \textcolor{red}{-0.000} & \textcolor{red}{-0.006} & \textcolor{red}{-0.008} & 0.001 & \textcolor{red}{-0.003} & 0.004 & \textcolor{red}{-0.001} & 0.005 & \textcolor{red}{-0.004} & 0.001 & 0.002\\
BB & 1 & \textcolor{red}{-0.001} & 0.002 & 0.008 & \textcolor{red}{-0.001} & 0.005 & 0.003 & 0.009 & \textcolor{red}{-0.003} & \textcolor{red}{-0.001} & 0.002 & 0.008 & 0.000\\
BB & 2 & \textcolor{red}{-0.006} & 0.004 & 0.007 & 0.001 & 0.004 & 0.002 & 0.002 & 0.000 & 0.008 & 0.007 & 0.010 & 0.008\\
BC & 1 & 0.003 & \textcolor{red}{-0.011} & \textcolor{red}{-0.001} & 0.004 & 0.007 & 0.002 & 0.004 & 0.004 & 0.014 & 0.001 & 0.003 & 0.006\\
BC & 2 & \textcolor{red}{-0.002} & \textcolor{red}{-0.002} & \textcolor{red}{-0.001} & 0.000 & 0.002 & 0.009 & 0.009 & 0.000 & 0.005 & 0.002 & 0.002 & \textcolor{red}{-0.001}\\
\hline
    \end{tabular}%
    }
\end{table}

\begin{table}[!htb]
\caption{$\RelRMSE_{\Base}$ for all series in the hierarchy across different forecast horizons. ARIMA, ETS, and VAR models for base forecasts. Using the shrinkage approach to estimate $\boldsymbol{W}$ under scenario 4. Values in red indicate a $\RelRMSE_{\Base}$ less than 0.}
\centering\resizebox{1\textwidth}{!}{%
\begin{tabular}{llcccccccccccc}\hline
\multirow{2}{*}{Series} & \multirow{2}{*}{Variable}
& \multicolumn{12}{c}{Forecast horizons ($h$)} \\ \cline{3-14}
& & 1 & 2 & 3 & 4 & 5 & 6 & 7 & 8 & 9 & 10 & 11 & 12 \\ \hline
ARIMA &  &  &  &  &  &  &  &  &  &  &  &  & \\
\hline
Total & 1 & 0.041 & 0.029 & 0.004 & 0.013 & 0.028 & 0.032 & 0.036 & 0.032 & 0.041 & 0.047 & 0.048 & 0.044\\
Total & 2 & 0.040 & 0.021 & 0.019 & 0.029 & 0.030 & 0.041 & 0.036 & 0.035 & 0.058 & 0.058 & 0.067 & 0.061\\
A & 1 & 0.030 & 0.044 & 0.030 & 0.017 & 0.025 & 0.027 & 0.033 & 0.027 & 0.029 & 0.033 & 0.028 & 0.022\\
A & 2 & 0.039 & 0.030 & 0.025 & 0.013 & 0.023 & 0.026 & 0.026 & 0.023 & 0.020 & 0.035 & 0.033 & 0.029\\
B & 1 & 0.053 & 0.042 & 0.025 & 0.033 & 0.041 & 0.040 & 0.034 & 0.034 & 0.034 & 0.035 & 0.036 & 0.039\\
B & 2 & 0.026 & 0.029 & 0.027 & 0.029 & 0.041 & 0.030 & 0.030 & 0.037 & 0.037 & 0.034 & 0.044 & 0.036\\
AA & 1 & 0.013 & 0.009 & 0.012 & 0.015 & 0.018 & 0.011 & 0.010 & 0.013 & 0.023 & 0.016 & 0.017 & 0.015\\
AA & 2 & 0.013 & \textcolor{red}{-0.003} & \textcolor{red}{-0.006} & \textcolor{red}{-0.004} & 0.013 & 0.002 & 0.011 & 0.010 & 0.016 & 0.006 & 0.008 & 0.010\\
AB & 1 & 0.020 & 0.012 & 0.014 & 0.012 & 0.019 & 0.014 & 0.015 & 0.013 & 0.022 & 0.012 & 0.016 & 0.015\\
AB & 2 & 0.009 & 0.001 & \textcolor{red}{-0.001} & 0.000 & 0.005 & 0.006 & 0.017 & 0.012 & 0.010 & 0.006 & 0.014 & 0.013\\
BA & 1 & 0.011 & 0.007 & 0.004 & 0.003 & 0.005 & 0.004 & 0.001 & 0.009 & 0.023 & 0.019 & 0.016 & 0.014\\
BA & 2 & 0.011 & 0.001 & 0.004 & 0.004 & 0.008 & 0.009 & 0.018 & 0.009 & 0.014 & 0.020 & 0.025 & 0.016\\
BB & 1 & 0.001 & 0.008 & 0.001 & 0.006 & 0.008 & 0.012 & 0.007 & 0.011 & 0.014 & 0.019 & 0.018 & 0.014\\
BB & 2 & 0.014 & \textcolor{red}{-0.000} & 0.004 & 0.008 & 0.017 & 0.011 & 0.017 & 0.012 & 0.019 & 0.011 & 0.014 & 0.008\\
BC & 1 & 0.018 & 0.009 & 0.005 & 0.004 & 0.012 & 0.010 & 0.004 & 0.011 & 0.019 & 0.015 & 0.014 & 0.016\\
BC & 2 & 0.015 & 0.007 & 0.009 & 0.009 & 0.013 & 0.009 & 0.016 & 0.012 & 0.018 & 0.010 & 0.018 & 0.013\\
\hline
ETS &  &  &  &  &  &  &  &  &  &  &  &  & \\
\hline
Total & 1 & 0.005 & 0.008 & 0.003 & 0.003 & 0.003 & 0.005 & 0.004 & 0.005 & 0.005 & 0.005 & 0.004 & 0.004\\
Total & 2 & 0.010 & 0.005 & 0.004 & 0.002 & 0.002 & 0.002 & 0.002 & 0.001 & \textcolor{red}{-0.000} & 0.000 & 0.000 & 0.000\\
A & 1 & 0.006 & 0.004 & 0.002 & 0.002 & 0.003 & 0.003 & 0.000 & 0.001 & 0.000 & 0.001 & \textcolor{red}{-0.000} & 0.000\\
A & 2 & 0.006 & 0.002 & 0.001 & 0.001 & 0.003 & 0.002 & 0.001 & 0.000 & 0.001 & 0.000 & 0.000 & 0.000\\
B & 1 & 0.007 & 0.005 & 0.002 & 0.001 & \textcolor{red}{-0.000} & 0.002 & 0.000 & \textcolor{red}{-0.000} & 0.000 & 0.001 & \textcolor{red}{-0.000} & \textcolor{red}{-0.000}\\
B & 2 & 0.002 & 0.003 & 0.001 & 0.001 & 0.003 & 0.002 & 0.002 & 0.001 & 0.004 & 0.003 & 0.003 & 0.001\\
AA & 1 & 0.001 & \textcolor{red}{-0.000} & 0.001 & 0.001 & \textcolor{red}{-0.000} & \textcolor{red}{-0.000} & 0.001 & 0.000 & 0.000 & \textcolor{red}{-0.001} & 0.001 & 0.000\\
AA & 2 & \textcolor{red}{-0.001} & \textcolor{red}{-0.002} & \textcolor{red}{-0.001} & \textcolor{red}{-0.001} & \textcolor{red}{-0.002} & \textcolor{red}{-0.002} & \textcolor{red}{-0.000} & 0.000 & \textcolor{red}{-0.001} & \textcolor{red}{-0.001} & \textcolor{red}{-0.000} & 0.000\\
AB & 1 & \textcolor{red}{-0.001} & \textcolor{red}{-0.000} & 0.000 & 0.000 & 0.001 & \textcolor{red}{-0.000} & 0.000 & 0.000 & \textcolor{red}{-0.001} & \textcolor{red}{-0.002} & \textcolor{red}{-0.001} & \textcolor{red}{-0.001}\\
AB & 2 & \textcolor{red}{-0.002} & 0.000 & 0.001 & \textcolor{red}{-0.000} & \textcolor{red}{-0.001} & \textcolor{red}{-0.000} & 0.001 & 0.001 & \textcolor{red}{-0.001} & \textcolor{red}{-0.000} & 0.000 & 0.001\\
BA & 1 & \textcolor{red}{-0.000} & \textcolor{red}{-0.002} & \textcolor{red}{-0.001} & \textcolor{red}{-0.002} & \textcolor{red}{-0.000} & \textcolor{red}{-0.001} & \textcolor{red}{-0.000} & \textcolor{red}{-0.001} & \textcolor{red}{-0.000} & \textcolor{red}{-0.001} & 0.000 & \textcolor{red}{-0.000}\\
BA & 2 & \textcolor{red}{-0.000} & \textcolor{red}{-0.001} & \textcolor{red}{-0.001} & \textcolor{red}{-0.001} & \textcolor{red}{-0.001} & \textcolor{red}{-0.001} & \textcolor{red}{-0.002} & \textcolor{red}{-0.002} & \textcolor{red}{-0.001} & \textcolor{red}{-0.002} & \textcolor{red}{-0.001} & \textcolor{red}{-0.001}\\
BB & 1 & \textcolor{red}{-0.001} & \textcolor{red}{-0.002} & \textcolor{red}{-0.001} & \textcolor{red}{-0.001} & 0.001 & \textcolor{red}{-0.001} & \textcolor{red}{-0.001} & \textcolor{red}{-0.001} & \textcolor{red}{-0.000} & \textcolor{red}{-0.001} & \textcolor{red}{-0.001} & \textcolor{red}{-0.001}\\
BB & 2 & 0.000 & \textcolor{red}{-0.000} & \textcolor{red}{-0.001} & \textcolor{red}{-0.001} & \textcolor{red}{-0.001} & \textcolor{red}{-0.001} & \textcolor{red}{-0.002} & \textcolor{red}{-0.002} & \textcolor{red}{-0.001} & \textcolor{red}{-0.001} & \textcolor{red}{-0.001} & \textcolor{red}{-0.002}\\
BC & 1 & 0.001 & \textcolor{red}{-0.003} & \textcolor{red}{-0.002} & \textcolor{red}{-0.001} & 0.001 & 0.000 & \textcolor{red}{-0.001} & \textcolor{red}{-0.000} & 0.002 & 0.000 & \textcolor{red}{-0.000} & 0.000\\
BC & 2 & 0.001 & \textcolor{red}{-0.001} & \textcolor{red}{-0.001} & \textcolor{red}{-0.000} & \textcolor{red}{-0.001} & \textcolor{red}{-0.002} & \textcolor{red}{-0.001} & \textcolor{red}{-0.001} & \textcolor{red}{-0.001} & \textcolor{red}{-0.002} & \textcolor{red}{-0.001} & \textcolor{red}{-0.001}\\
\hline
VAR &  &  &  &  &  &  &  &  &  &  &  &  & \\
\hline
Total & 1 & 0.021 & 0.039 & 0.022 & 0.027 & 0.024 & 0.036 & 0.017 & 0.021 & 0.005 & 0.024 & 0.009 & 0.024\\
Total & 2 & 0.026 & 0.031 & 0.018 & 0.017 & 0.026 & 0.041 & 0.025 & 0.031 & 0.004 & 0.032 & 0.010 & 0.017\\
A & 1 & 0.022 & 0.024 & 0.023 & 0.013 & 0.016 & 0.017 & 0.022 & 0.014 & 0.014 & 0.009 & 0.013 & \textcolor{red}{-0.002}\\
A & 2 & 0.020 & 0.026 & 0.018 & 0.017 & 0.020 & 0.020 & 0.023 & 0.012 & 0.017 & 0.012 & 0.019 & 0.007\\
B & 1 & 0.037 & 0.036 & 0.028 & 0.016 & 0.032 & 0.032 & 0.022 & 0.017 & 0.023 & 0.022 & 0.014 & 0.012\\
B & 2 & 0.033 & 0.035 & 0.025 & 0.022 & 0.031 & 0.029 & 0.026 & 0.021 & 0.029 & 0.023 & 0.019 & 0.015\\
AA & 1 & \textcolor{red}{-0.004} & \textcolor{red}{-0.000} & 0.005 & 0.004 & 0.005 & \textcolor{red}{-0.001} & 0.007 & \textcolor{red}{-0.003} & 0.004 & \textcolor{red}{-0.006} & 0.007 & \textcolor{red}{-0.003}\\
AA & 2 & \textcolor{red}{-0.001} & \textcolor{red}{-0.007} & 0.001 & \textcolor{red}{-0.001} & 0.007 & 0.001 & 0.012 & 0.004 & 0.008 & \textcolor{red}{-0.002} & 0.012 & 0.001\\
AB & 1 & 0.000 & \textcolor{red}{-0.001} & 0.005 & \textcolor{red}{-0.001} & \textcolor{red}{-0.000} & \textcolor{red}{-0.005} & 0.003 & \textcolor{red}{-0.003} & 0.002 & \textcolor{red}{-0.008} & 0.001 & \textcolor{red}{-0.008}\\
AB & 2 & 0.000 & \textcolor{red}{-0.005} & 0.003 & \textcolor{red}{-0.002} & 0.001 & \textcolor{red}{-0.004} & 0.002 & 0.001 & 0.009 & \textcolor{red}{-0.004} & 0.006 & \textcolor{red}{-0.001}\\
BA & 1 & \textcolor{red}{-0.004} & \textcolor{red}{-0.003} & \textcolor{red}{-0.001} & \textcolor{red}{-0.006} & 0.000 & \textcolor{red}{-0.002} & \textcolor{red}{-0.000} & \textcolor{red}{-0.003} & 0.005 & \textcolor{red}{-0.003} & 0.002 & \textcolor{red}{-0.003}\\
BA & 2 & 0.003 & \textcolor{red}{-0.003} & \textcolor{red}{-0.003} & \textcolor{red}{-0.004} & 0.002 & \textcolor{red}{-0.003} & 0.001 & \textcolor{red}{-0.002} & 0.005 & \textcolor{red}{-0.004} & 0.004 & \textcolor{red}{-0.001}\\
BB & 1 & 0.003 & 0.000 & \textcolor{red}{-0.004} & \textcolor{red}{-0.006} & 0.003 & \textcolor{red}{-0.003} & \textcolor{red}{-0.005} & \textcolor{red}{-0.006} & 0.006 & \textcolor{red}{-0.003} & \textcolor{red}{-0.001} & \textcolor{red}{-0.004}\\
BB & 2 & 0.005 & \textcolor{red}{-0.006} & \textcolor{red}{-0.003} & 0.000 & 0.004 & \textcolor{red}{-0.007} & \textcolor{red}{-0.003} & \textcolor{red}{-0.002} & 0.008 & \textcolor{red}{-0.002} & 0.001 & \textcolor{red}{-0.002}\\
BC & 1 & 0.002 & \textcolor{red}{-0.007} & \textcolor{red}{-0.006} & \textcolor{red}{-0.005} & 0.002 & \textcolor{red}{-0.005} & \textcolor{red}{-0.000} & 0.001 & 0.008 & \textcolor{red}{-0.003} & 0.002 & \textcolor{red}{-0.001}\\
BC & 2 & 0.007 & \textcolor{red}{-0.003} & \textcolor{red}{-0.004} & \textcolor{red}{-0.004} & 0.000 & \textcolor{red}{-0.008} & \textcolor{red}{-0.004} & \textcolor{red}{-0.001} & 0.005 & \textcolor{red}{-0.005} & \textcolor{red}{-0.004} & \textcolor{red}{-0.001}\\
\hline
    \end{tabular}%
    }
\end{table}

\begin{table}[!htb]
\caption{$\RelRMSE_{\Base}$ for all series in the hierarchy across different forecast horizons. ARIMA, ETS, and VAR models for base forecasts. Using the shrinkage approach to estimate $\boldsymbol{W}$ under scenario 5. Values in red indicate a $\RelRMSE_{\Base}$ less than 0.}
\centering\resizebox{1\textwidth}{!}{%
\begin{tabular}{llcccccccccccc}\hline
\multirow{2}{*}{Series} & \multirow{2}{*}{Variable}
& \multicolumn{12}{c}{Forecast horizons ($h$)} \\ \cline{3-14}
& & 1 & 2 & 3 & 4 & 5 & 6 & 7 & 8 & 9 & 10 & 11 & 12 \\ \hline
ARIMA &  &  &  &  &  &  &  &  &  &  &  &  & \\
\hline
Total & 1 & 0.045 & 0.042 & 0.034 & 0.038 & 0.049 & 0.054 & 0.049 & 0.056 & 0.065 & 0.071 & 0.072 & 0.075\\
Total & 2 & 0.035 & 0.032 & 0.021 & 0.037 & 0.047 & 0.045 & 0.037 & 0.046 & 0.048 & 0.050 & 0.049 & 0.062\\
A & 1 & 0.040 & 0.036 & 0.028 & 0.018 & 0.028 & 0.030 & 0.030 & 0.021 & 0.039 & 0.031 & 0.034 & 0.037\\
A & 2 & 0.031 & 0.036 & 0.033 & 0.022 & 0.020 & 0.024 & 0.024 & 0.017 & 0.023 & 0.026 & 0.023 & 0.022\\
B & 1 & 0.023 & 0.035 & 0.030 & 0.024 & 0.035 & 0.033 & 0.038 & 0.037 & 0.048 & 0.054 & 0.055 & 0.058\\
B & 2 & 0.042 & 0.035 & 0.022 & 0.018 & 0.032 & 0.044 & 0.034 & 0.035 & 0.037 & 0.045 & 0.037 & 0.040\\
AA & 1 & 0.021 & 0.010 & 0.007 & 0.003 & 0.012 & 0.014 & 0.009 & 0.014 & 0.025 & 0.020 & 0.021 & 0.022\\
AA & 2 & 0.017 & 0.012 & 0.013 & 0.010 & 0.011 & 0.010 & 0.016 & 0.014 & 0.022 & 0.014 & 0.016 & 0.013\\
AB & 1 & 0.017 & 0.014 & 0.007 & 0.002 & 0.004 & 0.011 & 0.009 & 0.001 & 0.014 & 0.017 & 0.016 & 0.012\\
AB & 2 & 0.032 & 0.015 & 0.010 & 0.000 & 0.009 & 0.006 & 0.008 & 0.004 & 0.015 & 0.010 & 0.008 & 0.010\\
BA & 1 & 0.004 & 0.013 & 0.011 & 0.007 & 0.010 & 0.012 & 0.016 & 0.018 & 0.023 & 0.025 & 0.027 & 0.018\\
BA & 2 & 0.001 & 0.007 & 0.010 & 0.001 & 0.002 & 0.007 & 0.013 & 0.012 & 0.010 & 0.019 & 0.030 & 0.027\\
BB & 1 & 0.007 & 0.010 & 0.001 & \textcolor{red}{-0.005} & 0.002 & 0.003 & 0.007 & 0.004 & 0.013 & 0.014 & 0.015 & 0.011\\
BB & 2 & 0.011 & 0.004 & 0.002 & 0.003 & 0.011 & 0.006 & 0.007 & 0.013 & 0.015 & 0.018 & 0.021 & 0.019\\
BC & 1 & 0.019 & 0.014 & 0.010 & 0.004 & 0.010 & 0.012 & 0.023 & 0.020 & 0.016 & 0.019 & 0.024 & 0.020\\
BC & 2 & 0.001 & 0.003 & 0.002 & \textcolor{red}{-0.002} & 0.003 & 0.004 & 0.008 & 0.009 & 0.005 & 0.010 & 0.016 & 0.016\\
\hline
ETS &  &  &  &  &  &  &  &  &  &  &  &  & \\
\hline
Total & 1 & 0.005 & 0.002 & 0.001 & 0.001 & 0.001 & 0.001 & 0.000 & 0.001 & 0.001 & 0.000 & 0.000 & 0.001\\
Total & 2 & 0.008 & 0.005 & 0.003 & 0.003 & 0.003 & 0.004 & 0.002 & 0.002 & 0.002 & 0.003 & 0.001 & 0.001\\
A & 1 & 0.003 & \textcolor{red}{-0.001} & \textcolor{red}{-0.000} & \textcolor{red}{-0.001} & \textcolor{red}{-0.000} & 0.000 & 0.001 & \textcolor{red}{-0.000} & 0.001 & 0.000 & 0.001 & 0.001\\
A & 2 & 0.003 & 0.001 & 0.000 & 0.001 & \textcolor{red}{-0.001} & 0.000 & 0.001 & 0.000 & \textcolor{red}{-0.000} & 0.001 & 0.001 & 0.000\\
B & 1 & 0.004 & 0.004 & 0.002 & \textcolor{red}{-0.000} & 0.001 & 0.000 & \textcolor{red}{-0.000} & \textcolor{red}{-0.001} & \textcolor{red}{-0.001} & 0.001 & \textcolor{red}{-0.000} & \textcolor{red}{-0.001}\\
B & 2 & 0.001 & 0.001 & \textcolor{red}{-0.001} & \textcolor{red}{-0.000} & 0.001 & 0.001 & \textcolor{red}{-0.000} & \textcolor{red}{-0.001} & \textcolor{red}{-0.001} & 0.000 & \textcolor{red}{-0.000} & 0.000\\
AA & 1 & 0.005 & 0.002 & 0.000 & 0.001 & 0.002 & 0.001 & \textcolor{red}{-0.001} & 0.000 & 0.001 & 0.000 & \textcolor{red}{-0.000} & 0.000\\
AA & 2 & 0.001 & 0.000 & 0.002 & 0.000 & 0.001 & 0.000 & 0.000 & \textcolor{red}{-0.001} & \textcolor{red}{-0.000} & \textcolor{red}{-0.001} & \textcolor{red}{-0.000} & \textcolor{red}{-0.000}\\
AB & 1 & 0.001 & 0.000 & \textcolor{red}{-0.001} & 0.000 & \textcolor{red}{-0.001} & \textcolor{red}{-0.002} & \textcolor{red}{-0.002} & \textcolor{red}{-0.001} & \textcolor{red}{-0.001} & \textcolor{red}{-0.002} & \textcolor{red}{-0.002} & \textcolor{red}{-0.000}\\
AB & 2 & 0.002 & 0.000 & 0.001 & 0.000 & 0.001 & \textcolor{red}{-0.001} & \textcolor{red}{-0.001} & \textcolor{red}{-0.001} & 0.000 & \textcolor{red}{-0.001} & \textcolor{red}{-0.001} & \textcolor{red}{-0.000}\\
BA & 1 & 0.001 & \textcolor{red}{-0.000} & 0.000 & \textcolor{red}{-0.000} & \textcolor{red}{-0.000} & 0.000 & 0.000 & \textcolor{red}{-0.001} & \textcolor{red}{-0.001} & \textcolor{red}{-0.001} & 0.000 & \textcolor{red}{-0.000}\\
BA & 2 & \textcolor{red}{-0.001} & \textcolor{red}{-0.000} & 0.001 & 0.001 & \textcolor{red}{-0.000} & \textcolor{red}{-0.000} & 0.001 & \textcolor{red}{-0.000} & \textcolor{red}{-0.000} & \textcolor{red}{-0.001} & 0.001 & \textcolor{red}{-0.000}\\
BB & 1 & \textcolor{red}{-0.000} & 0.001 & 0.000 & \textcolor{red}{-0.001} & \textcolor{red}{-0.000} & \textcolor{red}{-0.000} & \textcolor{red}{-0.001} & \textcolor{red}{-0.001} & \textcolor{red}{-0.000} & \textcolor{red}{-0.001} & \textcolor{red}{-0.000} & \textcolor{red}{-0.000}\\
BB & 2 & \textcolor{red}{-0.001} & \textcolor{red}{-0.001} & 0.000 & \textcolor{red}{-0.000} & \textcolor{red}{-0.001} & \textcolor{red}{-0.002} & \textcolor{red}{-0.000} & \textcolor{red}{-0.000} & 0.000 & \textcolor{red}{-0.000} & \textcolor{red}{-0.000} & 0.000\\
BC & 1 & 0.001 & \textcolor{red}{-0.001} & 0.001 & \textcolor{red}{-0.000} & \textcolor{red}{-0.000} & 0.000 & 0.001 & \textcolor{red}{-0.000} & 0.000 & \textcolor{red}{-0.000} & 0.001 & 0.000\\
BC & 2 & \textcolor{red}{-0.001} & \textcolor{red}{-0.001} & 0.000 & 0.000 & \textcolor{red}{-0.001} & \textcolor{red}{-0.001} & 0.000 & \textcolor{red}{-0.000} & \textcolor{red}{-0.000} & \textcolor{red}{-0.001} & \textcolor{red}{-0.000} & \textcolor{red}{-0.001}\\
\hline
VAR &  &  &  &  &  &  &  &  &  &  &  &  & \\
\hline
Total & 1 & 0.036 & 0.060 & 0.028 & 0.027 & 0.031 & 0.048 & 0.020 & 0.031 & 0.012 & 0.038 & 0.012 & 0.025\\
Total & 2 & 0.041 & 0.065 & 0.030 & 0.028 & 0.031 & 0.049 & 0.022 & 0.028 & 0.014 & 0.035 & 0.012 & 0.025\\
A & 1 & 0.026 & 0.023 & 0.027 & 0.019 & 0.023 & 0.020 & 0.037 & 0.018 & 0.028 & 0.019 & 0.037 & 0.015\\
A & 2 & 0.026 & 0.025 & 0.024 & 0.017 & 0.022 & 0.024 & 0.034 & 0.017 & 0.025 & 0.022 & 0.033 & 0.010\\
B & 1 & 0.051 & 0.055 & 0.031 & 0.018 & 0.037 & 0.041 & 0.034 & 0.026 & 0.033 & 0.022 & 0.029 & 0.008\\
B & 2 & 0.054 & 0.050 & 0.025 & 0.016 & 0.042 & 0.043 & 0.034 & 0.021 & 0.032 & 0.020 & 0.027 & 0.006\\
AA & 1 & \textcolor{red}{-0.001} & \textcolor{red}{-0.009} & \textcolor{red}{-0.002} & \textcolor{red}{-0.005} & 0.000 & \textcolor{red}{-0.007} & 0.003 & \textcolor{red}{-0.003} & 0.009 & \textcolor{red}{-0.003} & 0.008 & \textcolor{red}{-0.005}\\
AA & 2 & 0.005 & \textcolor{red}{-0.009} & \textcolor{red}{-0.002} & \textcolor{red}{-0.003} & \textcolor{red}{-0.001} & \textcolor{red}{-0.011} & 0.003 & \textcolor{red}{-0.005} & 0.006 & \textcolor{red}{-0.006} & 0.004 & \textcolor{red}{-0.009}\\
AB & 1 & 0.008 & \textcolor{red}{-0.009} & \textcolor{red}{-0.006} & 0.001 & 0.008 & \textcolor{red}{-0.007} & 0.000 & \textcolor{red}{-0.004} & 0.015 & \textcolor{red}{-0.007} & 0.009 & \textcolor{red}{-0.004}\\
AB & 2 & 0.010 & \textcolor{red}{-0.007} & \textcolor{red}{-0.007} & \textcolor{red}{-0.002} & 0.009 & \textcolor{red}{-0.008} & \textcolor{red}{-0.000} & \textcolor{red}{-0.004} & 0.013 & \textcolor{red}{-0.008} & 0.006 & \textcolor{red}{-0.005}\\
BA & 1 & 0.012 & 0.003 & 0.006 & 0.001 & 0.002 & \textcolor{red}{-0.006} & 0.002 & \textcolor{red}{-0.004} & 0.004 & \textcolor{red}{-0.006} & 0.004 & \textcolor{red}{-0.004}\\
BA & 2 & 0.012 & \textcolor{red}{-0.001} & 0.002 & \textcolor{red}{-0.001} & 0.004 & \textcolor{red}{-0.003} & 0.001 & \textcolor{red}{-0.003} & 0.009 & \textcolor{red}{-0.005} & 0.004 & \textcolor{red}{-0.005}\\
BB & 1 & 0.013 & 0.003 & 0.003 & \textcolor{red}{-0.002} & 0.006 & \textcolor{red}{-0.004} & 0.005 & 0.001 & 0.008 & \textcolor{red}{-0.007} & 0.006 & \textcolor{red}{-0.004}\\
BB & 2 & 0.011 & \textcolor{red}{-0.003} & 0.000 & \textcolor{red}{-0.005} & 0.005 & \textcolor{red}{-0.002} & 0.004 & \textcolor{red}{-0.000} & 0.007 & \textcolor{red}{-0.006} & 0.007 & \textcolor{red}{-0.004}\\
BC & 1 & 0.012 & \textcolor{red}{-0.004} & 0.001 & \textcolor{red}{-0.004} & 0.004 & \textcolor{red}{-0.004} & 0.004 & \textcolor{red}{-0.001} & 0.008 & \textcolor{red}{-0.009} & 0.006 & \textcolor{red}{-0.004}\\
BC & 2 & 0.013 & \textcolor{red}{-0.006} & 0.001 & \textcolor{red}{-0.005} & 0.008 & \textcolor{red}{-0.002} & 0.000 & \textcolor{red}{-0.004} & 0.005 & \textcolor{red}{-0.008} & 0.007 & \textcolor{red}{-0.006}\\
\hline
    \end{tabular}%
    }
\end{table}

\begin{table}[!htb]
\caption{$\RelRMSE_{\Base}$ for all series in the hierarchy across different forecast horizons. ARIMA, ETS, and VAR models for base forecasts. Using the shrinkage approach to estimate $\boldsymbol{W}$ under scenario 6. Values in red indicate a $\RelRMSE_{\Base}$ less than 0.}
\centering\resizebox{1\textwidth}{!}{%
\begin{tabular}{llcccccccccccc}\hline
\multirow{2}{*}{Series} & \multirow{2}{*}{Variable}
& \multicolumn{12}{c}{Forecast horizons ($h$)} \\ \cline{3-14}
& & 1 & 2 & 3 & 4 & 5 & 6 & 7 & 8 & 9 & 10 & 11 & 12 \\ \hline
ARIMA &  &  &  &  &  &  &  &  &  &  &  &  & \\
\hline
Total & 1 & 0.052 & 0.041 & 0.013 & 0.025 & 0.039 & 0.036 & 0.024 & 0.040 & 0.038 & 0.031 & 0.031 & 0.032\\
Total & 2 & 0.016 & 0.043 & 0.035 & 0.025 & 0.032 & 0.045 & 0.032 & 0.047 & 0.035 & 0.050 & 0.045 & 0.032\\
A & 1 & 0.029 & 0.025 & 0.021 & 0.015 & 0.034 & 0.023 & 0.027 & 0.025 & 0.047 & 0.039 & 0.038 & 0.029\\
A & 2 & 0.029 & 0.034 & 0.015 & 0.013 & 0.025 & 0.010 & 0.019 & 0.015 & 0.019 & 0.027 & 0.029 & 0.022\\
B & 1 & 0.046 & 0.039 & 0.026 & 0.027 & 0.051 & 0.046 & 0.045 & 0.046 & 0.046 & 0.031 & 0.038 & 0.045\\
B & 2 & 0.028 & 0.038 & 0.036 & 0.024 & 0.033 & 0.036 & 0.036 & 0.027 & 0.026 & 0.038 & 0.038 & 0.028\\
AA & 1 & 0.024 & 0.010 & 0.013 & 0.014 & 0.014 & 0.010 & 0.011 & 0.003 & 0.004 & 0.007 & 0.011 & 0.016\\
AA & 2 & 0.023 & 0.016 & 0.004 & 0.005 & 0.014 & 0.011 & 0.008 & 0.011 & 0.017 & 0.019 & 0.022 & 0.018\\
AB & 1 & 0.016 & 0.011 & 0.007 & 0.011 & 0.018 & 0.012 & 0.012 & 0.007 & 0.015 & 0.019 & 0.020 & 0.017\\
AB & 2 & 0.016 & 0.018 & 0.005 & 0.010 & 0.020 & 0.015 & 0.015 & 0.013 & 0.028 & 0.016 & 0.023 & 0.012\\
BA & 1 & 0.017 & \textcolor{red}{-0.003} & 0.006 & 0.002 & 0.017 & 0.009 & 0.013 & 0.011 & 0.027 & 0.009 & 0.015 & 0.008\\
BA & 2 & 0.023 & 0.011 & 0.013 & 0.006 & 0.019 & 0.012 & 0.023 & 0.014 & 0.023 & 0.015 & 0.016 & 0.002\\
BB & 1 & 0.022 & 0.010 & 0.007 & 0.005 & 0.025 & 0.021 & 0.017 & 0.013 & 0.025 & 0.016 & 0.019 & 0.017\\
BB & 2 & 0.024 & 0.012 & 0.008 & 0.008 & 0.018 & 0.025 & 0.025 & 0.009 & 0.020 & 0.021 & 0.020 & 0.001\\
BC & 1 & 0.030 & 0.006 & 0.012 & 0.009 & 0.019 & 0.012 & 0.013 & 0.016 & 0.036 & 0.015 & 0.014 & 0.017\\
BC & 2 & 0.014 & 0.017 & \textcolor{red}{-0.003} & 0.003 & 0.020 & 0.018 & 0.020 & 0.015 & 0.020 & 0.015 & 0.019 & 0.006\\
\hline
ETS &  &  &  &  &  &  &  &  &  &  &  &  & \\
\hline
Total & 1 & 0.014 & 0.009 & 0.009 & 0.010 & 0.007 & 0.002 & 0.004 & 0.004 & 0.003 & 0.004 & 0.004 & 0.001\\
Total & 2 & 0.009 & 0.009 & 0.012 & 0.011 & 0.010 & 0.005 & 0.005 & 0.005 & 0.005 & 0.003 & 0.003 & 0.003\\
A & 1 & 0.005 & 0.008 & 0.009 & 0.009 & 0.010 & 0.010 & 0.008 & 0.008 & 0.011 & 0.010 & 0.008 & 0.009\\
A & 2 & 0.011 & 0.010 & 0.011 & 0.010 & 0.007 & 0.007 & 0.007 & 0.006 & 0.007 & 0.009 & 0.006 & 0.007\\
B & 1 & 0.002 & 0.003 & 0.006 & 0.008 & 0.010 & 0.008 & 0.006 & 0.007 & 0.007 & 0.007 & 0.009 & 0.005\\
B & 2 & 0.001 & 0.004 & 0.006 & 0.008 & 0.008 & 0.004 & 0.007 & 0.009 & 0.007 & 0.006 & 0.005 & 0.006\\
AA & 1 & 0.006 & 0.009 & 0.008 & 0.007 & 0.007 & 0.008 & 0.008 & 0.008 & 0.007 & 0.007 & 0.006 & 0.005\\
AA & 2 & 0.007 & 0.006 & 0.004 & 0.007 & 0.005 & 0.006 & 0.007 & 0.006 & 0.004 & 0.006 & 0.005 & 0.008\\
AB & 1 & 0.007 & 0.009 & 0.008 & 0.007 & 0.009 & 0.005 & 0.005 & 0.006 & 0.007 & 0.007 & 0.006 & 0.007\\
AB & 2 & 0.007 & 0.006 & 0.007 & 0.005 & 0.003 & 0.005 & 0.006 & 0.005 & 0.004 & 0.004 & 0.003 & 0.005\\
BA & 1 & 0.001 & 0.000 & 0.001 & 0.002 & 0.000 & \textcolor{red}{-0.001} & \textcolor{red}{-0.000} & 0.000 & 0.001 & 0.001 & \textcolor{red}{-0.001} & 0.001\\
BA & 2 & 0.001 & 0.002 & 0.001 & 0.001 & 0.001 & 0.000 & 0.002 & 0.003 & 0.003 & 0.003 & 0.002 & 0.002\\
BB & 1 & 0.000 & 0.001 & 0.001 & 0.004 & 0.001 & 0.001 & 0.002 & 0.003 & 0.002 & 0.002 & 0.001 & 0.000\\
BB & 2 & 0.002 & 0.001 & \textcolor{red}{-0.002} & 0.004 & 0.002 & 0.001 & 0.002 & 0.002 & 0.002 & 0.001 & 0.002 & 0.002\\
BC & 1 & \textcolor{red}{-0.001} & 0.002 & 0.002 & 0.000 & 0.000 & 0.002 & 0.001 & 0.001 & 0.003 & 0.001 & 0.001 & 0.003\\
BC & 2 & \textcolor{red}{-0.001} & 0.001 & 0.001 & 0.001 & 0.001 & \textcolor{red}{-0.000} & \textcolor{red}{-0.000} & 0.000 & 0.001 & 0.001 & 0.001 & 0.001\\
\hline
VAR &  &  &  &  &  &  &  &  &  &  &  &  & \\
\hline
Total & 1 & 0.016 & 0.019 & 0.019 & 0.020 & 0.014 & 0.011 & 0.015 & 0.012 & 0.016 & 0.011 & 0.009 & 0.017\\
Total & 2 & 0.001 & 0.022 & 0.008 & 0.006 & 0.018 & 0.026 & 0.019 & 0.022 & 0.017 & 0.028 & 0.008 & 0.011\\
A & 1 & 0.008 & 0.020 & 0.012 & 0.008 & 0.006 & 0.009 & 0.004 & 0.003 & 0.005 & 0.003 & \textcolor{red}{-0.000} & \textcolor{red}{-0.002}\\
A & 2 & 0.009 & 0.016 & 0.017 & 0.012 & 0.017 & 0.016 & 0.014 & 0.009 & 0.006 & 0.006 & 0.006 & 0.002\\
B & 1 & 0.013 & 0.014 & 0.012 & 0.005 & 0.010 & 0.013 & 0.005 & 0.008 & 0.012 & 0.015 & 0.003 & 0.004\\
B & 2 & 0.015 & 0.007 & 0.004 & 0.005 & 0.011 & 0.018 & 0.013 & 0.008 & 0.004 & 0.005 & 0.009 & 0.003\\
AA & 1 & \textcolor{red}{-0.000} & 0.005 & 0.000 & 0.000 & \textcolor{red}{-0.003} & \textcolor{red}{-0.007} & \textcolor{red}{-0.004} & \textcolor{red}{-0.003} & \textcolor{red}{-0.004} & \textcolor{red}{-0.005} & \textcolor{red}{-0.003} & \textcolor{red}{-0.003}\\
AA & 2 & \textcolor{red}{-0.000} & 0.005 & 0.003 & 0.002 & \textcolor{red}{-0.002} & \textcolor{red}{-0.004} & \textcolor{red}{-0.002} & \textcolor{red}{-0.004} & \textcolor{red}{-0.000} & \textcolor{red}{-0.002} & 0.003 & 0.001\\
AB & 1 & \textcolor{red}{-0.004} & 0.005 & \textcolor{red}{-0.000} & \textcolor{red}{-0.005} & 0.000 & \textcolor{red}{-0.002} & \textcolor{red}{-0.002} & \textcolor{red}{-0.002} & \textcolor{red}{-0.001} & \textcolor{red}{-0.005} & \textcolor{red}{-0.005} & \textcolor{red}{-0.006}\\
AB & 2 & \textcolor{red}{-0.002} & \textcolor{red}{-0.001} & 0.002 & 0.000 & \textcolor{red}{-0.000} & \textcolor{red}{-0.003} & \textcolor{red}{-0.000} & \textcolor{red}{-0.005} & 0.002 & \textcolor{red}{-0.003} & 0.001 & 0.002\\
BA & 1 & \textcolor{red}{-0.001} & \textcolor{red}{-0.006} & \textcolor{red}{-0.002} & 0.000 & \textcolor{red}{-0.001} & 0.001 & 0.002 & 0.000 & 0.002 & \textcolor{red}{-0.001} & 0.002 & \textcolor{red}{-0.000}\\
BA & 2 & 0.001 & \textcolor{red}{-0.012} & \textcolor{red}{-0.008} & \textcolor{red}{-0.004} & \textcolor{red}{-0.003} & \textcolor{red}{-0.008} & \textcolor{red}{-0.003} & \textcolor{red}{-0.004} & \textcolor{red}{-0.001} & \textcolor{red}{-0.001} & \textcolor{red}{-0.001} & \textcolor{red}{-0.002}\\
BB & 1 & \textcolor{red}{-0.001} & \textcolor{red}{-0.006} & \textcolor{red}{-0.006} & \textcolor{red}{-0.002} & \textcolor{red}{-0.001} & 0.001 & \textcolor{red}{-0.000} & \textcolor{red}{-0.001} & \textcolor{red}{-0.000} & \textcolor{red}{-0.004} & 0.000 & 0.000\\
BB & 2 & 0.003 & \textcolor{red}{-0.005} & \textcolor{red}{-0.005} & 0.001 & \textcolor{red}{-0.001} & \textcolor{red}{-0.005} & \textcolor{red}{-0.003} & 0.001 & 0.000 & \textcolor{red}{-0.004} & \textcolor{red}{-0.001} & \textcolor{red}{-0.003}\\
BC & 1 & 0.001 & \textcolor{red}{-0.002} & \textcolor{red}{-0.003} & \textcolor{red}{-0.004} & 0.000 & 0.003 & \textcolor{red}{-0.003} & \textcolor{red}{-0.002} & 0.002 & \textcolor{red}{-0.004} & \textcolor{red}{-0.004} & \textcolor{red}{-0.008}\\
BC & 2 & 0.003 & \textcolor{red}{-0.004} & \textcolor{red}{-0.005} & \textcolor{red}{-0.001} & 0.001 & \textcolor{red}{-0.001} & \textcolor{red}{-0.000} & 0.001 & 0.001 & \textcolor{red}{-0.001} & \textcolor{red}{-0.001} & \textcolor{red}{-0.003}\\
\hline
    \end{tabular}%
    }
\end{table}

\begin{table}[!htb]
\caption{$\RelRMSE_{\Base}$ for all series in the hierarchy across different forecast horizons. ARIMA, ETS, and VAR models for base forecasts. Using the shrinkage approach to estimate $\boldsymbol{W}$ under scenario 7. Values in red indicate a $\RelRMSE_{\Base}$ less than 0.}
\centering\resizebox{1\textwidth}{!}{%
\begin{tabular}{llcccccccccccc}\hline
\multirow{2}{*}{Series} & \multirow{2}{*}{Variable}
& \multicolumn{12}{c}{Forecast horizons ($h$)} \\ \cline{3-14}
& & 1 & 2 & 3 & 4 & 5 & 6 & 7 & 8 & 9 & 10 & 11 & 12 \\ \hline
ARIMA &  &  &  &  &  &  &  &  &  &  &  &  & \\
\hline
Total & 1 & 0.034 & 0.039 & 0.038 & 0.032 & 0.041 & 0.039 & 0.041 & 0.038 & 0.041 & 0.046 & 0.043 & 0.043\\
Total & 2 & 0.034 & 0.023 & 0.016 & 0.014 & 0.020 & 0.012 & 0.017 & 0.017 & 0.021 & 0.017 & 0.013 & 0.026\\
A & 1 & 0.017 & 0.029 & 0.033 & 0.033 & 0.033 & 0.044 & 0.048 & 0.050 & 0.047 & 0.047 & 0.041 & 0.047\\
A & 2 & 0.022 & 0.031 & 0.018 & 0.013 & 0.024 & 0.027 & 0.021 & 0.031 & 0.030 & 0.035 & 0.031 & 0.022\\
B & 1 & 0.009 & 0.022 & 0.014 & 0.007 & 0.011 & 0.017 & 0.012 & 0.005 & 0.004 & 0.012 & 0.013 & 0.012\\
B & 2 & 0.015 & 0.015 & 0.006 & 0.012 & 0.010 & 0.024 & 0.019 & 0.015 & 0.016 & 0.019 & 0.019 & 0.019\\
AA & 1 & 0.032 & 0.017 & 0.007 & 0.007 & 0.020 & 0.020 & 0.019 & 0.017 & 0.024 & 0.024 & 0.022 & 0.019\\
AA & 2 & 0.017 & 0.020 & 0.020 & 0.015 & 0.025 & 0.025 & 0.029 & 0.017 & 0.022 & 0.022 & 0.019 & 0.024\\
AB & 1 & 0.018 & 0.017 & 0.022 & 0.010 & 0.019 & 0.014 & 0.023 & 0.012 & 0.015 & 0.018 & 0.028 & 0.020\\
AB & 2 & 0.019 & 0.011 & 0.018 & 0.013 & 0.013 & 0.013 & 0.020 & 0.018 & 0.026 & 0.016 & 0.026 & 0.028\\
BA & 1 & 0.028 & 0.007 & 0.009 & 0.011 & 0.021 & 0.007 & 0.012 & 0.014 & 0.033 & 0.024 & 0.026 & 0.019\\
BA & 2 & 0.019 & 0.013 & 0.021 & 0.017 & 0.018 & 0.015 & 0.021 & 0.015 & 0.014 & 0.018 & 0.023 & 0.016\\
BB & 1 & 0.017 & 0.004 & 0.005 & 0.005 & 0.018 & 0.015 & 0.022 & 0.022 & 0.022 & 0.012 & 0.012 & 0.010\\
BB & 2 & 0.022 & 0.017 & 0.023 & 0.017 & 0.026 & 0.014 & 0.017 & 0.022 & 0.027 & 0.019 & 0.019 & 0.019\\
BC & 1 & 0.026 & 0.030 & 0.025 & 0.020 & 0.025 & 0.026 & 0.023 & 0.021 & 0.022 & 0.021 & 0.023 & 0.019\\
BC & 2 & 0.021 & 0.015 & 0.012 & 0.011 & 0.025 & 0.017 & 0.018 & 0.013 & 0.024 & 0.020 & 0.024 & 0.019\\
\hline
ETS &  &  &  &  &  &  &  &  &  &  &  &  & \\
\hline
Total & 1 & 0.011 & 0.005 & 0.004 & 0.003 & 0.005 & 0.003 & 0.004 & 0.002 & 0.002 & 0.002 & 0.002 & \textcolor{red}{-0.000}\\
Total & 2 & 0.007 & 0.003 & 0.002 & 0.001 & 0.001 & \textcolor{red}{-0.001} & 0.001 & 0.001 & 0.002 & 0.001 & \textcolor{red}{-0.001} & \textcolor{red}{-0.001}\\
A & 1 & 0.003 & 0.003 & 0.003 & 0.001 & 0.000 & 0.001 & 0.000 & 0.000 & \textcolor{red}{-0.000} & \textcolor{red}{-0.000} & 0.001 & 0.000\\
A & 2 & 0.003 & 0.004 & 0.004 & 0.001 & 0.003 & 0.004 & 0.003 & 0.003 & 0.005 & 0.005 & 0.005 & 0.004\\
B & 1 & \textcolor{red}{-0.003} & \textcolor{red}{-0.000} & \textcolor{red}{-0.001} & \textcolor{red}{-0.001} & \textcolor{red}{-0.001} & \textcolor{red}{-0.001} & \textcolor{red}{-0.001} & \textcolor{red}{-0.000} & \textcolor{red}{-0.001} & \textcolor{red}{-0.002} & \textcolor{red}{-0.001} & \textcolor{red}{-0.000}\\
B & 2 & 0.001 & 0.001 & 0.001 & 0.002 & 0.002 & 0.004 & 0.002 & 0.002 & 0.001 & 0.003 & 0.002 & 0.003\\
AA & 1 & 0.003 & 0.003 & \textcolor{red}{-0.000} & 0.001 & 0.000 & 0.000 & 0.001 & 0.001 & 0.002 & 0.000 & 0.000 & 0.002\\
AA & 2 & 0.004 & 0.004 & 0.002 & 0.004 & 0.005 & 0.003 & 0.000 & 0.001 & 0.001 & \textcolor{red}{-0.000} & \textcolor{red}{-0.001} & 0.001\\
AB & 1 & \textcolor{red}{-0.002} & \textcolor{red}{-0.002} & \textcolor{red}{-0.001} & 0.000 & \textcolor{red}{-0.000} & \textcolor{red}{-0.000} & \textcolor{red}{-0.001} & 0.000 & 0.000 & 0.001 & \textcolor{red}{-0.000} & \textcolor{red}{-0.001}\\
AB & 2 & 0.000 & 0.000 & \textcolor{red}{-0.000} & 0.000 & \textcolor{red}{-0.002} & \textcolor{red}{-0.002} & \textcolor{red}{-0.001} & \textcolor{red}{-0.002} & \textcolor{red}{-0.002} & \textcolor{red}{-0.003} & \textcolor{red}{-0.001} & \textcolor{red}{-0.002}\\
BA & 1 & 0.004 & 0.001 & 0.002 & 0.002 & 0.001 & 0.001 & 0.001 & 0.000 & 0.001 & 0.001 & 0.002 & 0.001\\
BA & 2 & 0.000 & \textcolor{red}{-0.001} & \textcolor{red}{-0.005} & \textcolor{red}{-0.003} & \textcolor{red}{-0.004} & \textcolor{red}{-0.003} & \textcolor{red}{-0.002} & \textcolor{red}{-0.003} & \textcolor{red}{-0.003} & \textcolor{red}{-0.002} & \textcolor{red}{-0.001} & \textcolor{red}{-0.003}\\
BB & 1 & 0.003 & 0.004 & 0.001 & 0.001 & 0.002 & 0.001 & 0.002 & 0.001 & 0.001 & 0.000 & 0.000 & 0.001\\
BB & 2 & 0.008 & 0.006 & 0.007 & 0.006 & 0.007 & 0.005 & 0.006 & 0.004 & 0.005 & 0.003 & 0.003 & 0.004\\
BC & 1 & 0.002 & 0.001 & 0.001 & 0.001 & 0.001 & 0.001 & \textcolor{red}{-0.000} & 0.000 & 0.002 & 0.002 & 0.000 & 0.000\\
BC & 2 & 0.002 & \textcolor{red}{-0.000} & 0.002 & 0.001 & 0.002 & 0.001 & \textcolor{red}{-0.000} & 0.001 & 0.001 & 0.000 & 0.001 & 0.001\\
\hline
VAR &  &  &  &  &  &  &  &  &  &  &  &  & \\
\hline
Total & 1 & 0.011 & 0.025 & 0.015 & 0.012 & 0.009 & 0.026 & 0.018 & 0.022 & 0.002 & 0.029 & 0.002 & 0.020\\
Total & 2 & 0.024 & 0.021 & 0.023 & 0.025 & 0.017 & 0.029 & 0.016 & 0.017 & 0.003 & 0.018 & \textcolor{red}{-0.001} & 0.015\\
A & 1 & 0.020 & 0.017 & 0.014 & 0.013 & 0.019 & 0.022 & 0.012 & 0.016 & 0.028 & 0.024 & 0.017 & 0.019\\
A & 2 & 0.018 & 0.024 & 0.013 & 0.011 & 0.024 & 0.028 & 0.013 & 0.017 & 0.022 & 0.016 & 0.013 & 0.013\\
B & 1 & 0.007 & 0.009 & 0.004 & 0.005 & 0.005 & 0.016 & 0.011 & 0.015 & 0.013 & 0.022 & 0.012 & 0.014\\
B & 2 & 0.007 & 0.018 & 0.013 & 0.009 & 0.008 & 0.026 & 0.026 & 0.019 & 0.021 & 0.027 & 0.018 & 0.014\\
AA & 1 & 0.016 & 0.009 & 0.009 & 0.008 & 0.015 & 0.005 & 0.015 & 0.010 & 0.035 & 0.005 & 0.021 & \textcolor{red}{-0.000}\\
AA & 2 & 0.016 & 0.013 & 0.018 & 0.009 & 0.022 & 0.008 & 0.021 & 0.004 & 0.021 & \textcolor{red}{-0.006} & 0.017 & 0.003\\
AB & 1 & 0.005 & 0.007 & 0.029 & 0.007 & 0.019 & 0.008 & 0.035 & 0.008 & 0.012 & 0.010 & 0.033 & 0.007\\
AB & 2 & 0.004 & 0.009 & 0.009 & 0.001 & 0.007 & 0.001 & 0.023 & 0.012 & 0.024 & 0.017 & 0.034 & 0.014\\
BA & 1 & 0.025 & 0.003 & 0.012 & 0.012 & 0.032 & 0.016 & 0.019 & 0.017 & 0.062 & 0.021 & 0.030 & 0.014\\
BA & 2 & 0.026 & 0.012 & 0.013 & 0.013 & 0.029 & 0.016 & 0.022 & 0.015 & 0.032 & 0.006 & 0.025 & 0.007\\
BB & 1 & 0.021 & 0.014 & 0.014 & 0.011 & 0.027 & 0.016 & 0.034 & 0.005 & 0.022 & 0.012 & 0.025 & 0.002\\
BB & 2 & 0.025 & 0.016 & 0.030 & 0.015 & 0.037 & 0.017 & 0.015 & 0.008 & 0.039 & 0.009 & 0.029 & 0.009\\
BC & 1 & 0.022 & 0.019 & 0.026 & 0.007 & 0.022 & 0.005 & 0.019 & 0.005 & 0.017 & \textcolor{red}{-0.003} & 0.029 & 0.008\\
BC & 2 & 0.013 & 0.010 & 0.014 & 0.015 & 0.030 & 0.008 & 0.025 & 0.003 & 0.021 & 0.000 & 0.024 & 0.006\\
\hline
    \end{tabular}%
    }
\end{table}

\begin{table}[!htb]
\caption{$\RelRMSE_{\Base}$ for all series in the hierarchy across different forecast horizons. ARIMA, ETS, and VAR models for base forecasts. Using the shrinkage approach to estimate $\boldsymbol{W}$ under scenario 8. Values in red indicate a $\RelRMSE_{\Base}$ less than 0.}
\centering\resizebox{1\textwidth}{!}{%
\begin{tabular}{llcccccccccccc}\hline
\multirow{2}{*}{Series} & \multirow{2}{*}{Variable}
& \multicolumn{12}{c}{Forecast horizons ($h$)} \\ \cline{3-14}
& & 1 & 2 & 3 & 4 & 5 & 6 & 7 & 8 & 9 & 10 & 11 & 12 \\ \hline
ARIMA &  &  &  &  &  &  &  &  &  &  &  &  & \\
\hline
Total & 1 & 0.034 & 0.038 & 0.029 & 0.021 & 0.034 & 0.037 & 0.030 & 0.028 & 0.031 & 0.043 & 0.039 & 0.048\\
Total & 2 & 0.022 & 0.034 & 0.034 & 0.022 & 0.024 & 0.038 & 0.038 & 0.032 & 0.028 & 0.038 & 0.041 & 0.045\\
A & 1 & 0.021 & 0.024 & 0.023 & 0.012 & 0.031 & 0.030 & 0.026 & 0.033 & 0.040 & 0.041 & 0.044 & 0.044\\
A & 2 & 0.020 & 0.016 & 0.013 & 0.014 & 0.027 & 0.034 & 0.034 & 0.039 & 0.048 & 0.045 & 0.041 & 0.036\\
B & 1 & 0.017 & 0.007 & 0.012 & 0.012 & 0.011 & 0.016 & 0.015 & 0.012 & 0.015 & 0.017 & 0.015 & 0.017\\
B & 2 & 0.024 & 0.012 & 0.014 & 0.007 & 0.007 & 0.013 & 0.010 & 0.009 & 0.014 & 0.013 & 0.014 & 0.008\\
AA & 1 & 0.015 & 0.010 & 0.010 & 0.011 & 0.021 & 0.022 & 0.023 & 0.017 & 0.015 & 0.021 & 0.033 & 0.026\\
AA & 2 & 0.016 & 0.016 & 0.008 & 0.010 & 0.024 & 0.021 & 0.020 & 0.019 & 0.019 & 0.024 & 0.031 & 0.028\\
AB & 1 & 0.017 & 0.011 & 0.019 & 0.019 & 0.018 & 0.013 & 0.017 & 0.017 & 0.024 & 0.024 & 0.020 & 0.017\\
AB & 2 & 0.030 & 0.012 & 0.024 & 0.021 & 0.023 & 0.021 & 0.023 & 0.016 & 0.030 & 0.030 & 0.028 & 0.022\\
BA & 1 & 0.022 & 0.016 & 0.010 & 0.015 & 0.016 & 0.020 & 0.028 & 0.027 & 0.035 & 0.034 & 0.041 & 0.032\\
BA & 2 & 0.020 & 0.021 & 0.014 & 0.014 & 0.018 & 0.019 & 0.019 & 0.017 & 0.024 & 0.018 & 0.022 & 0.023\\
BB & 1 & 0.017 & 0.015 & 0.013 & 0.015 & 0.024 & 0.020 & 0.015 & 0.019 & 0.017 & 0.012 & 0.009 & 0.017\\
BB & 2 & 0.022 & 0.013 & 0.014 & 0.020 & 0.026 & 0.020 & 0.022 & 0.022 & 0.024 & 0.022 & 0.025 & 0.020\\
BC & 1 & 0.030 & 0.015 & 0.021 & 0.007 & 0.023 & 0.023 & 0.027 & 0.012 & 0.018 & 0.015 & 0.019 & 0.009\\
BC & 2 & 0.017 & 0.020 & 0.013 & 0.003 & 0.018 & 0.019 & 0.017 & 0.012 & 0.017 & 0.014 & 0.013 & 0.012\\
\hline
ETS &  &  &  &  &  &  &  &  &  &  &  &  & \\
\hline
Total & 1 & 0.005 & 0.002 & 0.002 & 0.001 & 0.002 & 0.002 & 0.002 & 0.002 & 0.001 & 0.001 & \textcolor{red}{-0.001} & \textcolor{red}{-0.001}\\
Total & 2 & 0.003 & 0.005 & 0.000 & 0.001 & 0.001 & 0.002 & 0.001 & 0.002 & 0.002 & 0.003 & 0.001 & 0.001\\
A & 1 & 0.003 & 0.002 & \textcolor{red}{-0.000} & 0.000 & 0.001 & 0.001 & \textcolor{red}{-0.000} & \textcolor{red}{-0.000} & 0.000 & 0.001 & 0.000 & \textcolor{red}{-0.000}\\
A & 2 & \textcolor{red}{-0.001} & \textcolor{red}{-0.003} & \textcolor{red}{-0.005} & \textcolor{red}{-0.002} & \textcolor{red}{-0.002} & \textcolor{red}{-0.001} & \textcolor{red}{-0.002} & \textcolor{red}{-0.002} & \textcolor{red}{-0.000} & 0.000 & \textcolor{red}{-0.001} & \textcolor{red}{-0.001}\\
B & 1 & 0.002 & 0.002 & 0.002 & 0.001 & 0.001 & 0.001 & 0.002 & 0.002 & 0.001 & 0.001 & 0.002 & 0.001\\
B & 2 & 0.005 & 0.002 & 0.002 & 0.001 & 0.002 & 0.001 & 0.001 & 0.001 & 0.000 & \textcolor{red}{-0.000} & 0.000 & 0.000\\
AA & 1 & 0.004 & 0.002 & 0.002 & 0.003 & 0.001 & 0.001 & 0.001 & 0.001 & 0.001 & 0.000 & 0.001 & 0.001\\
AA & 2 & 0.003 & 0.002 & 0.002 & 0.001 & 0.001 & 0.001 & 0.001 & 0.001 & 0.002 & 0.002 & 0.002 & 0.002\\
AB & 1 & \textcolor{red}{-0.000} & \textcolor{red}{-0.001} & \textcolor{red}{-0.001} & \textcolor{red}{-0.002} & \textcolor{red}{-0.002} & \textcolor{red}{-0.003} & \textcolor{red}{-0.002} & \textcolor{red}{-0.002} & \textcolor{red}{-0.001} & \textcolor{red}{-0.001} & \textcolor{red}{-0.000} & \textcolor{red}{-0.001}\\
AB & 2 & 0.005 & 0.002 & 0.004 & 0.002 & 0.002 & 0.000 & 0.001 & 0.000 & \textcolor{red}{-0.000} & \textcolor{red}{-0.002} & \textcolor{red}{-0.001} & \textcolor{red}{-0.001}\\
BA & 1 & 0.000 & 0.003 & 0.001 & \textcolor{red}{-0.001} & \textcolor{red}{-0.002} & 0.001 & 0.000 & \textcolor{red}{-0.002} & \textcolor{red}{-0.001} & \textcolor{red}{-0.000} & \textcolor{red}{-0.000} & \textcolor{red}{-0.001}\\
BA & 2 & 0.001 & 0.002 & 0.001 & 0.001 & 0.000 & 0.001 & 0.001 & 0.000 & 0.001 & 0.001 & 0.002 & 0.001\\
BB & 1 & 0.003 & 0.002 & 0.002 & 0.002 & 0.004 & 0.001 & 0.002 & 0.002 & 0.003 & 0.001 & 0.001 & 0.000\\
BB & 2 & 0.003 & 0.004 & 0.002 & 0.001 & 0.000 & 0.000 & \textcolor{red}{-0.000} & \textcolor{red}{-0.000} & \textcolor{red}{-0.001} & \textcolor{red}{-0.000} & \textcolor{red}{-0.002} & \textcolor{red}{-0.001}\\
BC & 1 & 0.006 & \textcolor{red}{-0.001} & \textcolor{red}{-0.000} & 0.000 & 0.000 & \textcolor{red}{-0.000} & \textcolor{red}{-0.001} & 0.001 & 0.001 & \textcolor{red}{-0.000} & \textcolor{red}{-0.000} & 0.001\\
BC & 2 & 0.003 & 0.003 & 0.003 & 0.001 & 0.002 & 0.002 & 0.001 & 0.001 & 0.002 & 0.001 & 0.001 & 0.001\\
\hline
VAR &  &  &  &  &  &  &  &  &  &  &  &  & \\
\hline
Total & 1 & 0.016 & 0.026 & 0.016 & 0.021 & 0.015 & 0.020 & 0.005 & 0.012 & \textcolor{red}{-0.009} & 0.020 & \textcolor{red}{-0.010} & 0.017\\
Total & 2 & 0.007 & 0.026 & 0.013 & 0.010 & 0.008 & 0.024 & 0.010 & 0.010 & \textcolor{red}{-0.011} & 0.022 & \textcolor{red}{-0.007} & 0.014\\
A & 1 & 0.022 & 0.026 & 0.020 & 0.013 & 0.024 & 0.021 & 0.016 & 0.014 & 0.018 & 0.008 & 0.004 & 0.004\\
A & 2 & 0.017 & 0.020 & 0.028 & 0.015 & 0.014 & 0.015 & 0.020 & 0.015 & 0.020 & 0.012 & 0.007 & \textcolor{red}{-0.000}\\
B & 1 & 0.005 & 0.017 & 0.020 & 0.012 & 0.009 & 0.025 & 0.022 & 0.019 & 0.020 & 0.019 & 0.009 & 0.010\\
B & 2 & 0.010 & 0.018 & 0.015 & 0.011 & 0.007 & 0.020 & 0.018 & 0.020 & 0.018 & 0.022 & 0.009 & 0.009\\
AA & 1 & 0.026 & 0.012 & 0.024 & 0.006 & 0.026 & 0.008 & 0.028 & \textcolor{red}{-0.002} & 0.017 & 0.004 & 0.037 & \textcolor{red}{-0.001}\\
AA & 2 & 0.021 & 0.008 & 0.015 & 0.001 & 0.025 & 0.003 & 0.023 & \textcolor{red}{-0.007} & 0.015 & 0.004 & 0.034 & \textcolor{red}{-0.001}\\
AB & 1 & 0.025 & 0.004 & 0.008 & 0.010 & 0.025 & 0.003 & 0.023 & 0.009 & 0.042 & 0.013 & 0.038 & 0.010\\
AB & 2 & 0.031 & \textcolor{red}{-0.003} & 0.006 & 0.006 & 0.029 & 0.003 & 0.020 & 0.010 & 0.044 & 0.007 & 0.034 & 0.008\\
BA & 1 & 0.039 & 0.014 & 0.010 & 0.001 & 0.033 & 0.009 & 0.015 & \textcolor{red}{-0.007} & 0.032 & 0.002 & 0.025 & \textcolor{red}{-0.005}\\
BA & 2 & 0.035 & 0.020 & 0.022 & 0.009 & 0.038 & 0.010 & 0.022 & \textcolor{red}{-0.004} & 0.035 & 0.003 & 0.027 & \textcolor{red}{-0.002}\\
BB & 1 & 0.037 & 0.025 & 0.040 & 0.025 & 0.043 & 0.021 & 0.039 & 0.016 & 0.046 & 0.011 & 0.054 & 0.016\\
BB & 2 & 0.029 & 0.015 & 0.033 & 0.018 & 0.033 & 0.018 & 0.039 & 0.010 & 0.037 & 0.011 & 0.050 & 0.009\\
BC & 1 & 0.038 & 0.011 & 0.025 & \textcolor{red}{-0.000} & 0.034 & 0.012 & 0.029 & 0.007 & 0.048 & 0.012 & 0.035 & 0.011\\
BC & 2 & 0.039 & 0.019 & 0.023 & 0.004 & 0.039 & 0.019 & 0.022 & 0.004 & 0.048 & 0.012 & 0.039 & 0.012\\
\hline
    \end{tabular}%
    }
\end{table}

\begin{table}[!htb]
\caption{$\RelRMSE_{\Base}$ for all series in the hierarchy across different forecast horizons. ARIMA, ETS, and VAR models for base forecasts. Using the shrinkage approach to estimate $\boldsymbol{W}$ under scenario 9. Values in red indicate a $\RelRMSE_{\Base}$ less than 0.}
\centering\resizebox{1\textwidth}{!}{%
\begin{tabular}{llcccccccccccc}\hline
\multirow{2}{*}{Series} & \multirow{2}{*}{Variable}
& \multicolumn{12}{c}{Forecast horizons ($h$)} \\ \cline{3-14}
& & 1 & 2 & 3 & 4 & 5 & 6 & 7 & 8 & 9 & 10 & 11 & 12 \\ \hline
ARIMA &  &  &  &  &  &  &  &  &  &  &  &  & \\
\hline
Total & 1 & 0.018 & 0.042 & 0.028 & 0.021 & 0.020 & 0.022 & 0.022 & 0.031 & 0.033 & 0.038 & 0.035 & 0.026\\
Total & 2 & 0.034 & 0.039 & 0.026 & 0.022 & 0.029 & 0.024 & 0.029 & 0.026 & 0.029 & 0.025 & 0.037 & 0.029\\
A & 1 & 0.014 & 0.018 & 0.027 & 0.033 & 0.021 & 0.021 & 0.034 & 0.027 & 0.027 & 0.038 & 0.032 & 0.035\\
A & 2 & 0.032 & 0.033 & 0.026 & 0.024 & 0.024 & 0.029 & 0.029 & 0.041 & 0.040 & 0.040 & 0.035 & 0.030\\
B & 1 & 0.015 & 0.024 & 0.012 & 0.010 & 0.009 & 0.007 & 0.006 & 0.003 & 0.010 & 0.001 & 0.009 & 0.009\\
B & 2 & 0.023 & 0.025 & 0.025 & 0.016 & 0.010 & 0.016 & 0.012 & 0.019 & 0.024 & 0.014 & 0.004 & 0.014\\
AA & 1 & 0.016 & 0.022 & 0.013 & 0.009 & 0.031 & 0.019 & 0.021 & 0.018 & 0.021 & 0.015 & 0.026 & 0.022\\
AA & 2 & 0.006 & 0.004 & 0.011 & 0.009 & 0.021 & 0.010 & 0.028 & 0.015 & 0.022 & 0.019 & 0.024 & 0.015\\
AB & 1 & 0.025 & 0.009 & 0.015 & 0.004 & 0.030 & 0.015 & 0.030 & 0.023 & 0.042 & 0.016 & 0.023 & 0.018\\
AB & 2 & 0.016 & 0.009 & 0.008 & 0.009 & 0.022 & 0.008 & 0.002 & 0.006 & 0.019 & 0.013 & 0.015 & 0.013\\
BA & 1 & 0.025 & 0.018 & 0.008 & 0.004 & 0.022 & 0.011 & 0.016 & 0.005 & 0.008 & 0.008 & 0.018 & 0.006\\
BA & 2 & 0.019 & 0.028 & 0.021 & 0.016 & 0.022 & 0.017 & 0.013 & 0.014 & 0.029 & 0.020 & 0.022 & 0.008\\
BB & 1 & 0.019 & 0.008 & 0.015 & 0.005 & 0.019 & 0.017 & 0.012 & 0.014 & 0.029 & 0.027 & 0.020 & 0.010\\
BB & 2 & 0.020 & 0.021 & 0.028 & 0.015 & 0.032 & 0.012 & 0.022 & 0.017 & 0.022 & 0.019 & 0.031 & 0.021\\
BC & 1 & 0.027 & 0.018 & 0.022 & 0.014 & 0.029 & 0.014 & 0.022 & 0.010 & 0.025 & 0.013 & 0.027 & 0.032\\
BC & 2 & 0.023 & 0.005 & 0.006 & 0.005 & 0.017 & 0.008 & 0.015 & 0.001 & 0.016 & 0.002 & 0.012 & 0.013\\
\hline
ETS &  &  &  &  &  &  &  &  &  &  &  &  & \\
\hline
Total & 1 & 0.007 & 0.010 & 0.005 & 0.008 & 0.007 & 0.007 & 0.008 & 0.007 & 0.005 & 0.004 & 0.004 & 0.002\\
Total & 2 & 0.007 & 0.011 & 0.012 & 0.013 & 0.010 & 0.009 & 0.006 & 0.008 & 0.006 & 0.008 & 0.005 & 0.007\\
A & 1 & 0.007 & \textcolor{red}{-0.001} & \textcolor{red}{-0.000} & \textcolor{red}{-0.001} & 0.001 & 0.000 & 0.003 & 0.002 & 0.003 & 0.003 & 0.002 & 0.003\\
A & 2 & 0.014 & 0.008 & 0.000 & 0.001 & 0.000 & 0.000 & 0.001 & 0.003 & 0.003 & 0.003 & \textcolor{red}{-0.000} & 0.001\\
B & 1 & 0.008 & 0.011 & 0.012 & 0.008 & 0.006 & 0.005 & 0.002 & 0.004 & 0.003 & 0.004 & 0.003 & 0.003\\
B & 2 & 0.010 & 0.005 & 0.005 & 0.003 & 0.001 & 0.002 & 0.001 & 0.001 & 0.002 & 0.001 & \textcolor{red}{-0.000} & \textcolor{red}{-0.000}\\
AA & 1 & 0.007 & 0.005 & 0.006 & 0.004 & 0.005 & 0.003 & 0.000 & 0.002 & 0.003 & 0.003 & 0.002 & 0.002\\
AA & 2 & 0.005 & 0.004 & 0.001 & 0.002 & 0.003 & 0.006 & 0.008 & 0.007 & 0.006 & 0.004 & 0.004 & 0.003\\
AB & 1 & 0.008 & 0.005 & 0.006 & 0.004 & 0.007 & 0.001 & 0.001 & 0.005 & 0.006 & 0.002 & 0.002 & 0.006\\
AB & 2 & 0.002 & 0.004 & 0.007 & 0.003 & 0.005 & \textcolor{red}{-0.000} & 0.002 & 0.001 & 0.003 & 0.005 & 0.006 & 0.005\\
BA & 1 & \textcolor{red}{-0.001} & 0.005 & 0.004 & 0.000 & 0.002 & 0.003 & 0.004 & 0.003 & 0.001 & 0.001 & 0.003 & \textcolor{red}{-0.000}\\
BA & 2 & 0.004 & 0.006 & 0.006 & 0.005 & 0.003 & 0.005 & 0.002 & 0.005 & 0.006 & 0.004 & 0.004 & 0.005\\
BB & 1 & 0.006 & 0.001 & 0.004 & 0.001 & 0.002 & 0.000 & 0.001 & 0.001 & 0.003 & 0.004 & \textcolor{red}{-0.001} & \textcolor{red}{-0.000}\\
BB & 2 & 0.007 & 0.005 & 0.005 & 0.003 & 0.002 & 0.005 & 0.001 & \textcolor{red}{-0.002} & \textcolor{red}{-0.001} & 0.001 & 0.002 & 0.001\\
BC & 1 & 0.008 & 0.004 & 0.004 & 0.008 & 0.012 & 0.008 & 0.004 & 0.002 & 0.005 & 0.002 & 0.007 & 0.008\\
BC & 2 & 0.005 & 0.002 & 0.001 & 0.005 & 0.010 & 0.005 & 0.006 & 0.006 & 0.007 & 0.006 & 0.008 & 0.007\\
\hline
VAR &  &  &  &  &  &  &  &  &  &  &  &  & \\
\hline
Total & 1 & 0.009 & 0.019 & 0.006 & \textcolor{red}{-0.006} & 0.008 & 0.016 & 0.016 & 0.007 & 0.012 & 0.019 & 0.010 & 0.011\\
Total & 2 & 0.009 & 0.019 & 0.008 & 0.008 & 0.009 & 0.003 & 0.001 & 0.002 & 0.002 & 0.011 & 0.001 & 0.011\\
A & 1 & 0.012 & 0.016 & 0.019 & 0.012 & 0.008 & 0.008 & 0.008 & 0.008 & 0.010 & 0.017 & 0.008 & 0.007\\
A & 2 & 0.009 & 0.008 & 0.009 & 0.003 & 0.003 & 0.011 & 0.004 & 0.006 & 0.004 & \textcolor{red}{-0.001} & \textcolor{red}{-0.004} & 0.002\\
B & 1 & 0.002 & 0.006 & 0.005 & 0.008 & 0.007 & 0.017 & 0.012 & 0.011 & 0.003 & 0.006 & 0.008 & 0.005\\
B & 2 & 0.009 & 0.010 & 0.007 & 0.006 & 0.008 & 0.013 & 0.014 & 0.012 & 0.011 & 0.008 & \textcolor{red}{-0.002} & 0.000\\
AA & 1 & \textcolor{red}{-0.000} & 0.013 & 0.014 & 0.006 & 0.016 & 0.007 & 0.008 & 0.003 & 0.001 & 0.002 & 0.010 & 0.002\\
AA & 2 & \textcolor{red}{-0.003} & 0.012 & 0.012 & 0.008 & 0.009 & 0.005 & 0.011 & 0.005 & 0.006 & 0.003 & 0.016 & 0.004\\
AB & 1 & 0.019 & \textcolor{red}{-0.010} & 0.005 & 0.011 & 0.010 & 0.011 & 0.014 & 0.008 & 0.020 & 0.008 & 0.010 & 0.011\\
AB & 2 & 0.008 & 0.001 & 0.007 & 0.005 & 0.012 & 0.012 & 0.012 & 0.003 & 0.016 & 0.012 & 0.014 & 0.006\\
BA & 1 & 0.003 & 0.010 & 0.003 & 0.013 & 0.016 & 0.012 & 0.014 & 0.004 & 0.013 & 0.005 & 0.006 & 0.001\\
BA & 2 & 0.007 & 0.014 & 0.015 & 0.015 & 0.022 & 0.019 & 0.003 & 0.004 & 0.028 & 0.012 & 0.009 & 0.005\\
BB & 1 & 0.017 & 0.014 & 0.022 & 0.007 & 0.007 & \textcolor{red}{-0.002} & 0.005 & 0.001 & 0.008 & 0.007 & 0.016 & 0.006\\
BB & 2 & 0.008 & 0.015 & 0.013 & 0.015 & 0.016 & 0.010 & 0.018 & 0.007 & 0.010 & 0.005 & 0.011 & 0.008\\
BC & 1 & 0.010 & 0.009 & 0.008 & 0.003 & 0.005 & \textcolor{red}{-0.001} & 0.001 & 0.003 & 0.015 & 0.008 & 0.013 & 0.008\\
BC & 2 & 0.003 & 0.004 & 0.005 & 0.007 & 0.005 & \textcolor{red}{-0.003} & 0.009 & 0.005 & 0.008 & 0.001 & 0.025 & 0.010\\
\hline
    \end{tabular}%
    }
\end{table}

\begin{table}[!htb]
\caption{$\RelRMSE_{\Uni}$ for all series in the hierarchy across different forecast horizons. ARIMA, ETS, and VAR models for base forecasts. Using the shrinkage approach to estimate $\boldsymbol{W}$ under scenario 1. Values in red indicate a $\RelRMSE_{\Uni}$ less than 0.}
\centering\resizebox{1\textwidth}{!}{%
\begin{tabular}{llcccccccccccc}\hline
\multirow{2}{*}{Series} & \multirow{2}{*}{Variable}
& \multicolumn{12}{c}{Forecast horizons ($h$)} \\ \cline{3-14}
& & 1 & 2 & 3 & 4 & 5 & 6 & 7 & 8 & 9 & 10 & 11 & 12 \\ \hline
ARIMA &  &  &  &  &  &  &  &  &  &  &  &  & \\
\hline
Total & 1 & 0.003 & 0.005 & 0.007 & 0.005 & 0.004 & 0.005 & 0.002 & \textcolor{red}{-0.001} & 0.000 & \textcolor{red}{-0.001} & \textcolor{red}{-0.004} & \textcolor{red}{-0.008}\\
Total & 2 & 0.004 & 0.002 & 0.003 & 0.002 & 0.001 & 0.002 & \textcolor{red}{-0.001} & \textcolor{red}{-0.001} & \textcolor{red}{-0.001} & \textcolor{red}{-0.004} & \textcolor{red}{-0.005} & \textcolor{red}{-0.007}\\
A & 1 & 0.006 & 0.005 & 0.005 & 0.003 & 0.004 & 0.005 & 0.004 & 0.002 & 0.001 & 0.002 & 0.001 & \textcolor{red}{-0.002}\\
A & 2 & 0.005 & 0.002 & 0.003 & 0.002 & \textcolor{red}{-0.000} & 0.000 & \textcolor{red}{-0.002} & 0.000 & 0.002 & \textcolor{red}{-0.002} & \textcolor{red}{-0.003} & \textcolor{red}{-0.005}\\
B & 1 & \textcolor{red}{-0.001} & 0.002 & 0.006 & 0.006 & 0.006 & 0.003 & 0.002 & 0.002 & 0.003 & \textcolor{red}{-0.001} & \textcolor{red}{-0.002} & \textcolor{red}{-0.005}\\
B & 2 & \textcolor{red}{-0.000} & \textcolor{red}{-0.001} & 0.002 & 0.002 & 0.002 & \textcolor{red}{-0.000} & \textcolor{red}{-0.003} & \textcolor{red}{-0.003} & \textcolor{red}{-0.000} & \textcolor{red}{-0.004} & \textcolor{red}{-0.003} & \textcolor{red}{-0.005}\\
AA & 1 & 0.000 & 0.003 & 0.001 & \textcolor{red}{-0.001} & \textcolor{red}{-0.002} & \textcolor{red}{-0.000} & \textcolor{red}{-0.001} & \textcolor{red}{-0.002} & \textcolor{red}{-0.003} & 0.000 & \textcolor{red}{-0.000} & \textcolor{red}{-0.000}\\
AA & 2 & 0.005 & 0.003 & 0.004 & 0.004 & \textcolor{red}{-0.000} & 0.002 & 0.001 & 0.003 & 0.002 & 0.001 & \textcolor{red}{-0.001} & 0.001\\
AB & 1 & 0.007 & 0.006 & 0.006 & 0.004 & 0.007 & 0.004 & 0.003 & \textcolor{red}{-0.000} & 0.003 & 0.002 & 0.001 & \textcolor{red}{-0.002}\\
AB & 2 & \textcolor{red}{-0.002} & \textcolor{red}{-0.000} & 0.000 & \textcolor{red}{-0.000} & \textcolor{red}{-0.002} & \textcolor{red}{-0.003} & \textcolor{red}{-0.004} & \textcolor{red}{-0.006} & \textcolor{red}{-0.003} & \textcolor{red}{-0.006} & \textcolor{red}{-0.007} & \textcolor{red}{-0.010}\\
BA & 1 & 0.000 & 0.002 & 0.004 & 0.006 & 0.005 & 0.000 & 0.001 & 0.002 & 0.005 & 0.002 & 0.001 & 0.001\\
BA & 2 & 0.001 & 0.001 & 0.004 & 0.005 & 0.003 & 0.003 & 0.002 & \textcolor{red}{-0.001} & \textcolor{red}{-0.001} & \textcolor{red}{-0.001} & 0.001 & \textcolor{red}{-0.001}\\
BB & 1 & \textcolor{red}{-0.001} & \textcolor{red}{-0.001} & 0.003 & 0.001 & 0.001 & 0.001 & \textcolor{red}{-0.000} & 0.001 & 0.001 & \textcolor{red}{-0.001} & \textcolor{red}{-0.000} & \textcolor{red}{-0.001}\\
BB & 2 & \textcolor{red}{-0.002} & \textcolor{red}{-0.003} & \textcolor{red}{-0.002} & \textcolor{red}{-0.002} & \textcolor{red}{-0.001} & \textcolor{red}{-0.003} & \textcolor{red}{-0.007} & \textcolor{red}{-0.007} & \textcolor{red}{-0.004} & \textcolor{red}{-0.004} & \textcolor{red}{-0.006} & \textcolor{red}{-0.007}\\
BC & 1 & \textcolor{red}{-0.001} & \textcolor{red}{-0.001} & 0.001 & 0.001 & 0.003 & \textcolor{red}{-0.000} & 0.000 & \textcolor{red}{-0.001} & \textcolor{red}{-0.001} & \textcolor{red}{-0.003} & \textcolor{red}{-0.004} & \textcolor{red}{-0.005}\\
BC & 2 & 0.000 & \textcolor{red}{-0.001} & 0.001 & 0.001 & 0.001 & 0.000 & 0.001 & \textcolor{red}{-0.000} & 0.001 & 0.000 & 0.001 & \textcolor{red}{-0.000}\\
\hline
ETS &  &  &  &  &  &  &  &  &  &  &  &  & \\
\hline
Total & 1 & 0.003 & 0.001 & 0.001 & 0.001 & 0.001 & 0.001 & 0.001 & 0.001 & 0.001 & 0.001 & 0.001 & 0.001\\
Total & 2 & 0.003 & 0.002 & 0.001 & 0.001 & 0.001 & 0.000 & 0.000 & 0.000 & 0.001 & 0.001 & 0.001 & 0.001\\
A & 1 & 0.002 & 0.000 & 0.000 & 0.000 & \textcolor{red}{-0.000} & \textcolor{red}{-0.000} & 0.000 & 0.000 & \textcolor{red}{-0.000} & 0.000 & 0.001 & 0.000\\
A & 2 & 0.001 & 0.000 & 0.000 & 0.000 & 0.000 & \textcolor{red}{-0.001} & \textcolor{red}{-0.000} & \textcolor{red}{-0.000} & 0.000 & \textcolor{red}{-0.000} & \textcolor{red}{-0.000} & \textcolor{red}{-0.000}\\
B & 1 & 0.002 & 0.001 & 0.001 & 0.001 & 0.001 & 0.001 & 0.001 & 0.001 & 0.001 & 0.001 & 0.001 & 0.001\\
B & 2 & 0.002 & 0.002 & 0.001 & 0.001 & 0.000 & 0.001 & 0.001 & 0.000 & 0.000 & 0.001 & 0.001 & 0.001\\
AA & 1 & 0.001 & \textcolor{red}{-0.000} & \textcolor{red}{-0.000} & \textcolor{red}{-0.000} & \textcolor{red}{-0.000} & \textcolor{red}{-0.001} & \textcolor{red}{-0.001} & \textcolor{red}{-0.001} & \textcolor{red}{-0.000} & \textcolor{red}{-0.000} & \textcolor{red}{-0.000} & \textcolor{red}{-0.000}\\
AA & 2 & 0.001 & 0.001 & 0.000 & 0.000 & 0.000 & \textcolor{red}{-0.000} & \textcolor{red}{-0.001} & \textcolor{red}{-0.000} & \textcolor{red}{-0.000} & \textcolor{red}{-0.000} & \textcolor{red}{-0.000} & \textcolor{red}{-0.000}\\
AB & 1 & 0.001 & 0.000 & 0.000 & 0.000 & 0.000 & \textcolor{red}{-0.000} & \textcolor{red}{-0.000} & 0.000 & 0.000 & \textcolor{red}{-0.000} & 0.000 & 0.000\\
AB & 2 & 0.000 & \textcolor{red}{-0.000} & 0.000 & 0.000 & \textcolor{red}{-0.000} & \textcolor{red}{-0.001} & \textcolor{red}{-0.001} & \textcolor{red}{-0.001} & \textcolor{red}{-0.000} & \textcolor{red}{-0.000} & \textcolor{red}{-0.000} & \textcolor{red}{-0.000}\\
BA & 1 & 0.000 & 0.000 & 0.000 & 0.000 & 0.001 & \textcolor{red}{-0.000} & \textcolor{red}{-0.000} & 0.000 & \textcolor{red}{-0.000} & 0.000 & 0.000 & 0.000\\
BA & 2 & 0.001 & 0.001 & 0.001 & 0.000 & 0.001 & 0.001 & 0.000 & 0.000 & 0.000 & \textcolor{red}{-0.000} & \textcolor{red}{-0.000} & 0.000\\
BB & 1 & 0.001 & 0.001 & 0.001 & 0.001 & 0.001 & 0.001 & 0.001 & 0.001 & 0.001 & 0.001 & 0.001 & 0.001\\
BB & 2 & 0.001 & 0.001 & 0.000 & 0.000 & \textcolor{red}{-0.001} & \textcolor{red}{-0.001} & \textcolor{red}{-0.000} & \textcolor{red}{-0.000} & 0.000 & 0.000 & \textcolor{red}{-0.000} & 0.000\\
BC & 1 & 0.000 & 0.000 & 0.000 & 0.000 & \textcolor{red}{-0.000} & \textcolor{red}{-0.000} & \textcolor{red}{-0.000} & 0.000 & 0.000 & 0.000 & 0.000 & \textcolor{red}{-0.000}\\
BC & 2 & 0.001 & 0.000 & \textcolor{red}{-0.000} & \textcolor{red}{-0.000} & \textcolor{red}{-0.000} & \textcolor{red}{-0.000} & \textcolor{red}{-0.000} & \textcolor{red}{-0.000} & \textcolor{red}{-0.000} & 0.000 & 0.000 & 0.000\\
\hline
VAR &  &  &  &  &  &  &  &  &  &  &  &  & \\
\hline
Total & 1 & \textcolor{red}{-0.021} & \textcolor{red}{-0.014} & \textcolor{red}{-0.005} & \textcolor{red}{-0.013} & \textcolor{red}{-0.017} & \textcolor{red}{-0.020} & \textcolor{red}{-0.004} & \textcolor{red}{-0.011} & \textcolor{red}{-0.011} & \textcolor{red}{-0.020} & \textcolor{red}{-0.011} & \textcolor{red}{-0.023}\\
Total & 2 & \textcolor{red}{-0.001} & \textcolor{red}{-0.001} & 0.006 & \textcolor{red}{-0.004} & \textcolor{red}{-0.002} & \textcolor{red}{-0.005} & \textcolor{red}{-0.006} & \textcolor{red}{-0.009} & \textcolor{red}{-0.003} & \textcolor{red}{-0.014} & 0.000 & \textcolor{red}{-0.013}\\
A & 1 & \textcolor{red}{-0.015} & \textcolor{red}{-0.015} & 0.004 & \textcolor{red}{-0.009} & \textcolor{red}{-0.014} & \textcolor{red}{-0.014} & 0.005 & \textcolor{red}{-0.009} & \textcolor{red}{-0.007} & \textcolor{red}{-0.006} & 0.014 & \textcolor{red}{-0.009}\\
A & 2 & \textcolor{red}{-0.026} & \textcolor{red}{-0.023} & \textcolor{red}{-0.018} & \textcolor{red}{-0.016} & \textcolor{red}{-0.025} & \textcolor{red}{-0.019} & \textcolor{red}{-0.027} & \textcolor{red}{-0.016} & \textcolor{red}{-0.023} & \textcolor{red}{-0.026} & \textcolor{red}{-0.036} & \textcolor{red}{-0.017}\\
B & 1 & \textcolor{red}{-0.014} & \textcolor{red}{-0.009} & \textcolor{red}{-0.015} & \textcolor{red}{-0.010} & \textcolor{red}{-0.017} & \textcolor{red}{-0.018} & \textcolor{red}{-0.014} & \textcolor{red}{-0.004} & \textcolor{red}{-0.011} & \textcolor{red}{-0.011} & \textcolor{red}{-0.024} & \textcolor{red}{-0.016}\\
B & 2 & 0.003 & 0.002 & 0.005 & \textcolor{red}{-0.001} & 0.003 & \textcolor{red}{-0.006} & \textcolor{red}{-0.005} & \textcolor{red}{-0.007} & 0.009 & \textcolor{red}{-0.004} & 0.012 & \textcolor{red}{-0.001}\\
AA & 1 & \textcolor{red}{-0.007} & \textcolor{red}{-0.001} & \textcolor{red}{-0.014} & \textcolor{red}{-0.001} & \textcolor{red}{-0.001} & 0.001 & \textcolor{red}{-0.021} & \textcolor{red}{-0.004} & \textcolor{red}{-0.013} & \textcolor{red}{-0.002} & \textcolor{red}{-0.025} & \textcolor{red}{-0.000}\\
AA & 2 & \textcolor{red}{-0.019} & \textcolor{red}{-0.020} & \textcolor{red}{-0.015} & \textcolor{red}{-0.014} & \textcolor{red}{-0.031} & \textcolor{red}{-0.025} & \textcolor{red}{-0.026} & \textcolor{red}{-0.013} & \textcolor{red}{-0.020} & \textcolor{red}{-0.023} & \textcolor{red}{-0.035} & \textcolor{red}{-0.016}\\
AB & 1 & \textcolor{red}{-0.051} & \textcolor{red}{-0.053} & \textcolor{red}{-0.031} & \textcolor{red}{-0.027} & \textcolor{red}{-0.057} & \textcolor{red}{-0.058} & \textcolor{red}{-0.032} & \textcolor{red}{-0.044} & \textcolor{red}{-0.048} & \textcolor{red}{-0.045} & \textcolor{red}{-0.019} & \textcolor{red}{-0.037}\\
AB & 2 & \textcolor{red}{-0.038} & \textcolor{red}{-0.036} & \textcolor{red}{-0.027} & \textcolor{red}{-0.028} & \textcolor{red}{-0.042} & \textcolor{red}{-0.036} & \textcolor{red}{-0.044} & \textcolor{red}{-0.040} & \textcolor{red}{-0.036} & \textcolor{red}{-0.034} & \textcolor{red}{-0.043} & \textcolor{red}{-0.035}\\
BA & 1 & \textcolor{red}{-0.058} & \textcolor{red}{-0.055} & \textcolor{red}{-0.045} & \textcolor{red}{-0.037} & \textcolor{red}{-0.072} & \textcolor{red}{-0.092} & \textcolor{red}{-0.074} & \textcolor{red}{-0.051} & \textcolor{red}{-0.057} & \textcolor{red}{-0.069} & \textcolor{red}{-0.068} & \textcolor{red}{-0.067}\\
BA & 2 & \textcolor{red}{-0.005} & \textcolor{red}{-0.009} & \textcolor{red}{-0.000} & \textcolor{red}{-0.003} & \textcolor{red}{-0.002} & \textcolor{red}{-0.006} & \textcolor{red}{-0.010} & \textcolor{red}{-0.010} & 0.002 & 0.001 & \textcolor{red}{-0.001} & \textcolor{red}{-0.002}\\
BB & 1 & \textcolor{red}{-0.013} & \textcolor{red}{-0.014} & \textcolor{red}{-0.038} & \textcolor{red}{-0.029} & \textcolor{red}{-0.052} & \textcolor{red}{-0.035} & \textcolor{red}{-0.039} & \textcolor{red}{-0.026} & \textcolor{red}{-0.051} & \textcolor{red}{-0.030} & \textcolor{red}{-0.043} & \textcolor{red}{-0.023}\\
BB & 2 & \textcolor{red}{-0.005} & \textcolor{red}{-0.001} & 0.002 & \textcolor{red}{-0.011} & \textcolor{red}{-0.006} & \textcolor{red}{-0.005} & \textcolor{red}{-0.009} & \textcolor{red}{-0.006} & 0.002 & 0.001 & 0.003 & \textcolor{red}{-0.003}\\
BC & 1 & \textcolor{red}{-0.017} & \textcolor{red}{-0.006} & \textcolor{red}{-0.016} & \textcolor{red}{-0.008} & \textcolor{red}{-0.015} & \textcolor{red}{-0.010} & \textcolor{red}{-0.021} & \textcolor{red}{-0.014} & \textcolor{red}{-0.021} & \textcolor{red}{-0.006} & \textcolor{red}{-0.024} & \textcolor{red}{-0.009}\\
BC & 2 & 0.002 & \textcolor{red}{-0.009} & \textcolor{red}{-0.003} & \textcolor{red}{-0.004} & \textcolor{red}{-0.009} & \textcolor{red}{-0.019} & \textcolor{red}{-0.016} & \textcolor{red}{-0.015} & \textcolor{red}{-0.009} & \textcolor{red}{-0.012} & \textcolor{red}{-0.009} & \textcolor{red}{-0.008}\\
\hline
    \end{tabular}%
    }
\end{table}

\begin{table}[!htb]
\caption{$\RelRMSE_{\Uni}$ for all series in the hierarchy across different forecast horizons. ARIMA, ETS, and VAR models for base forecasts. Using the shrinkage approach to estimate $\boldsymbol{W}$ under scenario 2. Values in red indicate a $\RelRMSE_{\Uni}$ less than 0.}
\centering\resizebox{1\textwidth}{!}{%
\begin{tabular}{llcccccccccccc}\hline
\multirow{2}{*}{Series} & \multirow{2}{*}{Variable}
& \multicolumn{12}{c}{Forecast horizons ($h$)} \\ \cline{3-14}
& & 1 & 2 & 3 & 4 & 5 & 6 & 7 & 8 & 9 & 10 & 11 & 12 \\ \hline
ARIMA &  &  &  &  &  &  &  &  &  &  &  &  & \\
\hline
Total & 1 & 0.004 & 0.001 & 0.004 & 0.007 & 0.008 & 0.009 & 0.009 & 0.008 & 0.012 & 0.010 & 0.009 & 0.005\\
Total & 2 & \textcolor{red}{-0.004} & \textcolor{red}{-0.007} & \textcolor{red}{-0.001} & 0.001 & 0.002 & 0.006 & 0.003 & 0.004 & 0.006 & 0.006 & 0.004 & 0.001\\
A & 1 & 0.001 & 0.004 & 0.008 & 0.006 & 0.006 & 0.005 & 0.005 & 0.002 & 0.004 & 0.001 & 0.003 & 0.002\\
A & 2 & 0.004 & \textcolor{red}{-0.003} & 0.001 & 0.003 & 0.003 & 0.006 & 0.004 & 0.002 & 0.006 & 0.003 & 0.004 & 0.002\\
B & 1 & 0.006 & \textcolor{red}{-0.001} & 0.001 & 0.006 & 0.008 & 0.007 & 0.005 & 0.007 & 0.010 & 0.012 & 0.009 & 0.004\\
B & 2 & \textcolor{red}{-0.002} & \textcolor{red}{-0.001} & 0.002 & 0.002 & 0.005 & 0.005 & 0.001 & 0.004 & 0.003 & 0.002 & 0.002 & 0.001\\
AA & 1 & \textcolor{red}{-0.003} & 0.001 & 0.004 & 0.002 & 0.004 & 0.004 & 0.004 & 0.000 & 0.003 & 0.003 & 0.005 & 0.003\\
AA & 2 & 0.001 & \textcolor{red}{-0.000} & 0.001 & 0.002 & 0.004 & 0.005 & 0.005 & 0.003 & 0.008 & 0.001 & 0.002 & 0.003\\
AB & 1 & \textcolor{red}{-0.000} & \textcolor{red}{-0.000} & 0.006 & 0.004 & 0.003 & 0.000 & 0.003 & \textcolor{red}{-0.002} & \textcolor{red}{-0.000} & \textcolor{red}{-0.002} & 0.002 & 0.003\\
AB & 2 & 0.003 & \textcolor{red}{-0.000} & 0.001 & 0.001 & 0.002 & 0.002 & 0.001 & \textcolor{red}{-0.002} & 0.002 & 0.001 & 0.000 & \textcolor{red}{-0.001}\\
BA & 1 & \textcolor{red}{-0.000} & \textcolor{red}{-0.001} & \textcolor{red}{-0.001} & 0.001 & 0.003 & 0.002 & 0.004 & \textcolor{red}{-0.002} & 0.001 & 0.004 & 0.004 & 0.002\\
BA & 2 & 0.001 & 0.003 & 0.002 & 0.004 & 0.004 & 0.006 & 0.003 & 0.004 & 0.006 & 0.009 & 0.008 & 0.006\\
BB & 1 & \textcolor{red}{-0.005} & \textcolor{red}{-0.004} & \textcolor{red}{-0.002} & 0.000 & \textcolor{red}{-0.002} & \textcolor{red}{-0.000} & \textcolor{red}{-0.001} & 0.003 & 0.005 & 0.005 & 0.005 & 0.002\\
BB & 2 & 0.003 & 0.000 & 0.000 & \textcolor{red}{-0.001} & 0.000 & 0.002 & 0.002 & 0.006 & 0.004 & 0.002 & 0.003 & 0.003\\
BC & 1 & 0.002 & 0.003 & 0.003 & 0.004 & 0.003 & 0.004 & 0.003 & 0.005 & 0.005 & 0.007 & 0.003 & 0.003\\
BC & 2 & \textcolor{red}{-0.001} & \textcolor{red}{-0.003} & 0.001 & 0.003 & 0.002 & 0.000 & \textcolor{red}{-0.001} & 0.001 & \textcolor{red}{-0.001} & \textcolor{red}{-0.003} & \textcolor{red}{-0.004} & \textcolor{red}{-0.001}\\
\hline
ETS &  &  &  &  &  &  &  &  &  &  &  &  & \\
\hline
Total & 1 & 0.002 & 0.001 & \textcolor{red}{-0.000} & \textcolor{red}{-0.000} & 0.001 & 0.001 & \textcolor{red}{-0.000} & 0.000 & 0.001 & 0.002 & 0.002 & 0.001\\
Total & 2 & 0.003 & 0.002 & 0.001 & 0.001 & 0.001 & 0.000 & 0.000 & 0.001 & 0.002 & 0.001 & 0.002 & 0.003\\
A & 1 & 0.002 & 0.001 & \textcolor{red}{-0.000} & \textcolor{red}{-0.000} & 0.000 & 0.001 & 0.000 & \textcolor{red}{-0.000} & 0.000 & 0.002 & 0.002 & 0.001\\
A & 2 & 0.001 & 0.001 & \textcolor{red}{-0.000} & \textcolor{red}{-0.001} & \textcolor{red}{-0.001} & \textcolor{red}{-0.002} & \textcolor{red}{-0.002} & \textcolor{red}{-0.001} & 0.000 & 0.001 & 0.002 & 0.001\\
B & 1 & 0.002 & 0.000 & 0.000 & 0.001 & 0.001 & 0.000 & 0.000 & 0.000 & 0.001 & 0.001 & 0.001 & 0.001\\
B & 2 & 0.003 & 0.002 & 0.002 & 0.002 & 0.002 & 0.001 & 0.000 & 0.001 & 0.003 & 0.002 & 0.001 & 0.002\\
AA & 1 & 0.002 & 0.002 & 0.000 & 0.000 & 0.001 & 0.001 & 0.001 & 0.000 & 0.000 & 0.001 & 0.001 & 0.001\\
AA & 2 & 0.000 & 0.000 & \textcolor{red}{-0.001} & \textcolor{red}{-0.001} & \textcolor{red}{-0.001} & \textcolor{red}{-0.001} & \textcolor{red}{-0.001} & \textcolor{red}{-0.000} & 0.000 & 0.001 & 0.001 & 0.001\\
AB & 1 & 0.001 & 0.001 & 0.000 & \textcolor{red}{-0.001} & \textcolor{red}{-0.000} & 0.001 & 0.000 & 0.000 & 0.001 & 0.002 & 0.001 & 0.001\\
AB & 2 & 0.001 & 0.001 & \textcolor{red}{-0.000} & \textcolor{red}{-0.001} & \textcolor{red}{-0.000} & \textcolor{red}{-0.001} & \textcolor{red}{-0.002} & \textcolor{red}{-0.001} & \textcolor{red}{-0.000} & \textcolor{red}{-0.001} & 0.001 & 0.001\\
BA & 1 & 0.001 & \textcolor{red}{-0.000} & \textcolor{red}{-0.000} & \textcolor{red}{-0.001} & \textcolor{red}{-0.000} & \textcolor{red}{-0.001} & \textcolor{red}{-0.002} & \textcolor{red}{-0.001} & \textcolor{red}{-0.001} & \textcolor{red}{-0.000} & \textcolor{red}{-0.001} & \textcolor{red}{-0.001}\\
BA & 2 & 0.002 & 0.001 & 0.001 & 0.001 & 0.000 & \textcolor{red}{-0.000} & \textcolor{red}{-0.000} & 0.000 & 0.002 & 0.002 & 0.001 & 0.000\\
BB & 1 & 0.000 & \textcolor{red}{-0.000} & \textcolor{red}{-0.000} & \textcolor{red}{-0.001} & \textcolor{red}{-0.000} & \textcolor{red}{-0.000} & \textcolor{red}{-0.001} & \textcolor{red}{-0.001} & 0.001 & 0.001 & 0.001 & 0.000\\
BB & 2 & 0.003 & 0.002 & 0.001 & 0.000 & 0.000 & \textcolor{red}{-0.000} & 0.000 & 0.000 & 0.001 & 0.002 & 0.002 & 0.002\\
BC & 1 & 0.001 & 0.000 & 0.000 & 0.000 & 0.002 & 0.002 & 0.001 & 0.000 & 0.001 & 0.001 & 0.001 & 0.001\\
BC & 2 & 0.001 & 0.001 & \textcolor{red}{-0.001} & \textcolor{red}{-0.001} & 0.000 & \textcolor{red}{-0.000} & \textcolor{red}{-0.000} & 0.000 & 0.002 & 0.001 & 0.001 & 0.002\\
\hline
VAR &  &  &  &  &  &  &  &  &  &  &  &  & \\
\hline
Total & 1 & \textcolor{red}{-0.034} & \textcolor{red}{-0.030} & \textcolor{red}{-0.014} & \textcolor{red}{-0.026} & \textcolor{red}{-0.040} & \textcolor{red}{-0.042} & \textcolor{red}{-0.026} & \textcolor{red}{-0.040} & \textcolor{red}{-0.041} & \textcolor{red}{-0.045} & \textcolor{red}{-0.008} & \textcolor{red}{-0.033}\\
Total & 2 & \textcolor{red}{-0.010} & \textcolor{red}{-0.009} & \textcolor{red}{-0.003} & \textcolor{red}{-0.012} & \textcolor{red}{-0.015} & \textcolor{red}{-0.032} & \textcolor{red}{-0.010} & \textcolor{red}{-0.016} & \textcolor{red}{-0.017} & \textcolor{red}{-0.023} & 0.005 & \textcolor{red}{-0.017}\\
A & 1 & \textcolor{red}{-0.041} & \textcolor{red}{-0.040} & \textcolor{red}{-0.016} & \textcolor{red}{-0.012} & \textcolor{red}{-0.011} & \textcolor{red}{-0.027} & \textcolor{red}{-0.017} & \textcolor{red}{-0.021} & \textcolor{red}{-0.007} & \textcolor{red}{-0.027} & 0.000 & \textcolor{red}{-0.012}\\
A & 2 & \textcolor{red}{-0.016} & \textcolor{red}{-0.014} & \textcolor{red}{-0.018} & \textcolor{red}{-0.012} & \textcolor{red}{-0.020} & \textcolor{red}{-0.027} & \textcolor{red}{-0.023} & \textcolor{red}{-0.018} & \textcolor{red}{-0.042} & \textcolor{red}{-0.028} & \textcolor{red}{-0.023} & \textcolor{red}{-0.034}\\
B & 1 & \textcolor{red}{-0.023} & \textcolor{red}{-0.008} & \textcolor{red}{-0.013} & \textcolor{red}{-0.020} & \textcolor{red}{-0.029} & \textcolor{red}{-0.025} & \textcolor{red}{-0.023} & \textcolor{red}{-0.035} & \textcolor{red}{-0.042} & \textcolor{red}{-0.028} & \textcolor{red}{-0.026} & \textcolor{red}{-0.030}\\
B & 2 & \textcolor{red}{-0.004} & \textcolor{red}{-0.019} & \textcolor{red}{-0.006} & \textcolor{red}{-0.016} & \textcolor{red}{-0.016} & \textcolor{red}{-0.032} & \textcolor{red}{-0.010} & \textcolor{red}{-0.016} & \textcolor{red}{-0.004} & \textcolor{red}{-0.007} & 0.024 & 0.002\\
AA & 1 & \textcolor{red}{-0.007} & \textcolor{red}{-0.009} & \textcolor{red}{-0.013} & 0.002 & \textcolor{red}{-0.001} & 0.005 & \textcolor{red}{-0.013} & \textcolor{red}{-0.001} & \textcolor{red}{-0.013} & \textcolor{red}{-0.005} & \textcolor{red}{-0.030} & \textcolor{red}{-0.003}\\
AA & 2 & \textcolor{red}{-0.008} & \textcolor{red}{-0.017} & \textcolor{red}{-0.023} & \textcolor{red}{-0.019} & \textcolor{red}{-0.014} & \textcolor{red}{-0.017} & \textcolor{red}{-0.019} & \textcolor{red}{-0.020} & \textcolor{red}{-0.026} & \textcolor{red}{-0.025} & \textcolor{red}{-0.031} & \textcolor{red}{-0.031}\\
AB & 1 & \textcolor{red}{-0.080} & \textcolor{red}{-0.093} & \textcolor{red}{-0.064} & \textcolor{red}{-0.050} & \textcolor{red}{-0.076} & \textcolor{red}{-0.103} & \textcolor{red}{-0.095} & \textcolor{red}{-0.069} & \textcolor{red}{-0.054} & \textcolor{red}{-0.086} & \textcolor{red}{-0.065} & \textcolor{red}{-0.046}\\
AB & 2 & \textcolor{red}{-0.040} & \textcolor{red}{-0.043} & \textcolor{red}{-0.030} & \textcolor{red}{-0.026} & \textcolor{red}{-0.035} & \textcolor{red}{-0.049} & \textcolor{red}{-0.041} & \textcolor{red}{-0.032} & \textcolor{red}{-0.059} & \textcolor{red}{-0.039} & \textcolor{red}{-0.024} & \textcolor{red}{-0.039}\\
BA & 1 & \textcolor{red}{-0.086} & \textcolor{red}{-0.097} & \textcolor{red}{-0.064} & \textcolor{red}{-0.062} & \textcolor{red}{-0.110} & \textcolor{red}{-0.120} & \textcolor{red}{-0.113} & \textcolor{red}{-0.110} & \textcolor{red}{-0.110} & \textcolor{red}{-0.113} & \textcolor{red}{-0.068} & \textcolor{red}{-0.077}\\
BA & 2 & \textcolor{red}{-0.008} & \textcolor{red}{-0.009} & \textcolor{red}{-0.005} & \textcolor{red}{-0.010} & \textcolor{red}{-0.006} & \textcolor{red}{-0.005} & \textcolor{red}{-0.005} & \textcolor{red}{-0.005} & \textcolor{red}{-0.006} & \textcolor{red}{-0.005} & 0.004 & \textcolor{red}{-0.001}\\
BB & 1 & \textcolor{red}{-0.032} & \textcolor{red}{-0.020} & \textcolor{red}{-0.018} & \textcolor{red}{-0.015} & \textcolor{red}{-0.025} & \textcolor{red}{-0.018} & \textcolor{red}{-0.026} & \textcolor{red}{-0.009} & \textcolor{red}{-0.030} & \textcolor{red}{-0.031} & \textcolor{red}{-0.050} & \textcolor{red}{-0.030}\\
BB & 2 & \textcolor{red}{-0.006} & \textcolor{red}{-0.016} & \textcolor{red}{-0.009} & \textcolor{red}{-0.013} & \textcolor{red}{-0.015} & \textcolor{red}{-0.032} & \textcolor{red}{-0.024} & \textcolor{red}{-0.020} & \textcolor{red}{-0.023} & \textcolor{red}{-0.019} & \textcolor{red}{-0.001} & \textcolor{red}{-0.005}\\
BC & 1 & \textcolor{red}{-0.000} & 0.002 & \textcolor{red}{-0.009} & \textcolor{red}{-0.007} & \textcolor{red}{-0.013} & 0.002 & \textcolor{red}{-0.008} & \textcolor{red}{-0.004} & \textcolor{red}{-0.017} & \textcolor{red}{-0.002} & \textcolor{red}{-0.014} & \textcolor{red}{-0.005}\\
BC & 2 & \textcolor{red}{-0.010} & \textcolor{red}{-0.016} & \textcolor{red}{-0.009} & \textcolor{red}{-0.010} & \textcolor{red}{-0.016} & \textcolor{red}{-0.021} & \textcolor{red}{-0.017} & \textcolor{red}{-0.013} & \textcolor{red}{-0.005} & \textcolor{red}{-0.012} & \textcolor{red}{-0.002} & \textcolor{red}{-0.006}\\
\hline
    \end{tabular}%
    }
\end{table}

\begin{table}[!htb]
\caption{$\RelRMSE_{\Uni}$ for all series in the hierarchy across different forecast horizons. ARIMA, ETS, and VAR models for base forecasts. Using the shrinkage approach to estimate $\boldsymbol{W}$ under scenario 3. Values in red indicate a $\RelRMSE_{\Uni}$ less than 0.}
\centering\resizebox{1\textwidth}{!}{%
\begin{tabular}{llcccccccccccc}\hline
\multirow{2}{*}{Series} & \multirow{2}{*}{Variable}
& \multicolumn{12}{c}{Forecast horizons ($h$)} \\ \cline{3-14}
& & 1 & 2 & 3 & 4 & 5 & 6 & 7 & 8 & 9 & 10 & 11 & 12 \\ \hline
ARIMA &  &  &  &  &  &  &  &  &  &  &  &  & \\
\hline
Total & 1 & \textcolor{red}{-0.004} & \textcolor{red}{-0.002} & \textcolor{red}{-0.000} & \textcolor{red}{-0.002} & \textcolor{red}{-0.006} & \textcolor{red}{-0.008} & \textcolor{red}{-0.010} & \textcolor{red}{-0.009} & \textcolor{red}{-0.012} & \textcolor{red}{-0.014} & \textcolor{red}{-0.016} & \textcolor{red}{-0.016}\\
Total & 2 & 0.001 & 0.001 & \textcolor{red}{-0.001} & \textcolor{red}{-0.004} & \textcolor{red}{-0.004} & \textcolor{red}{-0.004} & \textcolor{red}{-0.006} & \textcolor{red}{-0.011} & \textcolor{red}{-0.011} & \textcolor{red}{-0.013} & \textcolor{red}{-0.012} & \textcolor{red}{-0.014}\\
A & 1 & \textcolor{red}{-0.001} & \textcolor{red}{-0.004} & 0.001 & \textcolor{red}{-0.002} & \textcolor{red}{-0.006} & \textcolor{red}{-0.007} & \textcolor{red}{-0.005} & \textcolor{red}{-0.006} & \textcolor{red}{-0.010} & \textcolor{red}{-0.010} & \textcolor{red}{-0.009} & \textcolor{red}{-0.010}\\
A & 2 & 0.001 & 0.002 & 0.001 & \textcolor{red}{-0.004} & \textcolor{red}{-0.004} & \textcolor{red}{-0.004} & \textcolor{red}{-0.004} & \textcolor{red}{-0.005} & \textcolor{red}{-0.007} & \textcolor{red}{-0.008} & \textcolor{red}{-0.008} & \textcolor{red}{-0.007}\\
B & 1 & \textcolor{red}{-0.002} & 0.002 & \textcolor{red}{-0.001} & \textcolor{red}{-0.005} & \textcolor{red}{-0.005} & \textcolor{red}{-0.009} & \textcolor{red}{-0.010} & \textcolor{red}{-0.009} & \textcolor{red}{-0.010} & \textcolor{red}{-0.013} & \textcolor{red}{-0.014} & \textcolor{red}{-0.013}\\
B & 2 & 0.003 & 0.004 & \textcolor{red}{-0.001} & \textcolor{red}{-0.002} & \textcolor{red}{-0.004} & \textcolor{red}{-0.003} & \textcolor{red}{-0.004} & \textcolor{red}{-0.008} & \textcolor{red}{-0.011} & \textcolor{red}{-0.010} & \textcolor{red}{-0.008} & \textcolor{red}{-0.010}\\
AA & 1 & \textcolor{red}{-0.002} & \textcolor{red}{-0.002} & 0.002 & \textcolor{red}{-0.002} & \textcolor{red}{-0.005} & \textcolor{red}{-0.004} & \textcolor{red}{-0.002} & \textcolor{red}{-0.002} & \textcolor{red}{-0.006} & \textcolor{red}{-0.004} & \textcolor{red}{-0.003} & \textcolor{red}{-0.004}\\
AA & 2 & 0.001 & 0.002 & 0.003 & \textcolor{red}{-0.000} & \textcolor{red}{-0.001} & \textcolor{red}{-0.002} & \textcolor{red}{-0.000} & \textcolor{red}{-0.001} & \textcolor{red}{-0.004} & \textcolor{red}{-0.004} & \textcolor{red}{-0.003} & \textcolor{red}{-0.005}\\
AB & 1 & 0.001 & \textcolor{red}{-0.001} & 0.002 & \textcolor{red}{-0.001} & \textcolor{red}{-0.002} & \textcolor{red}{-0.002} & \textcolor{red}{-0.003} & \textcolor{red}{-0.004} & \textcolor{red}{-0.005} & \textcolor{red}{-0.006} & \textcolor{red}{-0.006} & \textcolor{red}{-0.007}\\
AB & 2 & 0.003 & 0.003 & \textcolor{red}{-0.000} & \textcolor{red}{-0.006} & \textcolor{red}{-0.005} & \textcolor{red}{-0.004} & \textcolor{red}{-0.003} & \textcolor{red}{-0.004} & \textcolor{red}{-0.005} & \textcolor{red}{-0.007} & \textcolor{red}{-0.006} & \textcolor{red}{-0.005}\\
BA & 1 & 0.000 & 0.002 & \textcolor{red}{-0.001} & \textcolor{red}{-0.002} & \textcolor{red}{-0.006} & \textcolor{red}{-0.007} & \textcolor{red}{-0.004} & \textcolor{red}{-0.005} & \textcolor{red}{-0.007} & \textcolor{red}{-0.007} & \textcolor{red}{-0.005} & \textcolor{red}{-0.006}\\
BA & 2 & 0.003 & 0.005 & \textcolor{red}{-0.001} & \textcolor{red}{-0.007} & \textcolor{red}{-0.001} & 0.000 & \textcolor{red}{-0.001} & \textcolor{red}{-0.003} & \textcolor{red}{-0.002} & \textcolor{red}{-0.001} & \textcolor{red}{-0.001} & \textcolor{red}{-0.006}\\
BB & 1 & \textcolor{red}{-0.002} & \textcolor{red}{-0.002} & \textcolor{red}{-0.003} & \textcolor{red}{-0.003} & \textcolor{red}{-0.003} & \textcolor{red}{-0.007} & \textcolor{red}{-0.006} & \textcolor{red}{-0.004} & \textcolor{red}{-0.004} & \textcolor{red}{-0.008} & \textcolor{red}{-0.009} & \textcolor{red}{-0.006}\\
BB & 2 & 0.001 & 0.001 & 0.002 & 0.003 & 0.001 & 0.001 & 0.002 & \textcolor{red}{-0.001} & \textcolor{red}{-0.001} & \textcolor{red}{-0.001} & 0.001 & \textcolor{red}{-0.002}\\
BC & 1 & \textcolor{red}{-0.000} & 0.001 & \textcolor{red}{-0.002} & \textcolor{red}{-0.001} & \textcolor{red}{-0.002} & \textcolor{red}{-0.004} & \textcolor{red}{-0.008} & \textcolor{red}{-0.003} & \textcolor{red}{-0.002} & \textcolor{red}{-0.003} & \textcolor{red}{-0.006} & \textcolor{red}{-0.004}\\
BC & 2 & \textcolor{red}{-0.001} & \textcolor{red}{-0.001} & \textcolor{red}{-0.002} & \textcolor{red}{-0.002} & \textcolor{red}{-0.004} & \textcolor{red}{-0.003} & \textcolor{red}{-0.001} & \textcolor{red}{-0.001} & \textcolor{red}{-0.001} & \textcolor{red}{-0.004} & \textcolor{red}{-0.001} & \textcolor{red}{-0.002}\\
\hline
ETS &  &  &  &  &  &  &  &  &  &  &  &  & \\
\hline
Total & 1 & 0.001 & 0.001 & 0.000 & 0.001 & 0.001 & 0.000 & \textcolor{red}{-0.000} & \textcolor{red}{-0.000} & \textcolor{red}{-0.001} & \textcolor{red}{-0.001} & \textcolor{red}{-0.001} & \textcolor{red}{-0.001}\\
Total & 2 & 0.001 & 0.001 & 0.001 & 0.001 & 0.000 & 0.000 & 0.000 & \textcolor{red}{-0.000} & \textcolor{red}{-0.000} & \textcolor{red}{-0.000} & \textcolor{red}{-0.000} & \textcolor{red}{-0.001}\\
A & 1 & 0.002 & 0.001 & 0.001 & 0.001 & 0.001 & 0.000 & 0.000 & 0.000 & 0.000 & \textcolor{red}{-0.000} & \textcolor{red}{-0.000} & \textcolor{red}{-0.000}\\
A & 2 & 0.002 & 0.002 & 0.001 & 0.001 & 0.000 & 0.001 & 0.001 & 0.000 & 0.000 & 0.000 & 0.000 & 0.000\\
B & 1 & 0.002 & 0.002 & 0.001 & 0.001 & 0.001 & 0.001 & 0.000 & \textcolor{red}{-0.000} & \textcolor{red}{-0.000} & 0.000 & \textcolor{red}{-0.000} & \textcolor{red}{-0.000}\\
B & 2 & 0.001 & 0.001 & 0.001 & 0.001 & 0.001 & 0.001 & 0.000 & 0.000 & \textcolor{red}{-0.000} & 0.000 & \textcolor{red}{-0.000} & \textcolor{red}{-0.000}\\
AA & 1 & 0.000 & 0.000 & 0.000 & 0.001 & 0.000 & 0.000 & 0.000 & 0.000 & 0.000 & \textcolor{red}{-0.000} & \textcolor{red}{-0.000} & \textcolor{red}{-0.000}\\
AA & 2 & 0.001 & 0.001 & 0.000 & 0.000 & \textcolor{red}{-0.000} & \textcolor{red}{-0.000} & 0.000 & \textcolor{red}{-0.000} & \textcolor{red}{-0.000} & \textcolor{red}{-0.000} & \textcolor{red}{-0.000} & \textcolor{red}{-0.000}\\
AB & 1 & 0.001 & 0.001 & 0.001 & 0.001 & 0.001 & 0.001 & 0.001 & 0.000 & \textcolor{red}{-0.000} & \textcolor{red}{-0.000} & 0.000 & \textcolor{red}{-0.000}\\
AB & 2 & 0.001 & 0.001 & 0.000 & 0.001 & 0.000 & 0.001 & 0.000 & 0.000 & 0.000 & 0.000 & 0.000 & 0.000\\
BA & 1 & 0.001 & 0.001 & 0.001 & 0.001 & 0.000 & 0.001 & 0.001 & 0.000 & \textcolor{red}{-0.000} & \textcolor{red}{-0.000} & \textcolor{red}{-0.000} & \textcolor{red}{-0.000}\\
BA & 2 & 0.001 & 0.001 & 0.001 & 0.000 & 0.000 & 0.000 & 0.001 & 0.000 & 0.000 & 0.000 & 0.000 & 0.000\\
BB & 1 & 0.001 & 0.002 & 0.001 & 0.001 & 0.001 & 0.000 & 0.000 & 0.000 & \textcolor{red}{-0.000} & 0.000 & \textcolor{red}{-0.000} & \textcolor{red}{-0.000}\\
BB & 2 & 0.001 & 0.001 & \textcolor{red}{-0.000} & 0.001 & 0.000 & \textcolor{red}{-0.000} & \textcolor{red}{-0.000} & \textcolor{red}{-0.000} & \textcolor{red}{-0.000} & \textcolor{red}{-0.000} & \textcolor{red}{-0.000} & \textcolor{red}{-0.000}\\
BC & 1 & 0.001 & 0.001 & 0.001 & 0.000 & \textcolor{red}{-0.000} & \textcolor{red}{-0.000} & \textcolor{red}{-0.000} & \textcolor{red}{-0.000} & \textcolor{red}{-0.001} & \textcolor{red}{-0.000} & \textcolor{red}{-0.000} & \textcolor{red}{-0.000}\\
BC & 2 & 0.001 & 0.001 & 0.000 & 0.000 & \textcolor{red}{-0.000} & \textcolor{red}{-0.000} & \textcolor{red}{-0.000} & \textcolor{red}{-0.000} & \textcolor{red}{-0.000} & \textcolor{red}{-0.000} & \textcolor{red}{-0.000} & \textcolor{red}{-0.000}\\
\hline
VAR &  &  &  &  &  &  &  &  &  &  &  &  & \\
\hline
Total & 1 & \textcolor{red}{-0.024} & \textcolor{red}{-0.031} & \textcolor{red}{-0.017} & \textcolor{red}{-0.015} & \textcolor{red}{-0.010} & \textcolor{red}{-0.032} & \textcolor{red}{-0.015} & \textcolor{red}{-0.014} & 0.009 & \textcolor{red}{-0.009} & 0.001 & \textcolor{red}{-0.007}\\
Total & 2 & \textcolor{red}{-0.007} & \textcolor{red}{-0.008} & \textcolor{red}{-0.004} & \textcolor{red}{-0.007} & 0.001 & \textcolor{red}{-0.014} & \textcolor{red}{-0.000} & \textcolor{red}{-0.004} & 0.011 & \textcolor{red}{-0.005} & 0.009 & \textcolor{red}{-0.001}\\
A & 1 & \textcolor{red}{-0.022} & \textcolor{red}{-0.010} & \textcolor{red}{-0.005} & \textcolor{red}{-0.002} & \textcolor{red}{-0.006} & \textcolor{red}{-0.011} & 0.002 & \textcolor{red}{-0.006} & 0.004 & \textcolor{red}{-0.008} & 0.006 & \textcolor{red}{-0.004}\\
A & 2 & \textcolor{red}{-0.016} & \textcolor{red}{-0.013} & \textcolor{red}{-0.014} & \textcolor{red}{-0.021} & \textcolor{red}{-0.019} & \textcolor{red}{-0.014} & \textcolor{red}{-0.007} & \textcolor{red}{-0.001} & \textcolor{red}{-0.006} & \textcolor{red}{-0.011} & \textcolor{red}{-0.013} & 0.005\\
B & 1 & \textcolor{red}{-0.019} & \textcolor{red}{-0.026} & \textcolor{red}{-0.019} & \textcolor{red}{-0.013} & \textcolor{red}{-0.011} & \textcolor{red}{-0.023} & \textcolor{red}{-0.028} & \textcolor{red}{-0.018} & 0.003 & \textcolor{red}{-0.004} & \textcolor{red}{-0.015} & \textcolor{red}{-0.011}\\
B & 2 & \textcolor{red}{-0.006} & \textcolor{red}{-0.009} & 0.004 & 0.001 & 0.006 & \textcolor{red}{-0.009} & 0.002 & \textcolor{red}{-0.004} & 0.010 & \textcolor{red}{-0.012} & 0.005 & \textcolor{red}{-0.006}\\
AA & 1 & \textcolor{red}{-0.003} & 0.009 & 0.001 & \textcolor{red}{-0.006} & \textcolor{red}{-0.004} & 0.003 & \textcolor{red}{-0.020} & \textcolor{red}{-0.010} & \textcolor{red}{-0.021} & 0.002 & \textcolor{red}{-0.023} & \textcolor{red}{-0.004}\\
AA & 2 & \textcolor{red}{-0.015} & \textcolor{red}{-0.011} & \textcolor{red}{-0.016} & \textcolor{red}{-0.019} & \textcolor{red}{-0.017} & \textcolor{red}{-0.014} & \textcolor{red}{-0.008} & \textcolor{red}{-0.013} & \textcolor{red}{-0.012} & \textcolor{red}{-0.008} & \textcolor{red}{-0.013} & \textcolor{red}{-0.003}\\
AB & 1 & \textcolor{red}{-0.044} & \textcolor{red}{-0.056} & \textcolor{red}{-0.038} & \textcolor{red}{-0.029} & \textcolor{red}{-0.051} & \textcolor{red}{-0.060} & \textcolor{red}{-0.032} & \textcolor{red}{-0.032} & \textcolor{red}{-0.036} & \textcolor{red}{-0.042} & \textcolor{red}{-0.029} & \textcolor{red}{-0.027}\\
AB & 2 & \textcolor{red}{-0.027} & \textcolor{red}{-0.025} & \textcolor{red}{-0.015} & \textcolor{red}{-0.030} & \textcolor{red}{-0.034} & \textcolor{red}{-0.031} & \textcolor{red}{-0.022} & \textcolor{red}{-0.008} & \textcolor{red}{-0.019} & \textcolor{red}{-0.028} & \textcolor{red}{-0.021} & \textcolor{red}{-0.009}\\
BA & 1 & \textcolor{red}{-0.047} & \textcolor{red}{-0.085} & \textcolor{red}{-0.054} & \textcolor{red}{-0.046} & \textcolor{red}{-0.069} & \textcolor{red}{-0.088} & \textcolor{red}{-0.074} & \textcolor{red}{-0.059} & \textcolor{red}{-0.052} & \textcolor{red}{-0.061} & \textcolor{red}{-0.045} & \textcolor{red}{-0.046}\\
BA & 2 & \textcolor{red}{-0.013} & \textcolor{red}{-0.003} & 0.006 & \textcolor{red}{-0.004} & \textcolor{red}{-0.001} & \textcolor{red}{-0.007} & \textcolor{red}{-0.003} & \textcolor{red}{-0.001} & 0.002 & \textcolor{red}{-0.002} & 0.002 & \textcolor{red}{-0.006}\\
BB & 1 & \textcolor{red}{-0.015} & \textcolor{red}{-0.020} & \textcolor{red}{-0.032} & \textcolor{red}{-0.026} & \textcolor{red}{-0.030} & \textcolor{red}{-0.028} & \textcolor{red}{-0.038} & \textcolor{red}{-0.021} & \textcolor{red}{-0.028} & \textcolor{red}{-0.022} & \textcolor{red}{-0.034} & \textcolor{red}{-0.025}\\
BB & 2 & \textcolor{red}{-0.000} & \textcolor{red}{-0.001} & \textcolor{red}{-0.001} & \textcolor{red}{-0.002} & \textcolor{red}{-0.002} & \textcolor{red}{-0.014} & \textcolor{red}{-0.006} & \textcolor{red}{-0.008} & \textcolor{red}{-0.002} & \textcolor{red}{-0.010} & \textcolor{red}{-0.005} & \textcolor{red}{-0.004}\\
BC & 1 & \textcolor{red}{-0.007} & \textcolor{red}{-0.001} & \textcolor{red}{-0.009} & \textcolor{red}{-0.010} & \textcolor{red}{-0.010} & \textcolor{red}{-0.006} & \textcolor{red}{-0.025} & \textcolor{red}{-0.011} & \textcolor{red}{-0.005} & 0.005 & \textcolor{red}{-0.020} & \textcolor{red}{-0.006}\\
BC & 2 & \textcolor{red}{-0.004} & \textcolor{red}{-0.009} & \textcolor{red}{-0.000} & \textcolor{red}{-0.004} & \textcolor{red}{-0.001} & \textcolor{red}{-0.010} & \textcolor{red}{-0.011} & \textcolor{red}{-0.010} & 0.002 & \textcolor{red}{-0.009} & \textcolor{red}{-0.007} & \textcolor{red}{-0.010}\\
\hline
    \end{tabular}%
    }
\end{table}

\begin{table}[!htb]
\caption{$\RelRMSE_{\Uni}$ for all series in the hierarchy across different forecast horizons. ARIMA, ETS, and VAR models for base forecasts. Using the shrinkage approach to estimate $\boldsymbol{W}$ under scenario 4. Values in red indicate a $\RelRMSE_{\Uni}$ less than 0.}
\centering\resizebox{1\textwidth}{!}{%
\begin{tabular}{llcccccccccccc}\hline
\multirow{2}{*}{Series} & \multirow{2}{*}{Variable}
& \multicolumn{12}{c}{Forecast horizons ($h$)} \\ \cline{3-14}
& & 1 & 2 & 3 & 4 & 5 & 6 & 7 & 8 & 9 & 10 & 11 & 12 \\ \hline
ARIMA &  &  &  &  &  &  &  &  &  &  &  &  & \\
\hline
Total & 1 & 0.006 & 0.009 & 0.009 & 0.006 & \textcolor{red}{-0.001} & 0.003 & 0.003 & \textcolor{red}{-0.000} & 0.001 & 0.001 & 0.000 & \textcolor{red}{-0.003}\\
Total & 2 & 0.004 & 0.006 & 0.008 & 0.004 & \textcolor{red}{-0.002} & 0.001 & 0.000 & \textcolor{red}{-0.004} & \textcolor{red}{-0.004} & \textcolor{red}{-0.003} & \textcolor{red}{-0.003} & \textcolor{red}{-0.003}\\
A & 1 & 0.010 & 0.013 & 0.005 & \textcolor{red}{-0.001} & \textcolor{red}{-0.001} & 0.003 & 0.001 & \textcolor{red}{-0.001} & \textcolor{red}{-0.002} & \textcolor{red}{-0.001} & \textcolor{red}{-0.003} & \textcolor{red}{-0.005}\\
A & 2 & 0.004 & 0.010 & 0.007 & \textcolor{red}{-0.001} & \textcolor{red}{-0.003} & 0.001 & \textcolor{red}{-0.001} & \textcolor{red}{-0.004} & \textcolor{red}{-0.003} & \textcolor{red}{-0.004} & \textcolor{red}{-0.004} & \textcolor{red}{-0.004}\\
B & 1 & 0.003 & 0.007 & 0.007 & 0.009 & 0.004 & 0.006 & 0.002 & 0.000 & \textcolor{red}{-0.000} & 0.000 & \textcolor{red}{-0.001} & \textcolor{red}{-0.005}\\
B & 2 & 0.006 & 0.003 & 0.007 & 0.006 & 0.003 & \textcolor{red}{-0.000} & \textcolor{red}{-0.001} & \textcolor{red}{-0.003} & \textcolor{red}{-0.000} & \textcolor{red}{-0.002} & \textcolor{red}{-0.003} & \textcolor{red}{-0.006}\\
AA & 1 & 0.011 & 0.007 & 0.006 & 0.001 & 0.002 & 0.002 & 0.000 & 0.001 & 0.004 & 0.003 & 0.000 & \textcolor{red}{-0.002}\\
AA & 2 & 0.011 & 0.011 & 0.004 & 0.000 & 0.002 & 0.001 & \textcolor{red}{-0.001} & \textcolor{red}{-0.002} & \textcolor{red}{-0.003} & \textcolor{red}{-0.006} & \textcolor{red}{-0.004} & \textcolor{red}{-0.005}\\
AB & 1 & 0.004 & 0.009 & 0.003 & 0.002 & 0.003 & 0.007 & 0.004 & \textcolor{red}{-0.000} & 0.000 & \textcolor{red}{-0.002} & \textcolor{red}{-0.003} & \textcolor{red}{-0.006}\\
AB & 2 & 0.004 & 0.007 & 0.004 & 0.003 & 0.001 & \textcolor{red}{-0.000} & \textcolor{red}{-0.002} & \textcolor{red}{-0.004} & 0.001 & \textcolor{red}{-0.002} & \textcolor{red}{-0.007} & \textcolor{red}{-0.005}\\
BA & 1 & 0.003 & 0.006 & 0.004 & 0.003 & 0.005 & 0.000 & 0.001 & \textcolor{red}{-0.002} & \textcolor{red}{-0.003} & 0.001 & 0.000 & \textcolor{red}{-0.001}\\
BA & 2 & 0.003 & 0.006 & 0.006 & 0.004 & 0.003 & 0.000 & 0.001 & \textcolor{red}{-0.002} & \textcolor{red}{-0.001} & 0.001 & \textcolor{red}{-0.002} & 0.001\\
BB & 1 & 0.003 & 0.003 & 0.006 & 0.007 & 0.004 & 0.006 & 0.003 & 0.003 & 0.000 & 0.002 & \textcolor{red}{-0.000} & \textcolor{red}{-0.000}\\
BB & 2 & 0.004 & 0.004 & 0.009 & 0.007 & 0.007 & 0.004 & \textcolor{red}{-0.001} & 0.001 & \textcolor{red}{-0.000} & \textcolor{red}{-0.000} & \textcolor{red}{-0.004} & \textcolor{red}{-0.005}\\
BC & 1 & 0.005 & 0.005 & 0.008 & 0.010 & 0.008 & 0.004 & 0.000 & \textcolor{red}{-0.000} & 0.002 & \textcolor{red}{-0.002} & \textcolor{red}{-0.003} & \textcolor{red}{-0.001}\\
BC & 2 & 0.005 & 0.004 & 0.004 & 0.000 & 0.002 & \textcolor{red}{-0.000} & \textcolor{red}{-0.002} & \textcolor{red}{-0.003} & 0.001 & 0.002 & 0.001 & 0.001\\
\hline
ETS &  &  &  &  &  &  &  &  &  &  &  &  & \\
\hline
Total & 1 & 0.005 & 0.003 & 0.002 & 0.001 & 0.000 & 0.000 & 0.000 & 0.000 & \textcolor{red}{-0.000} & 0.000 & 0.001 & 0.000\\
Total & 2 & 0.004 & 0.002 & 0.001 & 0.001 & \textcolor{red}{-0.001} & \textcolor{red}{-0.001} & \textcolor{red}{-0.002} & \textcolor{red}{-0.001} & \textcolor{red}{-0.001} & \textcolor{red}{-0.001} & \textcolor{red}{-0.001} & \textcolor{red}{-0.000}\\
A & 1 & 0.003 & 0.002 & 0.000 & \textcolor{red}{-0.000} & 0.001 & 0.001 & 0.000 & 0.000 & 0.000 & 0.000 & 0.000 & \textcolor{red}{-0.000}\\
A & 2 & 0.002 & 0.001 & \textcolor{red}{-0.000} & \textcolor{red}{-0.001} & \textcolor{red}{-0.000} & \textcolor{red}{-0.000} & \textcolor{red}{-0.001} & \textcolor{red}{-0.001} & \textcolor{red}{-0.000} & \textcolor{red}{-0.000} & \textcolor{red}{-0.001} & \textcolor{red}{-0.000}\\
B & 1 & 0.004 & 0.003 & 0.002 & 0.001 & 0.000 & 0.001 & \textcolor{red}{-0.000} & \textcolor{red}{-0.000} & \textcolor{red}{-0.001} & \textcolor{red}{-0.000} & \textcolor{red}{-0.000} & \textcolor{red}{-0.000}\\
B & 2 & 0.004 & 0.002 & 0.002 & 0.002 & \textcolor{red}{-0.000} & \textcolor{red}{-0.000} & \textcolor{red}{-0.000} & 0.000 & \textcolor{red}{-0.000} & 0.000 & \textcolor{red}{-0.000} & \textcolor{red}{-0.000}\\
AA & 1 & 0.003 & 0.001 & 0.001 & 0.001 & 0.002 & 0.002 & 0.002 & 0.001 & 0.000 & 0.001 & 0.001 & 0.000\\
AA & 2 & 0.002 & 0.001 & 0.000 & \textcolor{red}{-0.001} & 0.000 & \textcolor{red}{-0.000} & \textcolor{red}{-0.001} & \textcolor{red}{-0.000} & 0.000 & 0.000 & \textcolor{red}{-0.000} & \textcolor{red}{-0.001}\\
AB & 1 & 0.002 & 0.003 & 0.002 & 0.001 & 0.001 & 0.001 & 0.001 & 0.001 & 0.002 & 0.002 & 0.001 & 0.001\\
AB & 2 & 0.001 & 0.000 & 0.001 & 0.001 & 0.001 & 0.001 & 0.001 & 0.001 & 0.001 & 0.001 & 0.000 & 0.000\\
BA & 1 & 0.002 & 0.002 & 0.001 & 0.001 & 0.001 & 0.000 & \textcolor{red}{-0.001} & 0.000 & \textcolor{red}{-0.000} & 0.000 & \textcolor{red}{-0.000} & 0.000\\
BA & 2 & 0.003 & 0.001 & 0.002 & 0.001 & 0.000 & \textcolor{red}{-0.000} & 0.000 & \textcolor{red}{-0.000} & \textcolor{red}{-0.000} & 0.000 & 0.000 & \textcolor{red}{-0.000}\\
BB & 1 & 0.003 & 0.003 & 0.002 & 0.001 & 0.001 & 0.001 & 0.001 & 0.000 & \textcolor{red}{-0.000} & \textcolor{red}{-0.000} & 0.000 & 0.000\\
BB & 2 & 0.001 & 0.001 & 0.001 & 0.001 & \textcolor{red}{-0.001} & \textcolor{red}{-0.001} & \textcolor{red}{-0.001} & \textcolor{red}{-0.001} & \textcolor{red}{-0.001} & 0.000 & 0.000 & 0.000\\
BC & 1 & 0.003 & 0.002 & 0.001 & 0.001 & 0.000 & 0.001 & 0.000 & \textcolor{red}{-0.000} & \textcolor{red}{-0.000} & \textcolor{red}{-0.000} & 0.000 & 0.000\\
BC & 2 & 0.002 & 0.002 & 0.002 & 0.001 & 0.001 & 0.001 & 0.000 & 0.001 & 0.001 & 0.000 & 0.000 & 0.000\\
\hline
VAR &  &  &  &  &  &  &  &  &  &  &  &  & \\
\hline
Total & 1 & \textcolor{red}{-0.042} & \textcolor{red}{-0.054} & \textcolor{red}{-0.041} & \textcolor{red}{-0.045} & \textcolor{red}{-0.027} & \textcolor{red}{-0.039} & \textcolor{red}{-0.034} & \textcolor{red}{-0.035} & \textcolor{red}{-0.022} & \textcolor{red}{-0.036} & \textcolor{red}{-0.021} & \textcolor{red}{-0.035}\\
Total & 2 & \textcolor{red}{-0.028} & \textcolor{red}{-0.038} & \textcolor{red}{-0.032} & \textcolor{red}{-0.032} & \textcolor{red}{-0.030} & \textcolor{red}{-0.031} & \textcolor{red}{-0.013} & \textcolor{red}{-0.020} & \textcolor{red}{-0.013} & \textcolor{red}{-0.020} & \textcolor{red}{-0.001} & \textcolor{red}{-0.026}\\
A & 1 & \textcolor{red}{-0.032} & \textcolor{red}{-0.034} & \textcolor{red}{-0.011} & \textcolor{red}{-0.028} & \textcolor{red}{-0.017} & \textcolor{red}{-0.026} & \textcolor{red}{-0.028} & \textcolor{red}{-0.021} & \textcolor{red}{-0.008} & \textcolor{red}{-0.030} & \textcolor{red}{-0.021} & \textcolor{red}{-0.031}\\
A & 2 & \textcolor{red}{-0.028} & \textcolor{red}{-0.034} & \textcolor{red}{-0.032} & \textcolor{red}{-0.024} & \textcolor{red}{-0.034} & \textcolor{red}{-0.040} & \textcolor{red}{-0.044} & \textcolor{red}{-0.038} & \textcolor{red}{-0.033} & \textcolor{red}{-0.031} & \textcolor{red}{-0.028} & \textcolor{red}{-0.027}\\
B & 1 & \textcolor{red}{-0.029} & \textcolor{red}{-0.044} & \textcolor{red}{-0.049} & \textcolor{red}{-0.038} & \textcolor{red}{-0.030} & \textcolor{red}{-0.038} & \textcolor{red}{-0.038} & \textcolor{red}{-0.032} & \textcolor{red}{-0.037} & \textcolor{red}{-0.029} & \textcolor{red}{-0.030} & \textcolor{red}{-0.026}\\
B & 2 & \textcolor{red}{-0.026} & \textcolor{red}{-0.026} & \textcolor{red}{-0.012} & \textcolor{red}{-0.029} & \textcolor{red}{-0.024} & \textcolor{red}{-0.026} & \textcolor{red}{-0.009} & \textcolor{red}{-0.023} & \textcolor{red}{-0.013} & \textcolor{red}{-0.021} & 0.003 & \textcolor{red}{-0.015}\\
AA & 1 & \textcolor{red}{-0.017} & \textcolor{red}{-0.005} & \textcolor{red}{-0.009} & \textcolor{red}{-0.005} & \textcolor{red}{-0.026} & \textcolor{red}{-0.011} & \textcolor{red}{-0.029} & \textcolor{red}{-0.022} & \textcolor{red}{-0.044} & \textcolor{red}{-0.026} & \textcolor{red}{-0.025} & \textcolor{red}{-0.005}\\
AA & 2 & \textcolor{red}{-0.031} & \textcolor{red}{-0.031} & \textcolor{red}{-0.029} & \textcolor{red}{-0.030} & \textcolor{red}{-0.040} & \textcolor{red}{-0.023} & \textcolor{red}{-0.026} & \textcolor{red}{-0.028} & \textcolor{red}{-0.032} & \textcolor{red}{-0.011} & \textcolor{red}{-0.019} & \textcolor{red}{-0.011}\\
AB & 1 & \textcolor{red}{-0.078} & \textcolor{red}{-0.070} & \textcolor{red}{-0.060} & \textcolor{red}{-0.070} & \textcolor{red}{-0.070} & \textcolor{red}{-0.080} & \textcolor{red}{-0.091} & \textcolor{red}{-0.059} & \textcolor{red}{-0.050} & \textcolor{red}{-0.075} & \textcolor{red}{-0.085} & \textcolor{red}{-0.074}\\
AB & 2 & \textcolor{red}{-0.038} & \textcolor{red}{-0.055} & \textcolor{red}{-0.038} & \textcolor{red}{-0.022} & \textcolor{red}{-0.026} & \textcolor{red}{-0.037} & \textcolor{red}{-0.044} & \textcolor{red}{-0.036} & \textcolor{red}{-0.024} & \textcolor{red}{-0.046} & \textcolor{red}{-0.043} & \textcolor{red}{-0.044}\\
BA & 1 & \textcolor{red}{-0.089} & \textcolor{red}{-0.091} & \textcolor{red}{-0.080} & \textcolor{red}{-0.087} & \textcolor{red}{-0.114} & \textcolor{red}{-0.136} & \textcolor{red}{-0.122} & \textcolor{red}{-0.094} & \textcolor{red}{-0.100} & \textcolor{red}{-0.092} & \textcolor{red}{-0.074} & \textcolor{red}{-0.091}\\
BA & 2 & \textcolor{red}{-0.023} & \textcolor{red}{-0.021} & \textcolor{red}{-0.026} & \textcolor{red}{-0.029} & \textcolor{red}{-0.027} & \textcolor{red}{-0.022} & \textcolor{red}{-0.009} & \textcolor{red}{-0.012} & \textcolor{red}{-0.015} & \textcolor{red}{-0.016} & \textcolor{red}{-0.009} & \textcolor{red}{-0.014}\\
BB & 1 & \textcolor{red}{-0.029} & \textcolor{red}{-0.028} & \textcolor{red}{-0.033} & \textcolor{red}{-0.022} & \textcolor{red}{-0.029} & \textcolor{red}{-0.025} & \textcolor{red}{-0.037} & \textcolor{red}{-0.020} & \textcolor{red}{-0.039} & \textcolor{red}{-0.033} & \textcolor{red}{-0.036} & \textcolor{red}{-0.023}\\
BB & 2 & \textcolor{red}{-0.026} & \textcolor{red}{-0.023} & \textcolor{red}{-0.013} & \textcolor{red}{-0.022} & \textcolor{red}{-0.016} & \textcolor{red}{-0.022} & \textcolor{red}{-0.016} & \textcolor{red}{-0.024} & \textcolor{red}{-0.020} & \textcolor{red}{-0.020} & \textcolor{red}{-0.011} & \textcolor{red}{-0.014}\\
BC & 1 & \textcolor{red}{-0.014} & \textcolor{red}{-0.016} & \textcolor{red}{-0.027} & \textcolor{red}{-0.011} & \textcolor{red}{-0.014} & \textcolor{red}{-0.010} & \textcolor{red}{-0.026} & \textcolor{red}{-0.018} & \textcolor{red}{-0.026} & \textcolor{red}{-0.011} & \textcolor{red}{-0.030} & \textcolor{red}{-0.014}\\
BC & 2 & \textcolor{red}{-0.031} & \textcolor{red}{-0.018} & \textcolor{red}{-0.009} & \textcolor{red}{-0.017} & \textcolor{red}{-0.013} & \textcolor{red}{-0.022} & \textcolor{red}{-0.008} & \textcolor{red}{-0.012} & \textcolor{red}{-0.010} & \textcolor{red}{-0.010} & \textcolor{red}{-0.001} & \textcolor{red}{-0.007}\\
\hline
    \end{tabular}%
    }
\end{table}

\begin{table}[!htb]
\caption{$\RelRMSE_{\Uni}$ for all series in the hierarchy across different forecast horizons. ARIMA, ETS, and VAR models for base forecasts. Using the shrinkage approach to estimate $\boldsymbol{W}$ under scenario 5. Values in red indicate a $\RelRMSE_{\Uni}$ less than 0.}
\centering\resizebox{1\textwidth}{!}{%
\begin{tabular}{llcccccccccccc}\hline
\multirow{2}{*}{Series} & \multirow{2}{*}{Variable}
& \multicolumn{12}{c}{Forecast horizons ($h$)} \\ \cline{3-14}
& & 1 & 2 & 3 & 4 & 5 & 6 & 7 & 8 & 9 & 10 & 11 & 12 \\ \hline
ARIMA &  &  &  &  &  &  &  &  &  &  &  &  & \\
\hline
Total & 1 & 0.006 & 0.004 & 0.002 & \textcolor{red}{-0.005} & \textcolor{red}{-0.004} & \textcolor{red}{-0.001} & \textcolor{red}{-0.002} & \textcolor{red}{-0.001} & 0.000 & 0.002 & 0.001 & 0.001\\
Total & 2 & 0.010 & 0.008 & 0.010 & 0.002 & 0.001 & 0.002 & \textcolor{red}{-0.000} & 0.003 & 0.002 & 0.001 & \textcolor{red}{-0.000} & \textcolor{red}{-0.001}\\
A & 1 & \textcolor{red}{-0.000} & 0.003 & 0.001 & 0.000 & 0.001 & 0.000 & \textcolor{red}{-0.000} & \textcolor{red}{-0.000} & 0.000 & \textcolor{red}{-0.001} & \textcolor{red}{-0.001} & \textcolor{red}{-0.002}\\
A & 2 & 0.005 & 0.006 & 0.006 & 0.005 & 0.003 & 0.005 & \textcolor{red}{-0.001} & \textcolor{red}{-0.000} & 0.001 & \textcolor{red}{-0.001} & \textcolor{red}{-0.001} & 0.000\\
B & 1 & 0.007 & 0.004 & \textcolor{red}{-0.000} & \textcolor{red}{-0.004} & \textcolor{red}{-0.005} & 0.000 & 0.000 & 0.002 & 0.001 & 0.004 & 0.002 & 0.005\\
B & 2 & 0.006 & 0.006 & 0.006 & 0.000 & 0.003 & 0.003 & 0.000 & 0.002 & 0.001 & 0.003 & 0.003 & 0.001\\
AA & 1 & 0.002 & 0.005 & 0.003 & 0.001 & 0.000 & 0.003 & \textcolor{red}{-0.001} & \textcolor{red}{-0.000} & \textcolor{red}{-0.001} & \textcolor{red}{-0.003} & \textcolor{red}{-0.002} & \textcolor{red}{-0.001}\\
AA & 2 & 0.003 & 0.005 & 0.001 & \textcolor{red}{-0.000} & 0.001 & 0.001 & \textcolor{red}{-0.001} & \textcolor{red}{-0.002} & 0.002 & \textcolor{red}{-0.002} & \textcolor{red}{-0.000} & 0.004\\
AB & 1 & 0.002 & 0.002 & 0.001 & 0.002 & \textcolor{red}{-0.001} & 0.000 & \textcolor{red}{-0.001} & \textcolor{red}{-0.003} & \textcolor{red}{-0.002} & \textcolor{red}{-0.001} & 0.003 & 0.004\\
AB & 2 & 0.005 & 0.007 & 0.005 & 0.008 & \textcolor{red}{-0.000} & 0.001 & \textcolor{red}{-0.004} & \textcolor{red}{-0.003} & 0.001 & \textcolor{red}{-0.003} & \textcolor{red}{-0.005} & \textcolor{red}{-0.001}\\
BA & 1 & 0.007 & 0.004 & 0.005 & 0.004 & \textcolor{red}{-0.001} & 0.001 & 0.001 & \textcolor{red}{-0.000} & 0.001 & 0.005 & 0.002 & 0.007\\
BA & 2 & \textcolor{red}{-0.001} & 0.005 & \textcolor{red}{-0.001} & \textcolor{red}{-0.001} & \textcolor{red}{-0.001} & 0.003 & 0.005 & 0.005 & 0.002 & 0.002 & 0.001 & 0.006\\
BB & 1 & 0.005 & 0.000 & 0.002 & 0.001 & \textcolor{red}{-0.001} & \textcolor{red}{-0.001} & 0.005 & 0.006 & 0.002 & 0.004 & 0.004 & 0.006\\
BB & 2 & 0.008 & 0.001 & 0.002 & 0.001 & 0.002 & 0.002 & 0.001 & 0.002 & \textcolor{red}{-0.001} & 0.001 & 0.001 & 0.000\\
BC & 1 & 0.007 & 0.007 & 0.008 & 0.001 & \textcolor{red}{-0.001} & 0.000 & \textcolor{red}{-0.001} & \textcolor{red}{-0.002} & \textcolor{red}{-0.004} & \textcolor{red}{-0.004} & \textcolor{red}{-0.000} & \textcolor{red}{-0.002}\\
BC & 2 & 0.002 & 0.003 & 0.003 & 0.007 & 0.006 & 0.002 & \textcolor{red}{-0.001} & \textcolor{red}{-0.002} & \textcolor{red}{-0.000} & 0.001 & \textcolor{red}{-0.000} & 0.001\\
\hline
ETS &  &  &  &  &  &  &  &  &  &  &  &  & \\
\hline
Total & 1 & 0.005 & 0.002 & 0.004 & 0.003 & 0.002 & 0.001 & \textcolor{red}{-0.000} & \textcolor{red}{-0.001} & 0.001 & 0.001 & \textcolor{red}{-0.001} & \textcolor{red}{-0.001}\\
Total & 2 & 0.005 & 0.003 & 0.002 & 0.001 & 0.001 & \textcolor{red}{-0.001} & \textcolor{red}{-0.001} & \textcolor{red}{-0.000} & 0.003 & 0.002 & 0.003 & 0.002\\
A & 1 & 0.006 & 0.004 & 0.003 & 0.002 & 0.001 & 0.001 & 0.001 & 0.001 & 0.002 & 0.002 & 0.001 & \textcolor{red}{-0.000}\\
A & 2 & 0.002 & 0.002 & 0.003 & 0.002 & 0.001 & 0.001 & 0.000 & \textcolor{red}{-0.000} & 0.003 & 0.001 & \textcolor{red}{-0.000} & 0.001\\
B & 1 & 0.004 & 0.003 & 0.003 & 0.003 & 0.001 & \textcolor{red}{-0.000} & 0.002 & 0.001 & 0.001 & 0.001 & 0.002 & 0.000\\
B & 2 & 0.006 & 0.004 & 0.004 & 0.002 & 0.002 & 0.000 & 0.001 & 0.002 & 0.004 & 0.004 & 0.005 & 0.003\\
AA & 1 & 0.000 & \textcolor{red}{-0.001} & \textcolor{red}{-0.000} & 0.002 & 0.002 & 0.002 & 0.002 & 0.001 & 0.001 & 0.002 & 0.001 & \textcolor{red}{-0.001}\\
AA & 2 & 0.002 & 0.002 & 0.002 & 0.003 & 0.001 & 0.001 & 0.002 & \textcolor{red}{-0.000} & 0.002 & 0.003 & 0.001 & 0.001\\
AB & 1 & 0.009 & 0.005 & 0.004 & 0.003 & 0.001 & 0.001 & 0.001 & \textcolor{red}{-0.000} & 0.000 & \textcolor{red}{-0.001} & \textcolor{red}{-0.001} & \textcolor{red}{-0.001}\\
AB & 2 & 0.002 & 0.003 & 0.002 & \textcolor{red}{-0.001} & 0.001 & \textcolor{red}{-0.000} & \textcolor{red}{-0.000} & \textcolor{red}{-0.001} & 0.002 & 0.001 & \textcolor{red}{-0.001} & 0.000\\
BA & 1 & 0.004 & 0.004 & 0.003 & 0.001 & 0.001 & 0.002 & 0.004 & 0.000 & 0.002 & 0.003 & 0.003 & 0.001\\
BA & 2 & 0.005 & 0.004 & 0.005 & 0.002 & 0.002 & 0.002 & 0.002 & 0.002 & 0.002 & 0.004 & 0.003 & 0.002\\
BB & 1 & 0.006 & 0.002 & 0.001 & 0.002 & 0.001 & 0.001 & 0.002 & 0.002 & 0.002 & 0.002 & 0.001 & 0.001\\
BB & 2 & 0.007 & 0.002 & 0.001 & 0.002 & \textcolor{red}{-0.000} & \textcolor{red}{-0.000} & \textcolor{red}{-0.000} & 0.000 & \textcolor{red}{-0.001} & 0.001 & 0.000 & 0.001\\
BC & 1 & 0.006 & 0.005 & 0.004 & 0.003 & 0.001 & 0.002 & 0.001 & 0.001 & 0.002 & 0.001 & \textcolor{red}{-0.001} & \textcolor{red}{-0.000}\\
BC & 2 & 0.003 & 0.002 & 0.003 & 0.002 & 0.001 & 0.000 & 0.001 & 0.001 & 0.002 & 0.002 & 0.001 & 0.001\\
\hline
VAR &  &  &  &  &  &  &  &  &  &  &  &  & \\
\hline
Total & 1 & \textcolor{red}{-0.040} & \textcolor{red}{-0.056} & \textcolor{red}{-0.053} & \textcolor{red}{-0.060} & \textcolor{red}{-0.057} & \textcolor{red}{-0.066} & \textcolor{red}{-0.045} & \textcolor{red}{-0.051} & \textcolor{red}{-0.056} & \textcolor{red}{-0.057} & \textcolor{red}{-0.039} & \textcolor{red}{-0.059}\\
Total & 2 & \textcolor{red}{-0.016} & \textcolor{red}{-0.033} & \textcolor{red}{-0.019} & \textcolor{red}{-0.039} & \textcolor{red}{-0.038} & \textcolor{red}{-0.041} & \textcolor{red}{-0.026} & \textcolor{red}{-0.039} & \textcolor{red}{-0.023} & \textcolor{red}{-0.037} & \textcolor{red}{-0.017} & \textcolor{red}{-0.022}\\
A & 1 & \textcolor{red}{-0.048} & \textcolor{red}{-0.059} & \textcolor{red}{-0.028} & \textcolor{red}{-0.043} & \textcolor{red}{-0.044} & \textcolor{red}{-0.041} & \textcolor{red}{-0.027} & \textcolor{red}{-0.020} & \textcolor{red}{-0.018} & \textcolor{red}{-0.026} & \textcolor{red}{-0.009} & \textcolor{red}{-0.023}\\
A & 2 & \textcolor{red}{-0.031} & \textcolor{red}{-0.043} & \textcolor{red}{-0.031} & \textcolor{red}{-0.026} & \textcolor{red}{-0.022} & \textcolor{red}{-0.027} & \textcolor{red}{-0.033} & \textcolor{red}{-0.029} & \textcolor{red}{-0.040} & \textcolor{red}{-0.043} & \textcolor{red}{-0.050} & \textcolor{red}{-0.039}\\
B & 1 & \textcolor{red}{-0.033} & \textcolor{red}{-0.031} & \textcolor{red}{-0.036} & \textcolor{red}{-0.035} & \textcolor{red}{-0.047} & \textcolor{red}{-0.051} & \textcolor{red}{-0.042} & \textcolor{red}{-0.048} & \textcolor{red}{-0.058} & \textcolor{red}{-0.048} & \textcolor{red}{-0.051} & \textcolor{red}{-0.044}\\
B & 2 & \textcolor{red}{-0.020} & \textcolor{red}{-0.023} & \textcolor{red}{-0.011} & \textcolor{red}{-0.027} & \textcolor{red}{-0.026} & \textcolor{red}{-0.036} & \textcolor{red}{-0.032} & \textcolor{red}{-0.027} & \textcolor{red}{-0.013} & \textcolor{red}{-0.038} & \textcolor{red}{-0.013} & \textcolor{red}{-0.018}\\
AA & 1 & \textcolor{red}{-0.021} & \textcolor{red}{-0.013} & \textcolor{red}{-0.017} & \textcolor{red}{-0.006} & \textcolor{red}{-0.010} & \textcolor{red}{-0.005} & \textcolor{red}{-0.012} & \textcolor{red}{-0.002} & \textcolor{red}{-0.030} & \textcolor{red}{-0.023} & \textcolor{red}{-0.056} & \textcolor{red}{-0.026}\\
AA & 2 & \textcolor{red}{-0.021} & \textcolor{red}{-0.033} & \textcolor{red}{-0.041} & \textcolor{red}{-0.037} & \textcolor{red}{-0.042} & \textcolor{red}{-0.048} & \textcolor{red}{-0.048} & \textcolor{red}{-0.035} & \textcolor{red}{-0.047} & \textcolor{red}{-0.049} & \textcolor{red}{-0.056} & \textcolor{red}{-0.049}\\
AB & 1 & \textcolor{red}{-0.091} & \textcolor{red}{-0.108} & \textcolor{red}{-0.091} & \textcolor{red}{-0.099} & \textcolor{red}{-0.101} & \textcolor{red}{-0.131} & \textcolor{red}{-0.100} & \textcolor{red}{-0.069} & \textcolor{red}{-0.082} & \textcolor{red}{-0.090} & \textcolor{red}{-0.095} & \textcolor{red}{-0.073}\\
AB & 2 & \textcolor{red}{-0.043} & \textcolor{red}{-0.072} & \textcolor{red}{-0.049} & \textcolor{red}{-0.027} & \textcolor{red}{-0.033} & \textcolor{red}{-0.064} & \textcolor{red}{-0.058} & \textcolor{red}{-0.042} & \textcolor{red}{-0.051} & \textcolor{red}{-0.071} & \textcolor{red}{-0.062} & \textcolor{red}{-0.047}\\
BA & 1 & \textcolor{red}{-0.095} & \textcolor{red}{-0.125} & \textcolor{red}{-0.127} & \textcolor{red}{-0.118} & \textcolor{red}{-0.155} & \textcolor{red}{-0.154} & \textcolor{red}{-0.140} & \textcolor{red}{-0.129} & \textcolor{red}{-0.161} & \textcolor{red}{-0.197} & \textcolor{red}{-0.176} & \textcolor{red}{-0.145}\\
BA & 2 & \textcolor{red}{-0.020} & \textcolor{red}{-0.031} & \textcolor{red}{-0.026} & \textcolor{red}{-0.025} & \textcolor{red}{-0.031} & \textcolor{red}{-0.030} & \textcolor{red}{-0.032} & \textcolor{red}{-0.018} & \textcolor{red}{-0.012} & \textcolor{red}{-0.011} & 0.001 & \textcolor{red}{-0.011}\\
BB & 1 & \textcolor{red}{-0.044} & \textcolor{red}{-0.029} & \textcolor{red}{-0.024} & \textcolor{red}{-0.027} & \textcolor{red}{-0.042} & \textcolor{red}{-0.039} & \textcolor{red}{-0.063} & \textcolor{red}{-0.046} & \textcolor{red}{-0.053} & \textcolor{red}{-0.039} & \textcolor{red}{-0.052} & \textcolor{red}{-0.028}\\
BB & 2 & \textcolor{red}{-0.012} & \textcolor{red}{-0.027} & \textcolor{red}{-0.021} & \textcolor{red}{-0.015} & \textcolor{red}{-0.022} & \textcolor{red}{-0.027} & \textcolor{red}{-0.027} & \textcolor{red}{-0.009} & \textcolor{red}{-0.007} & \textcolor{red}{-0.024} & \textcolor{red}{-0.015} & \textcolor{red}{-0.011}\\
BC & 1 & \textcolor{red}{-0.026} & \textcolor{red}{-0.018} & \textcolor{red}{-0.016} & \textcolor{red}{-0.017} & \textcolor{red}{-0.027} & \textcolor{red}{-0.031} & \textcolor{red}{-0.034} & \textcolor{red}{-0.028} & \textcolor{red}{-0.027} & \textcolor{red}{-0.006} & \textcolor{red}{-0.024} & \textcolor{red}{-0.006}\\
BC & 2 & \textcolor{red}{-0.018} & \textcolor{red}{-0.017} & \textcolor{red}{-0.010} & \textcolor{red}{-0.030} & \textcolor{red}{-0.030} & \textcolor{red}{-0.030} & \textcolor{red}{-0.015} & \textcolor{red}{-0.024} & \textcolor{red}{-0.018} & \textcolor{red}{-0.023} & \textcolor{red}{-0.004} & \textcolor{red}{-0.013}\\
\hline
    \end{tabular}%
    }
\end{table}

\begin{table}[!htb]
\caption{$\RelRMSE_{\Uni}$ for all series in the hierarchy across different forecast horizons. ARIMA, ETS, and VAR models for base forecasts. Using the shrinkage approach to estimate $\boldsymbol{W}$ under scenario 6. Values in red indicate a $\RelRMSE_{\Uni}$ less than 0.}
\centering\resizebox{1\textwidth}{!}{%
\begin{tabular}{llcccccccccccc}\hline
\multirow{2}{*}{Series} & \multirow{2}{*}{Variable}
& \multicolumn{12}{c}{Forecast horizons ($h$)} \\ \cline{3-14}
& & 1 & 2 & 3 & 4 & 5 & 6 & 7 & 8 & 9 & 10 & 11 & 12 \\ \hline
ARIMA &  &  &  &  &  &  &  &  &  &  &  &  & \\
\hline
Total & 1 & 0.011 & 0.012 & 0.011 & 0.002 & 0.003 & \textcolor{red}{-0.000} & 0.000 & \textcolor{red}{-0.001} & \textcolor{red}{-0.003} & \textcolor{red}{-0.005} & \textcolor{red}{-0.006} & \textcolor{red}{-0.008}\\
Total & 2 & 0.016 & 0.012 & 0.011 & 0.003 & 0.008 & 0.001 & 0.002 & 0.002 & 0.003 & \textcolor{red}{-0.002} & \textcolor{red}{-0.002} & \textcolor{red}{-0.007}\\
A & 1 & 0.007 & 0.013 & 0.007 & 0.003 & 0.005 & 0.000 & \textcolor{red}{-0.001} & 0.001 & \textcolor{red}{-0.003} & \textcolor{red}{-0.009} & \textcolor{red}{-0.008} & \textcolor{red}{-0.007}\\
A & 2 & 0.010 & 0.013 & 0.014 & 0.010 & 0.011 & 0.006 & 0.006 & 0.003 & 0.002 & 0.001 & 0.003 & \textcolor{red}{-0.002}\\
B & 1 & 0.017 & 0.009 & 0.011 & 0.004 & 0.008 & 0.005 & 0.005 & 0.004 & 0.003 & 0.000 & \textcolor{red}{-0.001} & \textcolor{red}{-0.000}\\
B & 2 & 0.011 & 0.003 & 0.004 & \textcolor{red}{-0.001} & 0.002 & \textcolor{red}{-0.003} & \textcolor{red}{-0.003} & \textcolor{red}{-0.004} & \textcolor{red}{-0.004} & \textcolor{red}{-0.006} & \textcolor{red}{-0.004} & \textcolor{red}{-0.008}\\
AA & 1 & 0.006 & 0.013 & 0.004 & \textcolor{red}{-0.004} & 0.003 & 0.003 & 0.000 & \textcolor{red}{-0.001} & 0.000 & \textcolor{red}{-0.003} & \textcolor{red}{-0.005} & \textcolor{red}{-0.008}\\
AA & 2 & 0.006 & 0.009 & 0.010 & 0.009 & 0.012 & 0.005 & 0.002 & 0.003 & 0.003 & 0.001 & 0.001 & \textcolor{red}{-0.000}\\
AB & 1 & 0.015 & 0.019 & 0.016 & 0.008 & 0.014 & 0.008 & 0.004 & 0.002 & 0.002 & \textcolor{red}{-0.001} & \textcolor{red}{-0.002} & \textcolor{red}{-0.002}\\
AB & 2 & 0.009 & 0.011 & 0.011 & 0.006 & 0.008 & 0.005 & 0.003 & 0.001 & 0.001 & \textcolor{red}{-0.002} & 0.001 & \textcolor{red}{-0.003}\\
BA & 1 & 0.017 & 0.010 & 0.014 & 0.005 & 0.008 & 0.006 & 0.007 & 0.003 & 0.005 & 0.003 & 0.006 & 0.003\\
BA & 2 & 0.006 & 0.002 & 0.005 & 0.002 & 0.003 & \textcolor{red}{-0.000} & 0.004 & 0.004 & 0.002 & 0.000 & 0.008 & 0.006\\
BB & 1 & 0.007 & 0.005 & 0.009 & 0.004 & 0.004 & \textcolor{red}{-0.003} & 0.002 & 0.002 & 0.003 & \textcolor{red}{-0.001} & 0.001 & 0.000\\
BB & 2 & 0.013 & 0.007 & 0.001 & \textcolor{red}{-0.000} & 0.004 & 0.004 & 0.000 & \textcolor{red}{-0.004} & \textcolor{red}{-0.002} & \textcolor{red}{-0.001} & \textcolor{red}{-0.004} & \textcolor{red}{-0.010}\\
BC & 1 & 0.008 & 0.004 & 0.004 & 0.005 & 0.007 & 0.009 & 0.007 & 0.002 & 0.002 & 0.004 & 0.003 & 0.000\\
BC & 2 & 0.009 & 0.003 & 0.005 & 0.001 & 0.004 & 0.004 & 0.005 & 0.002 & 0.002 & 0.001 & 0.001 & \textcolor{red}{-0.004}\\
\hline
ETS &  &  &  &  &  &  &  &  &  &  &  &  & \\
\hline
Total & 1 & 0.006 & 0.005 & 0.008 & 0.007 & 0.006 & 0.004 & 0.004 & 0.004 & 0.004 & 0.004 & 0.004 & 0.003\\
Total & 2 & 0.006 & 0.006 & 0.007 & 0.004 & 0.004 & 0.002 & 0.002 & 0.002 & 0.002 & 0.003 & 0.002 & 0.002\\
A & 1 & 0.006 & 0.005 & 0.005 & 0.004 & 0.004 & 0.003 & 0.002 & 0.003 & 0.003 & 0.003 & 0.002 & 0.003\\
A & 2 & 0.004 & 0.004 & 0.004 & 0.002 & 0.003 & 0.003 & 0.002 & 0.002 & 0.002 & 0.002 & 0.001 & 0.001\\
B & 1 & 0.006 & 0.005 & 0.006 & 0.005 & 0.005 & 0.003 & 0.003 & 0.002 & 0.002 & 0.002 & 0.002 & 0.001\\
B & 2 & 0.006 & 0.005 & 0.006 & 0.005 & 0.005 & 0.001 & 0.001 & 0.001 & 0.001 & 0.000 & \textcolor{red}{-0.000} & 0.001\\
AA & 1 & 0.006 & 0.005 & 0.003 & 0.001 & 0.002 & 0.002 & 0.001 & 0.002 & 0.001 & 0.002 & 0.001 & 0.001\\
AA & 2 & 0.003 & 0.002 & 0.002 & 0.000 & 0.001 & 0.002 & 0.001 & 0.000 & 0.001 & 0.001 & 0.000 & 0.000\\
AB & 1 & 0.004 & 0.004 & 0.004 & 0.003 & 0.003 & 0.002 & 0.001 & 0.002 & 0.002 & 0.002 & 0.001 & 0.001\\
AB & 2 & 0.004 & 0.004 & 0.003 & 0.001 & 0.001 & 0.002 & 0.001 & 0.001 & 0.001 & 0.001 & 0.000 & \textcolor{red}{-0.000}\\
BA & 1 & 0.003 & 0.003 & 0.004 & 0.003 & 0.003 & 0.002 & 0.002 & 0.002 & 0.002 & 0.001 & 0.002 & 0.002\\
BA & 2 & 0.004 & 0.003 & 0.004 & 0.003 & 0.003 & 0.002 & 0.002 & 0.002 & 0.002 & 0.001 & 0.001 & 0.001\\
BB & 1 & 0.005 & 0.004 & 0.003 & 0.002 & 0.002 & 0.001 & 0.001 & 0.001 & 0.001 & 0.001 & 0.000 & 0.000\\
BB & 2 & 0.003 & 0.001 & 0.002 & 0.000 & 0.000 & \textcolor{red}{-0.000} & 0.001 & \textcolor{red}{-0.001} & \textcolor{red}{-0.001} & \textcolor{red}{-0.001} & \textcolor{red}{-0.001} & \textcolor{red}{-0.001}\\
BC & 1 & 0.003 & 0.001 & 0.003 & 0.001 & 0.001 & 0.000 & 0.001 & 0.000 & 0.000 & \textcolor{red}{-0.000} & 0.000 & \textcolor{red}{-0.000}\\
BC & 2 & 0.004 & 0.004 & 0.003 & 0.002 & 0.002 & 0.001 & \textcolor{red}{-0.000} & \textcolor{red}{-0.000} & \textcolor{red}{-0.000} & 0.000 & 0.000 & \textcolor{red}{-0.000}\\
\hline
VAR &  &  &  &  &  &  &  &  &  &  &  &  & \\
\hline
Total & 1 & \textcolor{red}{-0.052} & \textcolor{red}{-0.071} & \textcolor{red}{-0.045} & \textcolor{red}{-0.031} & \textcolor{red}{-0.024} & \textcolor{red}{-0.057} & \textcolor{red}{-0.044} & \textcolor{red}{-0.033} & \textcolor{red}{-0.007} & \textcolor{red}{-0.025} & \textcolor{red}{-0.016} & \textcolor{red}{-0.027}\\
Total & 2 & \textcolor{red}{-0.039} & \textcolor{red}{-0.045} & \textcolor{red}{-0.038} & \textcolor{red}{-0.038} & \textcolor{red}{-0.026} & \textcolor{red}{-0.036} & \textcolor{red}{-0.019} & \textcolor{red}{-0.024} & \textcolor{red}{-0.005} & \textcolor{red}{-0.017} & \textcolor{red}{-0.005} & \textcolor{red}{-0.014}\\
A & 1 & \textcolor{red}{-0.033} & \textcolor{red}{-0.037} & \textcolor{red}{-0.034} & \textcolor{red}{-0.022} & \textcolor{red}{-0.008} & \textcolor{red}{-0.028} & \textcolor{red}{-0.021} & \textcolor{red}{-0.025} & \textcolor{red}{-0.008} & \textcolor{red}{-0.015} & \textcolor{red}{-0.005} & \textcolor{red}{-0.012}\\
A & 2 & \textcolor{red}{-0.043} & \textcolor{red}{-0.042} & \textcolor{red}{-0.027} & \textcolor{red}{-0.033} & \textcolor{red}{-0.052} & \textcolor{red}{-0.050} & \textcolor{red}{-0.043} & \textcolor{red}{-0.033} & \textcolor{red}{-0.038} & \textcolor{red}{-0.020} & \textcolor{red}{-0.013} & \textcolor{red}{-0.014}\\
B & 1 & \textcolor{red}{-0.044} & \textcolor{red}{-0.059} & \textcolor{red}{-0.040} & \textcolor{red}{-0.033} & \textcolor{red}{-0.038} & \textcolor{red}{-0.059} & \textcolor{red}{-0.046} & \textcolor{red}{-0.033} & \textcolor{red}{-0.025} & \textcolor{red}{-0.028} & \textcolor{red}{-0.031} & \textcolor{red}{-0.029}\\
B & 2 & \textcolor{red}{-0.031} & \textcolor{red}{-0.046} & \textcolor{red}{-0.033} & \textcolor{red}{-0.030} & \textcolor{red}{-0.018} & \textcolor{red}{-0.030} & \textcolor{red}{-0.005} & \textcolor{red}{-0.011} & \textcolor{red}{-0.001} & \textcolor{red}{-0.011} & 0.003 & \textcolor{red}{-0.012}\\
AA & 1 & \textcolor{red}{-0.018} & \textcolor{red}{-0.026} & \textcolor{red}{-0.030} & \textcolor{red}{-0.013} & \textcolor{red}{-0.019} & \textcolor{red}{-0.021} & \textcolor{red}{-0.037} & \textcolor{red}{-0.021} & \textcolor{red}{-0.028} & \textcolor{red}{-0.003} & \textcolor{red}{-0.020} & \textcolor{red}{-0.011}\\
AA & 2 & \textcolor{red}{-0.035} & \textcolor{red}{-0.035} & \textcolor{red}{-0.024} & \textcolor{red}{-0.029} & \textcolor{red}{-0.032} & \textcolor{red}{-0.021} & \textcolor{red}{-0.035} & \textcolor{red}{-0.021} & \textcolor{red}{-0.017} & \textcolor{red}{-0.016} & \textcolor{red}{-0.021} & \textcolor{red}{-0.006}\\
AB & 1 & \textcolor{red}{-0.070} & \textcolor{red}{-0.091} & \textcolor{red}{-0.082} & \textcolor{red}{-0.044} & \textcolor{red}{-0.050} & \textcolor{red}{-0.072} & \textcolor{red}{-0.075} & \textcolor{red}{-0.060} & \textcolor{red}{-0.053} & \textcolor{red}{-0.055} & \textcolor{red}{-0.052} & \textcolor{red}{-0.041}\\
AB & 2 & \textcolor{red}{-0.055} & \textcolor{red}{-0.069} & \textcolor{red}{-0.030} & \textcolor{red}{-0.033} & \textcolor{red}{-0.063} & \textcolor{red}{-0.060} & \textcolor{red}{-0.036} & \textcolor{red}{-0.029} & \textcolor{red}{-0.047} & \textcolor{red}{-0.034} & \textcolor{red}{-0.025} & \textcolor{red}{-0.036}\\
BA & 1 & \textcolor{red}{-0.097} & \textcolor{red}{-0.115} & \textcolor{red}{-0.099} & \textcolor{red}{-0.082} & \textcolor{red}{-0.103} & \textcolor{red}{-0.140} & \textcolor{red}{-0.099} & \textcolor{red}{-0.079} & \textcolor{red}{-0.067} & \textcolor{red}{-0.081} & \textcolor{red}{-0.075} & \textcolor{red}{-0.069}\\
BA & 2 & \textcolor{red}{-0.029} & \textcolor{red}{-0.037} & \textcolor{red}{-0.033} & \textcolor{red}{-0.029} & \textcolor{red}{-0.023} & \textcolor{red}{-0.021} & \textcolor{red}{-0.010} & \textcolor{red}{-0.010} & \textcolor{red}{-0.007} & \textcolor{red}{-0.011} & \textcolor{red}{-0.005} & \textcolor{red}{-0.006}\\
BB & 1 & \textcolor{red}{-0.034} & \textcolor{red}{-0.036} & \textcolor{red}{-0.041} & \textcolor{red}{-0.030} & \textcolor{red}{-0.046} & \textcolor{red}{-0.040} & \textcolor{red}{-0.046} & \textcolor{red}{-0.038} & \textcolor{red}{-0.036} & \textcolor{red}{-0.007} & \textcolor{red}{-0.033} & \textcolor{red}{-0.026}\\
BB & 2 & \textcolor{red}{-0.013} & \textcolor{red}{-0.030} & \textcolor{red}{-0.037} & \textcolor{red}{-0.030} & \textcolor{red}{-0.022} & \textcolor{red}{-0.031} & \textcolor{red}{-0.022} & \textcolor{red}{-0.020} & \textcolor{red}{-0.014} & \textcolor{red}{-0.018} & \textcolor{red}{-0.008} & \textcolor{red}{-0.012}\\
BC & 1 & \textcolor{red}{-0.026} & \textcolor{red}{-0.028} & \textcolor{red}{-0.018} & \textcolor{red}{-0.010} & \textcolor{red}{-0.015} & \textcolor{red}{-0.010} & \textcolor{red}{-0.014} & \textcolor{red}{-0.008} & \textcolor{red}{-0.025} & \textcolor{red}{-0.007} & \textcolor{red}{-0.018} & \textcolor{red}{-0.007}\\
BC & 2 & \textcolor{red}{-0.026} & \textcolor{red}{-0.025} & \textcolor{red}{-0.029} & \textcolor{red}{-0.013} & \textcolor{red}{-0.013} & \textcolor{red}{-0.018} & \textcolor{red}{-0.005} & \textcolor{red}{-0.004} & \textcolor{red}{-0.002} & \textcolor{red}{-0.008} & \textcolor{red}{-0.006} & \textcolor{red}{-0.005}\\
\hline
    \end{tabular}%
    }
\end{table}

\begin{table}[!htb]
\caption{$\RelRMSE_{\Uni}$ for all series in the hierarchy across different forecast horizons. ARIMA, ETS, and VAR models for base forecasts. Using the shrinkage approach to estimate $\boldsymbol{W}$ under scenario 7. Values in red indicate a $\RelRMSE_{\Uni}$ less than 0.}
\centering\resizebox{1\textwidth}{!}{%
\begin{tabular}{llcccccccccccc}\hline
\multirow{2}{*}{Series} & \multirow{2}{*}{Variable}
& \multicolumn{12}{c}{Forecast horizons ($h$)} \\ \cline{3-14}
& & 1 & 2 & 3 & 4 & 5 & 6 & 7 & 8 & 9 & 10 & 11 & 12 \\ \hline
ARIMA &  &  &  &  &  &  &  &  &  &  &  &  & \\
\hline
Total & 1 & 0.005 & 0.002 & 0.004 & 0.000 & 0.000 & \textcolor{red}{-0.005} & 0.001 & \textcolor{red}{-0.002} & \textcolor{red}{-0.000} & \textcolor{red}{-0.003} & \textcolor{red}{-0.002} & \textcolor{red}{-0.004}\\
Total & 2 & 0.005 & 0.003 & 0.001 & \textcolor{red}{-0.001} & 0.002 & \textcolor{red}{-0.001} & 0.000 & 0.000 & 0.001 & 0.001 & 0.000 & \textcolor{red}{-0.002}\\
A & 1 & 0.004 & 0.008 & 0.005 & 0.003 & 0.004 & 0.002 & 0.005 & 0.003 & 0.001 & 0.002 & 0.001 & 0.001\\
A & 2 & 0.004 & 0.002 & 0.003 & 0.001 & 0.003 & 0.004 & 0.002 & 0.002 & 0.001 & 0.003 & 0.001 & 0.001\\
B & 1 & 0.004 & 0.000 & 0.002 & \textcolor{red}{-0.001} & 0.000 & \textcolor{red}{-0.003} & \textcolor{red}{-0.001} & \textcolor{red}{-0.002} & 0.000 & \textcolor{red}{-0.002} & \textcolor{red}{-0.000} & \textcolor{red}{-0.005}\\
B & 2 & 0.004 & 0.001 & 0.000 & 0.002 & 0.005 & 0.001 & \textcolor{red}{-0.001} & 0.002 & 0.001 & 0.001 & \textcolor{red}{-0.002} & \textcolor{red}{-0.003}\\
AA & 1 & 0.003 & 0.007 & 0.002 & 0.004 & 0.004 & 0.002 & 0.003 & 0.001 & 0.002 & 0.002 & 0.001 & 0.002\\
AA & 2 & 0.003 & 0.002 & 0.004 & 0.001 & 0.005 & 0.004 & 0.002 & 0.002 & 0.003 & 0.002 & 0.001 & 0.001\\
AB & 1 & \textcolor{red}{-0.000} & 0.001 & 0.003 & 0.001 & 0.004 & \textcolor{red}{-0.000} & 0.002 & 0.003 & 0.001 & 0.000 & 0.001 & 0.001\\
AB & 2 & \textcolor{red}{-0.000} & 0.002 & 0.001 & 0.002 & 0.001 & 0.003 & 0.002 & 0.001 & \textcolor{red}{-0.000} & 0.003 & 0.001 & \textcolor{red}{-0.002}\\
BA & 1 & 0.003 & \textcolor{red}{-0.003} & \textcolor{red}{-0.001} & 0.000 & 0.002 & \textcolor{red}{-0.003} & \textcolor{red}{-0.002} & \textcolor{red}{-0.001} & 0.001 & \textcolor{red}{-0.003} & \textcolor{red}{-0.001} & \textcolor{red}{-0.002}\\
BA & 2 & 0.002 & 0.000 & 0.002 & 0.003 & 0.002 & 0.000 & 0.001 & 0.003 & 0.001 & 0.002 & 0.001 & 0.001\\
BB & 1 & 0.005 & 0.002 & \textcolor{red}{-0.000} & \textcolor{red}{-0.001} & 0.001 & \textcolor{red}{-0.001} & 0.002 & 0.001 & 0.001 & 0.003 & 0.001 & \textcolor{red}{-0.002}\\
BB & 2 & 0.001 & 0.000 & \textcolor{red}{-0.001} & \textcolor{red}{-0.002} & \textcolor{red}{-0.000} & \textcolor{red}{-0.002} & \textcolor{red}{-0.001} & 0.001 & 0.002 & 0.000 & \textcolor{red}{-0.000} & \textcolor{red}{-0.003}\\
BC & 1 & \textcolor{red}{-0.000} & 0.003 & 0.002 & \textcolor{red}{-0.001} & 0.001 & 0.000 & 0.001 & 0.000 & 0.002 & 0.000 & 0.000 & \textcolor{red}{-0.004}\\
BC & 2 & 0.001 & 0.004 & \textcolor{red}{-0.000} & 0.000 & 0.001 & 0.002 & 0.000 & 0.002 & 0.001 & 0.002 & 0.001 & \textcolor{red}{-0.000}\\
\hline
ETS &  &  &  &  &  &  &  &  &  &  &  &  & \\
\hline
Total & 1 & 0.002 & 0.001 & 0.000 & 0.001 & 0.000 & \textcolor{red}{-0.000} & \textcolor{red}{-0.000} & 0.000 & 0.000 & 0.000 & 0.000 & \textcolor{red}{-0.000}\\
Total & 2 & 0.002 & 0.003 & 0.001 & 0.001 & 0.000 & 0.000 & 0.001 & 0.001 & 0.001 & 0.001 & 0.001 & \textcolor{red}{-0.000}\\
A & 1 & 0.001 & 0.001 & 0.001 & 0.001 & 0.001 & 0.000 & 0.001 & 0.000 & 0.000 & 0.000 & 0.000 & \textcolor{red}{-0.000}\\
A & 2 & 0.000 & 0.001 & 0.001 & 0.001 & 0.000 & 0.000 & 0.001 & 0.000 & 0.001 & 0.001 & 0.000 & \textcolor{red}{-0.000}\\
B & 1 & 0.001 & 0.000 & 0.000 & \textcolor{red}{-0.000} & \textcolor{red}{-0.000} & \textcolor{red}{-0.000} & 0.000 & 0.000 & 0.000 & 0.000 & 0.000 & 0.000\\
B & 2 & 0.001 & 0.002 & 0.001 & 0.001 & 0.000 & \textcolor{red}{-0.000} & 0.000 & 0.000 & 0.000 & 0.000 & 0.000 & \textcolor{red}{-0.000}\\
AA & 1 & 0.001 & 0.001 & 0.001 & 0.001 & \textcolor{red}{-0.000} & 0.000 & 0.000 & 0.000 & 0.001 & 0.001 & 0.000 & 0.000\\
AA & 2 & 0.001 & 0.001 & 0.001 & 0.000 & \textcolor{red}{-0.000} & 0.000 & 0.000 & 0.000 & 0.001 & 0.000 & \textcolor{red}{-0.000} & \textcolor{red}{-0.000}\\
AB & 1 & 0.001 & 0.001 & 0.001 & 0.001 & \textcolor{red}{-0.000} & \textcolor{red}{-0.000} & 0.000 & 0.000 & 0.000 & 0.000 & 0.000 & \textcolor{red}{-0.000}\\
AB & 2 & 0.000 & \textcolor{red}{-0.000} & 0.000 & \textcolor{red}{-0.000} & \textcolor{red}{-0.000} & \textcolor{red}{-0.001} & \textcolor{red}{-0.000} & \textcolor{red}{-0.000} & \textcolor{red}{-0.000} & \textcolor{red}{-0.000} & \textcolor{red}{-0.000} & \textcolor{red}{-0.001}\\
BA & 1 & 0.001 & 0.001 & 0.001 & 0.001 & 0.000 & 0.000 & 0.000 & 0.000 & 0.001 & 0.000 & 0.000 & \textcolor{red}{-0.000}\\
BA & 2 & 0.001 & 0.003 & 0.002 & 0.001 & 0.001 & 0.001 & 0.001 & 0.001 & 0.001 & 0.001 & 0.001 & 0.000\\
BB & 1 & 0.001 & 0.001 & 0.000 & 0.001 & 0.001 & 0.000 & 0.001 & 0.000 & 0.001 & 0.001 & 0.001 & 0.000\\
BB & 2 & 0.001 & 0.002 & 0.002 & 0.001 & 0.001 & 0.001 & 0.001 & 0.001 & 0.001 & 0.001 & 0.001 & 0.000\\
BC & 1 & 0.002 & 0.002 & 0.001 & 0.000 & 0.001 & 0.001 & 0.001 & 0.001 & 0.001 & 0.001 & 0.000 & 0.001\\
BC & 2 & 0.001 & 0.002 & 0.001 & 0.000 & 0.001 & 0.001 & 0.001 & 0.000 & 0.001 & 0.000 & \textcolor{red}{-0.000} & \textcolor{red}{-0.000}\\
\hline
VAR &  &  &  &  &  &  &  &  &  &  &  &  & \\
\hline
Total & 1 & \textcolor{red}{-0.021} & \textcolor{red}{-0.033} & \textcolor{red}{-0.006} & \textcolor{red}{-0.011} & \textcolor{red}{-0.014} & \textcolor{red}{-0.030} & \textcolor{red}{-0.007} & \textcolor{red}{-0.017} & \textcolor{red}{-0.013} & \textcolor{red}{-0.026} & \textcolor{red}{-0.000} & \textcolor{red}{-0.012}\\
Total & 2 & \textcolor{red}{-0.006} & \textcolor{red}{-0.019} & \textcolor{red}{-0.006} & \textcolor{red}{-0.012} & \textcolor{red}{-0.000} & \textcolor{red}{-0.006} & 0.002 & \textcolor{red}{-0.010} & 0.000 & \textcolor{red}{-0.007} & 0.006 & \textcolor{red}{-0.003}\\
A & 1 & \textcolor{red}{-0.017} & \textcolor{red}{-0.017} & \textcolor{red}{-0.004} & \textcolor{red}{-0.006} & 0.002 & \textcolor{red}{-0.006} & 0.008 & \textcolor{red}{-0.003} & 0.004 & \textcolor{red}{-0.013} & 0.003 & \textcolor{red}{-0.008}\\
A & 2 & \textcolor{red}{-0.014} & \textcolor{red}{-0.021} & \textcolor{red}{-0.012} & \textcolor{red}{-0.009} & \textcolor{red}{-0.020} & \textcolor{red}{-0.025} & \textcolor{red}{-0.019} & \textcolor{red}{-0.021} & \textcolor{red}{-0.026} & \textcolor{red}{-0.020} & \textcolor{red}{-0.020} & \textcolor{red}{-0.016}\\
B & 1 & \textcolor{red}{-0.022} & \textcolor{red}{-0.022} & \textcolor{red}{-0.009} & \textcolor{red}{-0.012} & \textcolor{red}{-0.023} & \textcolor{red}{-0.032} & \textcolor{red}{-0.019} & \textcolor{red}{-0.015} & \textcolor{red}{-0.023} & \textcolor{red}{-0.017} & \textcolor{red}{-0.007} & \textcolor{red}{-0.003}\\
B & 2 & \textcolor{red}{-0.007} & \textcolor{red}{-0.015} & \textcolor{red}{-0.005} & \textcolor{red}{-0.007} & 0.006 & \textcolor{red}{-0.004} & 0.006 & \textcolor{red}{-0.003} & 0.008 & \textcolor{red}{-0.006} & 0.011 & 0.000\\
AA & 1 & \textcolor{red}{-0.008} & \textcolor{red}{-0.004} & \textcolor{red}{-0.017} & \textcolor{red}{-0.005} & \textcolor{red}{-0.012} & \textcolor{red}{-0.002} & \textcolor{red}{-0.017} & \textcolor{red}{-0.008} & \textcolor{red}{-0.024} & \textcolor{red}{-0.005} & \textcolor{red}{-0.018} & \textcolor{red}{-0.001}\\
AA & 2 & \textcolor{red}{-0.020} & \textcolor{red}{-0.030} & \textcolor{red}{-0.026} & \textcolor{red}{-0.016} & \textcolor{red}{-0.014} & \textcolor{red}{-0.016} & \textcolor{red}{-0.023} & \textcolor{red}{-0.019} & \textcolor{red}{-0.027} & \textcolor{red}{-0.024} & \textcolor{red}{-0.024} & \textcolor{red}{-0.008}\\
AB & 1 & \textcolor{red}{-0.047} & \textcolor{red}{-0.052} & \textcolor{red}{-0.051} & \textcolor{red}{-0.038} & \textcolor{red}{-0.041} & \textcolor{red}{-0.055} & \textcolor{red}{-0.036} & \textcolor{red}{-0.029} & \textcolor{red}{-0.032} & \textcolor{red}{-0.052} & \textcolor{red}{-0.038} & \textcolor{red}{-0.028}\\
AB & 2 & \textcolor{red}{-0.015} & \textcolor{red}{-0.036} & \textcolor{red}{-0.027} & \textcolor{red}{-0.024} & \textcolor{red}{-0.033} & \textcolor{red}{-0.043} & \textcolor{red}{-0.031} & \textcolor{red}{-0.033} & \textcolor{red}{-0.038} & \textcolor{red}{-0.029} & \textcolor{red}{-0.025} & \textcolor{red}{-0.022}\\
BA & 1 & \textcolor{red}{-0.050} & \textcolor{red}{-0.058} & \textcolor{red}{-0.056} & \textcolor{red}{-0.054} & \textcolor{red}{-0.072} & \textcolor{red}{-0.099} & \textcolor{red}{-0.080} & \textcolor{red}{-0.072} & \textcolor{red}{-0.080} & \textcolor{red}{-0.082} & \textcolor{red}{-0.060} & \textcolor{red}{-0.041}\\
BA & 2 & \textcolor{red}{-0.003} & \textcolor{red}{-0.011} & \textcolor{red}{-0.012} & \textcolor{red}{-0.005} & 0.001 & \textcolor{red}{-0.001} & \textcolor{red}{-0.001} & \textcolor{red}{-0.003} & \textcolor{red}{-0.001} & \textcolor{red}{-0.001} & \textcolor{red}{-0.005} & \textcolor{red}{-0.007}\\
BB & 1 & \textcolor{red}{-0.023} & \textcolor{red}{-0.016} & \textcolor{red}{-0.026} & \textcolor{red}{-0.025} & \textcolor{red}{-0.025} & \textcolor{red}{-0.012} & \textcolor{red}{-0.030} & \textcolor{red}{-0.018} & \textcolor{red}{-0.030} & \textcolor{red}{-0.011} & \textcolor{red}{-0.025} & \textcolor{red}{-0.013}\\
BB & 2 & \textcolor{red}{-0.008} & \textcolor{red}{-0.011} & 0.002 & \textcolor{red}{-0.003} & \textcolor{red}{-0.002} & \textcolor{red}{-0.006} & 0.000 & \textcolor{red}{-0.002} & 0.000 & \textcolor{red}{-0.004} & 0.001 & 0.003\\
BC & 1 & \textcolor{red}{-0.020} & \textcolor{red}{-0.014} & \textcolor{red}{-0.015} & \textcolor{red}{-0.011} & \textcolor{red}{-0.013} & \textcolor{red}{-0.013} & \textcolor{red}{-0.016} & \textcolor{red}{-0.006} & \textcolor{red}{-0.016} & \textcolor{red}{-0.008} & \textcolor{red}{-0.019} & \textcolor{red}{-0.005}\\
BC & 2 & \textcolor{red}{-0.009} & \textcolor{red}{-0.013} & \textcolor{red}{-0.010} & \textcolor{red}{-0.007} & \textcolor{red}{-0.002} & \textcolor{red}{-0.008} & 0.000 & 0.003 & \textcolor{red}{-0.000} & \textcolor{red}{-0.002} & 0.003 & 0.002\\
\hline
    \end{tabular}%
    }
\end{table}

\begin{table}[!htb]
\caption{$\RelRMSE_{\Uni}$ for all series in the hierarchy across different forecast horizons. ARIMA, ETS, and VAR models for base forecasts. Using the shrinkage approach to estimate $\boldsymbol{W}$ under scenario 8. Values in red indicate a $\RelRMSE_{\Uni}$ less than 0.}
\centering\resizebox{1\textwidth}{!}{%
\begin{tabular}{llcccccccccccc}\hline
\multirow{2}{*}{Series} & \multirow{2}{*}{Variable}
& \multicolumn{12}{c}{Forecast horizons ($h$)} \\ \cline{3-14}
& & 1 & 2 & 3 & 4 & 5 & 6 & 7 & 8 & 9 & 10 & 11 & 12 \\ \hline
ARIMA &  &  &  &  &  &  &  &  &  &  &  &  & \\
\hline
Total & 1 & 0.004 & 0.003 & 0.004 & 0.003 & 0.004 & 0.004 & 0.001 & 0.002 & 0.003 & 0.004 & 0.004 & 0.003\\
Total & 2 & 0.006 & 0.005 & 0.004 & 0.005 & 0.005 & 0.004 & \textcolor{red}{-0.002} & \textcolor{red}{-0.003} & \textcolor{red}{-0.001} & 0.001 & 0.003 & 0.007\\
A & 1 & 0.011 & 0.005 & 0.006 & 0.003 & 0.007 & 0.004 & 0.005 & 0.006 & 0.005 & 0.008 & 0.006 & 0.002\\
A & 2 & 0.007 & 0.007 & 0.004 & 0.004 & 0.004 & 0.002 & \textcolor{red}{-0.000} & \textcolor{red}{-0.000} & 0.004 & 0.003 & 0.004 & 0.005\\
B & 1 & 0.001 & 0.004 & 0.005 & 0.004 & 0.006 & 0.006 & 0.004 & 0.007 & 0.006 & 0.008 & 0.009 & 0.008\\
B & 2 & \textcolor{red}{-0.002} & 0.001 & 0.001 & 0.005 & 0.004 & 0.005 & 0.003 & 0.002 & 0.003 & 0.003 & 0.002 & 0.006\\
AA & 1 & 0.006 & 0.004 & 0.004 & 0.001 & 0.004 & 0.004 & 0.003 & 0.006 & 0.005 & 0.010 & 0.007 & 0.007\\
AA & 2 & 0.005 & 0.005 & 0.004 & 0.001 & 0.003 & 0.005 & 0.007 & 0.010 & 0.008 & 0.010 & 0.010 & 0.008\\
AB & 1 & 0.010 & 0.006 & 0.006 & 0.002 & 0.007 & 0.002 & 0.004 & 0.007 & 0.009 & 0.009 & 0.007 & 0.004\\
AB & 2 & 0.006 & 0.006 & 0.002 & 0.005 & 0.002 & 0.001 & 0.002 & \textcolor{red}{-0.002} & 0.001 & \textcolor{red}{-0.000} & \textcolor{red}{-0.000} & 0.001\\
BA & 1 & 0.001 & 0.004 & 0.002 & 0.003 & 0.003 & 0.000 & 0.004 & 0.006 & 0.004 & 0.007 & 0.004 & 0.003\\
BA & 2 & 0.001 & 0.004 & 0.002 & 0.001 & 0.007 & 0.004 & \textcolor{red}{-0.000} & 0.001 & 0.004 & 0.001 & 0.002 & 0.003\\
BB & 1 & \textcolor{red}{-0.002} & 0.001 & 0.001 & \textcolor{red}{-0.000} & \textcolor{red}{-0.001} & \textcolor{red}{-0.001} & 0.000 & 0.003 & 0.002 & 0.002 & 0.003 & 0.003\\
BB & 2 & \textcolor{red}{-0.003} & 0.001 & 0.001 & 0.007 & 0.004 & 0.005 & 0.004 & 0.005 & 0.005 & 0.004 & \textcolor{red}{-0.001} & 0.001\\
BC & 1 & 0.000 & 0.001 & 0.003 & 0.001 & 0.004 & 0.003 & 0.001 & 0.005 & 0.006 & 0.007 & 0.003 & 0.003\\
BC & 2 & 0.003 & 0.001 & 0.000 & 0.000 & 0.005 & 0.004 & 0.004 & 0.001 & 0.004 & 0.003 & 0.002 & 0.004\\
\hline
ETS &  &  &  &  &  &  &  &  &  &  &  &  & \\
\hline
Total & 1 & 0.005 & 0.003 & 0.002 & 0.002 & 0.000 & \textcolor{red}{-0.000} & 0.001 & 0.002 & 0.002 & 0.004 & 0.002 & 0.002\\
Total & 2 & 0.004 & 0.003 & 0.001 & 0.001 & \textcolor{red}{-0.000} & \textcolor{red}{-0.002} & \textcolor{red}{-0.001} & 0.001 & 0.002 & 0.002 & 0.001 & 0.003\\
A & 1 & 0.005 & 0.002 & 0.002 & 0.001 & 0.002 & 0.002 & 0.002 & 0.002 & 0.002 & 0.003 & 0.002 & 0.002\\
A & 2 & 0.004 & 0.003 & 0.001 & 0.001 & 0.003 & 0.002 & \textcolor{red}{-0.000} & 0.000 & 0.001 & 0.001 & 0.001 & 0.004\\
B & 1 & 0.003 & 0.004 & 0.002 & 0.003 & 0.001 & \textcolor{red}{-0.000} & \textcolor{red}{-0.000} & 0.002 & 0.001 & 0.003 & 0.002 & 0.002\\
B & 2 & 0.004 & 0.004 & 0.001 & 0.002 & 0.001 & \textcolor{red}{-0.001} & \textcolor{red}{-0.000} & 0.001 & 0.001 & 0.001 & 0.000 & 0.001\\
AA & 1 & 0.004 & 0.002 & 0.002 & 0.002 & 0.001 & 0.002 & 0.002 & 0.001 & 0.002 & 0.002 & 0.002 & 0.000\\
AA & 2 & 0.003 & 0.002 & \textcolor{red}{-0.000} & 0.001 & 0.001 & 0.000 & \textcolor{red}{-0.000} & \textcolor{red}{-0.000} & 0.001 & 0.000 & 0.000 & 0.001\\
AB & 1 & 0.003 & 0.002 & 0.001 & \textcolor{red}{-0.000} & 0.001 & 0.001 & \textcolor{red}{-0.000} & 0.001 & 0.002 & 0.003 & 0.003 & 0.002\\
AB & 2 & 0.002 & 0.002 & \textcolor{red}{-0.000} & 0.001 & 0.001 & 0.001 & \textcolor{red}{-0.000} & 0.001 & 0.001 & 0.002 & 0.002 & 0.004\\
BA & 1 & 0.002 & 0.003 & 0.002 & 0.001 & 0.001 & \textcolor{red}{-0.000} & 0.001 & 0.001 & 0.001 & 0.002 & 0.001 & 0.001\\
BA & 2 & 0.001 & 0.002 & 0.002 & 0.000 & 0.001 & 0.001 & 0.001 & 0.000 & 0.000 & 0.001 & \textcolor{red}{-0.001} & 0.000\\
BB & 1 & 0.003 & 0.001 & 0.002 & 0.001 & \textcolor{red}{-0.001} & \textcolor{red}{-0.000} & 0.001 & 0.001 & \textcolor{red}{-0.000} & 0.002 & 0.002 & 0.001\\
BB & 2 & 0.003 & 0.002 & \textcolor{red}{-0.001} & 0.000 & 0.000 & 0.000 & 0.000 & 0.001 & 0.001 & 0.001 & 0.001 & \textcolor{red}{-0.001}\\
BC & 1 & 0.002 & 0.002 & 0.001 & 0.001 & 0.000 & \textcolor{red}{-0.001} & \textcolor{red}{-0.001} & 0.001 & 0.000 & 0.001 & 0.001 & 0.000\\
BC & 2 & 0.003 & 0.003 & 0.003 & 0.002 & 0.002 & 0.000 & 0.002 & 0.002 & 0.001 & 0.001 & 0.002 & 0.001\\
\hline
VAR &  &  &  &  &  &  &  &  &  &  &  &  & \\
\hline
Total & 1 & \textcolor{red}{-0.021} & \textcolor{red}{-0.044} & \textcolor{red}{-0.036} & \textcolor{red}{-0.020} & \textcolor{red}{-0.017} & \textcolor{red}{-0.041} & \textcolor{red}{-0.034} & \textcolor{red}{-0.039} & \textcolor{red}{-0.028} & \textcolor{red}{-0.047} & \textcolor{red}{-0.020} & \textcolor{red}{-0.029}\\
Total & 2 & \textcolor{red}{-0.008} & \textcolor{red}{-0.019} & \textcolor{red}{-0.011} & \textcolor{red}{-0.015} & 0.000 & \textcolor{red}{-0.013} & \textcolor{red}{-0.004} & \textcolor{red}{-0.009} & \textcolor{red}{-0.004} & \textcolor{red}{-0.017} & \textcolor{red}{-0.000} & \textcolor{red}{-0.012}\\
A & 1 & \textcolor{red}{-0.009} & \textcolor{red}{-0.020} & \textcolor{red}{-0.015} & \textcolor{red}{-0.017} & \textcolor{red}{-0.007} & \textcolor{red}{-0.029} & \textcolor{red}{-0.016} & \textcolor{red}{-0.025} & \textcolor{red}{-0.010} & \textcolor{red}{-0.026} & \textcolor{red}{-0.005} & \textcolor{red}{-0.016}\\
A & 2 & \textcolor{red}{-0.011} & \textcolor{red}{-0.023} & \textcolor{red}{-0.026} & \textcolor{red}{-0.020} & \textcolor{red}{-0.017} & \textcolor{red}{-0.022} & \textcolor{red}{-0.025} & \textcolor{red}{-0.023} & \textcolor{red}{-0.028} & \textcolor{red}{-0.027} & \textcolor{red}{-0.034} & \textcolor{red}{-0.027}\\
B & 1 & \textcolor{red}{-0.015} & \textcolor{red}{-0.035} & \textcolor{red}{-0.027} & \textcolor{red}{-0.011} & \textcolor{red}{-0.022} & \textcolor{red}{-0.032} & \textcolor{red}{-0.039} & \textcolor{red}{-0.035} & \textcolor{red}{-0.042} & \textcolor{red}{-0.044} & \textcolor{red}{-0.032} & \textcolor{red}{-0.026}\\
B & 2 & \textcolor{red}{-0.010} & \textcolor{red}{-0.017} & 0.005 & \textcolor{red}{-0.001} & 0.000 & \textcolor{red}{-0.008} & 0.014 & 0.007 & 0.017 & 0.001 & 0.018 & 0.007\\
AA & 1 & \textcolor{red}{-0.021} & \textcolor{red}{-0.010} & \textcolor{red}{-0.008} & \textcolor{red}{-0.008} & \textcolor{red}{-0.020} & \textcolor{red}{-0.002} & \textcolor{red}{-0.006} & \textcolor{red}{-0.006} & \textcolor{red}{-0.018} & 0.002 & \textcolor{red}{-0.015} & \textcolor{red}{-0.009}\\
AA & 2 & \textcolor{red}{-0.019} & \textcolor{red}{-0.022} & \textcolor{red}{-0.027} & \textcolor{red}{-0.024} & \textcolor{red}{-0.026} & \textcolor{red}{-0.027} & \textcolor{red}{-0.034} & \textcolor{red}{-0.027} & \textcolor{red}{-0.036} & \textcolor{red}{-0.028} & \textcolor{red}{-0.034} & \textcolor{red}{-0.031}\\
AB & 1 & \textcolor{red}{-0.035} & \textcolor{red}{-0.063} & \textcolor{red}{-0.060} & \textcolor{red}{-0.039} & \textcolor{red}{-0.044} & \textcolor{red}{-0.064} & \textcolor{red}{-0.056} & \textcolor{red}{-0.051} & \textcolor{red}{-0.056} & \textcolor{red}{-0.074} & \textcolor{red}{-0.053} & \textcolor{red}{-0.050}\\
AB & 2 & \textcolor{red}{-0.035} & \textcolor{red}{-0.049} & \textcolor{red}{-0.018} & \textcolor{red}{-0.014} & \textcolor{red}{-0.018} & \textcolor{red}{-0.043} & \textcolor{red}{-0.027} & \textcolor{red}{-0.026} & \textcolor{red}{-0.044} & \textcolor{red}{-0.047} & \textcolor{red}{-0.043} & \textcolor{red}{-0.023}\\
BA & 1 & \textcolor{red}{-0.069} & \textcolor{red}{-0.086} & \textcolor{red}{-0.062} & \textcolor{red}{-0.055} & \textcolor{red}{-0.081} & \textcolor{red}{-0.094} & \textcolor{red}{-0.094} & \textcolor{red}{-0.072} & \textcolor{red}{-0.087} & \textcolor{red}{-0.097} & \textcolor{red}{-0.066} & \textcolor{red}{-0.063}\\
BA & 2 & \textcolor{red}{-0.008} & \textcolor{red}{-0.012} & \textcolor{red}{-0.007} & \textcolor{red}{-0.007} & 0.005 & 0.007 & 0.004 & 0.004 & 0.010 & 0.004 & \textcolor{red}{-0.000} & \textcolor{red}{-0.002}\\
BB & 1 & \textcolor{red}{-0.019} & \textcolor{red}{-0.011} & \textcolor{red}{-0.016} & \textcolor{red}{-0.013} & \textcolor{red}{-0.004} & 0.000 & \textcolor{red}{-0.031} & \textcolor{red}{-0.024} & \textcolor{red}{-0.023} & \textcolor{red}{-0.020} & \textcolor{red}{-0.034} & \textcolor{red}{-0.022}\\
BB & 2 & \textcolor{red}{-0.009} & \textcolor{red}{-0.016} & \textcolor{red}{-0.004} & \textcolor{red}{-0.009} & \textcolor{red}{-0.017} & \textcolor{red}{-0.017} & \textcolor{red}{-0.003} & \textcolor{red}{-0.005} & \textcolor{red}{-0.006} & \textcolor{red}{-0.015} & \textcolor{red}{-0.004} & \textcolor{red}{-0.006}\\
BC & 1 & \textcolor{red}{-0.009} & \textcolor{red}{-0.011} & \textcolor{red}{-0.003} & 0.002 & \textcolor{red}{-0.012} & \textcolor{red}{-0.011} & \textcolor{red}{-0.021} & \textcolor{red}{-0.017} & \textcolor{red}{-0.024} & \textcolor{red}{-0.019} & \textcolor{red}{-0.023} & \textcolor{red}{-0.009}\\
BC & 2 & \textcolor{red}{-0.010} & \textcolor{red}{-0.005} & \textcolor{red}{-0.002} & \textcolor{red}{-0.005} & \textcolor{red}{-0.011} & \textcolor{red}{-0.009} & \textcolor{red}{-0.004} & \textcolor{red}{-0.006} & \textcolor{red}{-0.003} & \textcolor{red}{-0.006} & 0.005 & 0.001\\
\hline
    \end{tabular}%
    }
\end{table}

\begin{table}[!htb]
\caption{$\RelRMSE_{\Uni}$ for all series in the hierarchy across different forecast horizons. ARIMA, ETS, and VAR models for base forecasts. Using the shrinkage approach to estimate $\boldsymbol{W}$ under scenario 9. Values in red indicate a $\RelRMSE_{\Uni}$ less than 0.}
\centering\resizebox{1\textwidth}{!}{%
\begin{tabular}{llcccccccccccc}\hline
\multirow{2}{*}{Series} & \multirow{2}{*}{Variable}
& \multicolumn{12}{c}{Forecast horizons ($h$)} \\ \cline{3-14}
& & 1 & 2 & 3 & 4 & 5 & 6 & 7 & 8 & 9 & 10 & 11 & 12 \\ \hline
ARIMA &  &  &  &  &  &  &  &  &  &  &  &  & \\
\hline
Total & 1 & 0.002 & 0.004 & 0.002 & \textcolor{red}{-0.001} & \textcolor{red}{-0.001} & \textcolor{red}{-0.001} & \textcolor{red}{-0.004} & \textcolor{red}{-0.009} & \textcolor{red}{-0.006} & \textcolor{red}{-0.007} & \textcolor{red}{-0.009} & \textcolor{red}{-0.012}\\
Total & 2 & \textcolor{red}{-0.003} & \textcolor{red}{-0.005} & \textcolor{red}{-0.002} & \textcolor{red}{-0.003} & \textcolor{red}{-0.005} & \textcolor{red}{-0.006} & \textcolor{red}{-0.003} & \textcolor{red}{-0.005} & \textcolor{red}{-0.008} & \textcolor{red}{-0.010} & \textcolor{red}{-0.008} & \textcolor{red}{-0.008}\\
A & 1 & 0.003 & 0.004 & 0.003 & 0.002 & 0.000 & \textcolor{red}{-0.001} & \textcolor{red}{-0.003} & \textcolor{red}{-0.003} & \textcolor{red}{-0.003} & \textcolor{red}{-0.004} & \textcolor{red}{-0.006} & \textcolor{red}{-0.006}\\
A & 2 & 0.001 & 0.000 & 0.001 & \textcolor{red}{-0.001} & 0.002 & \textcolor{red}{-0.001} & 0.001 & \textcolor{red}{-0.003} & \textcolor{red}{-0.006} & \textcolor{red}{-0.005} & \textcolor{red}{-0.005} & \textcolor{red}{-0.006}\\
B & 1 & 0.002 & 0.005 & 0.002 & 0.000 & 0.001 & 0.001 & \textcolor{red}{-0.003} & \textcolor{red}{-0.006} & \textcolor{red}{-0.005} & \textcolor{red}{-0.006} & \textcolor{red}{-0.008} & \textcolor{red}{-0.010}\\
B & 2 & \textcolor{red}{-0.000} & \textcolor{red}{-0.003} & 0.001 & \textcolor{red}{-0.000} & 0.001 & \textcolor{red}{-0.004} & \textcolor{red}{-0.005} & \textcolor{red}{-0.005} & \textcolor{red}{-0.005} & \textcolor{red}{-0.008} & \textcolor{red}{-0.009} & \textcolor{red}{-0.008}\\
AA & 1 & 0.002 & 0.002 & 0.003 & 0.002 & 0.001 & \textcolor{red}{-0.001} & \textcolor{red}{-0.002} & \textcolor{red}{-0.001} & \textcolor{red}{-0.002} & \textcolor{red}{-0.003} & \textcolor{red}{-0.003} & \textcolor{red}{-0.003}\\
AA & 2 & 0.004 & 0.007 & 0.004 & 0.000 & 0.004 & 0.001 & 0.003 & \textcolor{red}{-0.001} & \textcolor{red}{-0.002} & \textcolor{red}{-0.001} & \textcolor{red}{-0.003} & \textcolor{red}{-0.004}\\
AB & 1 & 0.000 & 0.004 & 0.003 & 0.002 & 0.001 & 0.000 & \textcolor{red}{-0.002} & \textcolor{red}{-0.001} & \textcolor{red}{-0.003} & \textcolor{red}{-0.002} & \textcolor{red}{-0.003} & \textcolor{red}{-0.004}\\
AB & 2 & \textcolor{red}{-0.001} & \textcolor{red}{-0.003} & \textcolor{red}{-0.002} & \textcolor{red}{-0.003} & 0.001 & \textcolor{red}{-0.001} & \textcolor{red}{-0.001} & \textcolor{red}{-0.004} & \textcolor{red}{-0.003} & \textcolor{red}{-0.007} & \textcolor{red}{-0.006} & \textcolor{red}{-0.008}\\
BA & 1 & 0.002 & 0.003 & 0.002 & 0.002 & 0.002 & 0.003 & 0.002 & \textcolor{red}{-0.001} & \textcolor{red}{-0.001} & \textcolor{red}{-0.002} & \textcolor{red}{-0.000} & \textcolor{red}{-0.002}\\
BA & 2 & 0.003 & 0.001 & \textcolor{red}{-0.001} & \textcolor{red}{-0.001} & \textcolor{red}{-0.000} & \textcolor{red}{-0.002} & \textcolor{red}{-0.000} & \textcolor{red}{-0.005} & \textcolor{red}{-0.003} & \textcolor{red}{-0.004} & \textcolor{red}{-0.004} & \textcolor{red}{-0.003}\\
BB & 1 & 0.001 & 0.003 & 0.003 & 0.002 & 0.002 & 0.003 & 0.003 & 0.000 & 0.001 & \textcolor{red}{-0.001} & \textcolor{red}{-0.003} & \textcolor{red}{-0.004}\\
BB & 2 & \textcolor{red}{-0.003} & \textcolor{red}{-0.002} & 0.001 & 0.001 & 0.003 & 0.002 & 0.001 & 0.000 & 0.000 & \textcolor{red}{-0.000} & \textcolor{red}{-0.001} & \textcolor{red}{-0.000}\\
BC & 1 & 0.006 & 0.002 & 0.003 & \textcolor{red}{-0.001} & 0.001 & \textcolor{red}{-0.000} & \textcolor{red}{-0.001} & \textcolor{red}{-0.001} & \textcolor{red}{-0.000} & \textcolor{red}{-0.002} & \textcolor{red}{-0.003} & \textcolor{red}{-0.006}\\
BC & 2 & 0.001 & \textcolor{red}{-0.003} & 0.001 & \textcolor{red}{-0.001} & \textcolor{red}{-0.001} & \textcolor{red}{-0.003} & \textcolor{red}{-0.003} & \textcolor{red}{-0.002} & \textcolor{red}{-0.003} & \textcolor{red}{-0.004} & \textcolor{red}{-0.003} & \textcolor{red}{-0.002}\\
\hline
ETS &  &  &  &  &  &  &  &  &  &  &  &  & \\
\hline
Total & 1 & 0.001 & 0.003 & 0.002 & 0.002 & 0.002 & 0.001 & 0.001 & 0.001 & 0.001 & 0.001 & 0.001 & 0.001\\
Total & 2 & 0.001 & 0.001 & 0.001 & 0.000 & 0.000 & 0.000 & 0.000 & 0.001 & 0.001 & 0.000 & \textcolor{red}{-0.000} & \textcolor{red}{-0.000}\\
A & 1 & 0.002 & 0.003 & 0.002 & 0.001 & 0.001 & 0.000 & 0.001 & 0.001 & 0.001 & 0.001 & 0.000 & 0.000\\
A & 2 & 0.002 & 0.002 & 0.003 & 0.002 & 0.002 & 0.001 & 0.001 & 0.001 & 0.001 & 0.001 & 0.001 & 0.000\\
B & 1 & \textcolor{red}{-0.000} & 0.001 & 0.001 & 0.001 & 0.001 & 0.001 & 0.001 & 0.001 & 0.001 & 0.001 & 0.001 & 0.001\\
B & 2 & 0.001 & 0.001 & 0.001 & 0.000 & 0.001 & 0.001 & 0.000 & 0.000 & 0.000 & 0.000 & \textcolor{red}{-0.000} & \textcolor{red}{-0.000}\\
AA & 1 & 0.001 & 0.001 & 0.001 & 0.001 & 0.001 & 0.000 & 0.001 & 0.001 & 0.001 & 0.000 & 0.001 & 0.000\\
AA & 2 & 0.001 & 0.002 & 0.002 & 0.001 & 0.001 & 0.000 & 0.001 & 0.001 & 0.001 & 0.000 & 0.000 & 0.000\\
AB & 1 & \textcolor{red}{-0.000} & 0.001 & 0.001 & 0.000 & 0.001 & 0.000 & 0.000 & 0.000 & 0.000 & 0.000 & \textcolor{red}{-0.000} & \textcolor{red}{-0.000}\\
AB & 2 & 0.000 & 0.000 & 0.001 & 0.001 & 0.001 & \textcolor{red}{-0.000} & 0.000 & 0.001 & 0.001 & 0.000 & 0.000 & 0.000\\
BA & 1 & \textcolor{red}{-0.001} & 0.001 & 0.002 & 0.002 & 0.002 & 0.001 & 0.001 & 0.001 & 0.001 & 0.001 & 0.001 & 0.001\\
BA & 2 & 0.002 & 0.001 & 0.001 & 0.001 & 0.001 & 0.000 & 0.000 & 0.000 & 0.001 & 0.000 & \textcolor{red}{-0.000} & 0.000\\
BB & 1 & 0.001 & 0.000 & \textcolor{red}{-0.000} & 0.001 & 0.000 & 0.000 & 0.001 & 0.000 & 0.000 & 0.000 & \textcolor{red}{-0.000} & 0.000\\
BB & 2 & 0.000 & 0.000 & 0.000 & 0.000 & 0.001 & 0.000 & 0.000 & 0.000 & 0.000 & 0.000 & 0.000 & 0.000\\
BC & 1 & 0.001 & 0.002 & 0.000 & 0.001 & 0.001 & 0.002 & 0.002 & 0.001 & 0.001 & 0.001 & 0.001 & 0.001\\
BC & 2 & 0.001 & 0.001 & 0.001 & 0.000 & 0.000 & 0.001 & 0.000 & 0.000 & 0.000 & 0.000 & \textcolor{red}{-0.000} & \textcolor{red}{-0.000}\\
\hline
VAR &  &  &  &  &  &  &  &  &  &  &  &  & \\
\hline
Total & 1 & \textcolor{red}{-0.035} & \textcolor{red}{-0.026} & \textcolor{red}{-0.022} & \textcolor{red}{-0.030} & \textcolor{red}{-0.024} & \textcolor{red}{-0.038} & \textcolor{red}{-0.016} & \textcolor{red}{-0.007} & \textcolor{red}{-0.004} & \textcolor{red}{-0.017} & 0.009 & \textcolor{red}{-0.009}\\
Total & 2 & \textcolor{red}{-0.015} & \textcolor{red}{-0.021} & \textcolor{red}{-0.017} & \textcolor{red}{-0.008} & \textcolor{red}{-0.003} & \textcolor{red}{-0.017} & \textcolor{red}{-0.005} & 0.002 & 0.007 & \textcolor{red}{-0.007} & 0.012 & 0.002\\
A & 1 & \textcolor{red}{-0.035} & \textcolor{red}{-0.031} & \textcolor{red}{-0.020} & \textcolor{red}{-0.018} & \textcolor{red}{-0.016} & \textcolor{red}{-0.033} & \textcolor{red}{-0.011} & \textcolor{red}{-0.011} & \textcolor{red}{-0.005} & \textcolor{red}{-0.018} & 0.001 & \textcolor{red}{-0.002}\\
A & 2 & \textcolor{red}{-0.032} & \textcolor{red}{-0.022} & \textcolor{red}{-0.017} & \textcolor{red}{-0.015} & \textcolor{red}{-0.035} & \textcolor{red}{-0.028} & \textcolor{red}{-0.023} & \textcolor{red}{-0.018} & \textcolor{red}{-0.020} & \textcolor{red}{-0.009} & \textcolor{red}{-0.015} & \textcolor{red}{-0.008}\\
B & 1 & \textcolor{red}{-0.019} & \textcolor{red}{-0.017} & \textcolor{red}{-0.027} & \textcolor{red}{-0.025} & \textcolor{red}{-0.020} & \textcolor{red}{-0.024} & \textcolor{red}{-0.026} & \textcolor{red}{-0.014} & \textcolor{red}{-0.012} & \textcolor{red}{-0.016} & \textcolor{red}{-0.014} & \textcolor{red}{-0.010}\\
B & 2 & \textcolor{red}{-0.010} & \textcolor{red}{-0.016} & \textcolor{red}{-0.008} & \textcolor{red}{-0.006} & 0.003 & \textcolor{red}{-0.017} & \textcolor{red}{-0.004} & \textcolor{red}{-0.000} & 0.009 & \textcolor{red}{-0.009} & 0.017 & 0.000\\
AA & 1 & 0.006 & 0.006 & \textcolor{red}{-0.006} & 0.001 & \textcolor{red}{-0.005} & 0.003 & \textcolor{red}{-0.004} & 0.002 & \textcolor{red}{-0.015} & 0.004 & \textcolor{red}{-0.014} & 0.001\\
AA & 2 & \textcolor{red}{-0.025} & \textcolor{red}{-0.010} & \textcolor{red}{-0.013} & \textcolor{red}{-0.009} & \textcolor{red}{-0.022} & \textcolor{red}{-0.026} & \textcolor{red}{-0.030} & \textcolor{red}{-0.010} & \textcolor{red}{-0.014} & \textcolor{red}{-0.010} & \textcolor{red}{-0.010} & 0.001\\
AB & 1 & \textcolor{red}{-0.074} & \textcolor{red}{-0.092} & \textcolor{red}{-0.074} & \textcolor{red}{-0.064} & \textcolor{red}{-0.082} & \textcolor{red}{-0.081} & \textcolor{red}{-0.060} & \textcolor{red}{-0.051} & \textcolor{red}{-0.052} & \textcolor{red}{-0.058} & \textcolor{red}{-0.046} & \textcolor{red}{-0.036}\\
AB & 2 & \textcolor{red}{-0.033} & \textcolor{red}{-0.037} & \textcolor{red}{-0.025} & \textcolor{red}{-0.028} & \textcolor{red}{-0.039} & \textcolor{red}{-0.038} & \textcolor{red}{-0.025} & \textcolor{red}{-0.034} & \textcolor{red}{-0.026} & \textcolor{red}{-0.015} & \textcolor{red}{-0.020} & \textcolor{red}{-0.028}\\
BA & 1 & \textcolor{red}{-0.059} & \textcolor{red}{-0.099} & \textcolor{red}{-0.085} & \textcolor{red}{-0.082} & \textcolor{red}{-0.086} & \textcolor{red}{-0.101} & \textcolor{red}{-0.091} & \textcolor{red}{-0.069} & \textcolor{red}{-0.067} & \textcolor{red}{-0.074} & \textcolor{red}{-0.057} & \textcolor{red}{-0.049}\\
BA & 2 & \textcolor{red}{-0.013} & \textcolor{red}{-0.010} & \textcolor{red}{-0.003} & 0.003 & 0.004 & \textcolor{red}{-0.002} & 0.006 & 0.000 & \textcolor{red}{-0.004} & \textcolor{red}{-0.007} & 0.001 & \textcolor{red}{-0.006}\\
BB & 1 & \textcolor{red}{-0.029} & \textcolor{red}{-0.023} & \textcolor{red}{-0.028} & \textcolor{red}{-0.020} & \textcolor{red}{-0.033} & \textcolor{red}{-0.017} & \textcolor{red}{-0.029} & \textcolor{red}{-0.026} & \textcolor{red}{-0.046} & \textcolor{red}{-0.019} & \textcolor{red}{-0.034} & \textcolor{red}{-0.017}\\
BB & 2 & \textcolor{red}{-0.008} & \textcolor{red}{-0.003} & \textcolor{red}{-0.009} & \textcolor{red}{-0.008} & \textcolor{red}{-0.006} & \textcolor{red}{-0.013} & \textcolor{red}{-0.011} & \textcolor{red}{-0.005} & 0.002 & \textcolor{red}{-0.009} & 0.005 & 0.004\\
BC & 1 & \textcolor{red}{-0.016} & \textcolor{red}{-0.004} & \textcolor{red}{-0.020} & \textcolor{red}{-0.013} & \textcolor{red}{-0.011} & 0.001 & \textcolor{red}{-0.009} & 0.002 & \textcolor{red}{-0.008} & \textcolor{red}{-0.002} & \textcolor{red}{-0.018} & \textcolor{red}{-0.002}\\
BC & 2 & \textcolor{red}{-0.004} & \textcolor{red}{-0.019} & \textcolor{red}{-0.020} & \textcolor{red}{-0.017} & \textcolor{red}{-0.011} & \textcolor{red}{-0.015} & \textcolor{red}{-0.011} & \textcolor{red}{-0.013} & \textcolor{red}{-0.006} & \textcolor{red}{-0.012} & \textcolor{red}{-0.004} & \textcolor{red}{-0.010}\\
\hline
    \end{tabular}%
    }
\end{table}

\begin{table}[!htb]
  \caption{$\RelRMSE^{\Base}$ of admission series. VAR model for base forecasts and the shrinkage approach to estimate $\bm{W}$. Values in red indicate a $\RelRMSE^{\Base}$ less than 0.}
  \centering
  \resizebox{1\textwidth}{!}{%
    \begin{tabular}{lrrrrrrrrrrrr}
      \hline
      \multirow{2}{*}{Series} & \multicolumn{12}{c}{Forecast horizons ($h$)} \\ \cline{2-13}
      & 1 & 2 & 3 & 4 & 5 & 6 & 7 & 8 & 9 & 10 & 11 & 12 \\ \hline
      Total & \textcolor{red}{-0.020} & \textcolor{red}{-0.089} & \textcolor{red}{-0.049} & \textcolor{red}{-0.080} & \textcolor{red}{-0.118} & \textcolor{red}{-0.176} & \textcolor{red}{-0.147} & \textcolor{red}{-0.087} & \textcolor{red}{-0.062} & \textcolor{red}{-0.124} & \textcolor{red}{-0.212} & \textcolor{red}{-0.402} \\
      Midwest & \textcolor{red}{-0.045} & \textcolor{red}{-0.065} & \textcolor{red}{-0.035} & \textcolor{red}{-0.039} & \textcolor{red}{-0.083} & \textcolor{red}{-0.107} & \textcolor{red}{-0.091} & \textcolor{red}{-0.079} & \textcolor{red}{-0.063} & \textcolor{red}{-0.172} & \textcolor{red}{-0.178} & \textcolor{red}{-0.235} \\
      Northeast & 0.033 & \textcolor{red}{-0.047} & \textcolor{red}{-0.029} & \textcolor{red}{-0.109} & \textcolor{red}{-0.119} & \textcolor{red}{-0.160} & \textcolor{red}{-0.160} & \textcolor{red}{-0.145} & \textcolor{red}{-0.141} & \textcolor{red}{-0.173} & \textcolor{red}{-0.245} & \textcolor{red}{-0.363} \\
      North & 0.019 & \textcolor{red}{-0.029} & \textcolor{red}{-0.028} & \textcolor{red}{-0.031} & \textcolor{red}{-0.070} & \textcolor{red}{-0.081} & \textcolor{red}{-0.052} & \textcolor{red}{-0.042} & \textcolor{red}{-0.023} & \textcolor{red}{-0.141} & \textcolor{red}{-0.142} & \textcolor{red}{-0.123} \\
      Southeast & \textcolor{red}{-0.003} & \textcolor{red}{-0.065} & \textcolor{red}{-0.026} & \textcolor{red}{-0.044} & \textcolor{red}{-0.102} & \textcolor{red}{-0.169} & \textcolor{red}{-0.153} & \textcolor{red}{-0.131} & \textcolor{red}{-0.094} & \textcolor{red}{-0.171} & \textcolor{red}{-0.249} & \textcolor{red}{-0.403} \\
      South & \textcolor{red}{-0.033} & \textcolor{red}{-0.067} & 0.003 & 0.014 & \textcolor{red}{-0.003} & \textcolor{red}{-0.061} & \textcolor{red}{-0.060} & \textcolor{red}{-0.043} & \textcolor{red}{-0.035} & \textcolor{red}{-0.089} & \textcolor{red}{-0.172} & \textcolor{red}{-0.290} \\
      AC & \textcolor{red}{-0.019} & \textcolor{red}{-0.021} & \textcolor{red}{-0.017} & \textcolor{red}{-0.021} & \textcolor{red}{-0.020} & \textcolor{red}{-0.024} & \textcolor{red}{-0.022} & \textcolor{red}{-0.020} & \textcolor{red}{-0.025} & \textcolor{red}{-0.037} & \textcolor{red}{-0.063} & \textcolor{red}{-0.140} \\
      AL & 0.115 & 0.094 & 0.100 & 0.141 & 0.154 & 0.195 & 0.203 & 0.220 & 0.190 & 0.159 & 0.149 & 0.225 \\
      AM & \textcolor{red}{-0.034} & \textcolor{red}{-0.077} & \textcolor{red}{-0.086} & \textcolor{red}{-0.119} & \textcolor{red}{-0.157} & \textcolor{red}{-0.211} & \textcolor{red}{-0.216} & \textcolor{red}{-0.205} & \textcolor{red}{-0.177} & \textcolor{red}{-0.210} & \textcolor{red}{-0.287} & \textcolor{red}{-0.306} \\
      AP & \textcolor{red}{-0.016} & \textcolor{red}{-0.034} & \textcolor{red}{-0.038} & \textcolor{red}{-0.052} & \textcolor{red}{-0.057} & \textcolor{red}{-0.064} & \textcolor{red}{-0.070} & \textcolor{red}{-0.082} & \textcolor{red}{-0.091} & \textcolor{red}{-0.107} & \textcolor{red}{-0.113} & \textcolor{red}{-0.099} \\
      BA & \textcolor{red}{-0.009} & \textcolor{red}{-0.023} & \textcolor{red}{-0.026} & \textcolor{red}{-0.033} & \textcolor{red}{-0.064} & \textcolor{red}{-0.077} & \textcolor{red}{-0.077} & \textcolor{red}{-0.060} & \textcolor{red}{-0.050} & \textcolor{red}{-0.059} & \textcolor{red}{-0.067} & \textcolor{red}{-0.073} \\
      CE & 0.021 & \textcolor{red}{-0.009} & 0.000 & \textcolor{red}{-0.042} & \textcolor{red}{-0.065} & \textcolor{red}{-0.077} & \textcolor{red}{-0.060} & \textcolor{red}{-0.028} & \textcolor{red}{-0.006} & \textcolor{red}{-0.029} & \textcolor{red}{-0.063} & \textcolor{red}{-0.149} \\
      DF & 0.012 & 0.002 & \textcolor{red}{-0.007} & \textcolor{red}{-0.029} & \textcolor{red}{-0.056} & \textcolor{red}{-0.098} & \textcolor{red}{-0.066} & \textcolor{red}{-0.036} & \textcolor{red}{-0.037} & \textcolor{red}{-0.059} & \textcolor{red}{-0.139} & \textcolor{red}{-0.165} \\
      ES & \textcolor{red}{-0.012} & \textcolor{red}{-0.049} & \textcolor{red}{-0.036} & \textcolor{red}{-0.030} & \textcolor{red}{-0.030} & \textcolor{red}{-0.034} & \textcolor{red}{-0.036} & \textcolor{red}{-0.035} & \textcolor{red}{-0.030} & \textcolor{red}{-0.031} & \textcolor{red}{-0.027} & \textcolor{red}{-0.033} \\
      GO & \textcolor{red}{-0.027} & \textcolor{red}{-0.053} & \textcolor{red}{-0.031} & \textcolor{red}{-0.029} & \textcolor{red}{-0.046} & \textcolor{red}{-0.056} & \textcolor{red}{-0.067} & \textcolor{red}{-0.057} & \textcolor{red}{-0.067} & \textcolor{red}{-0.166} & \textcolor{red}{-0.206} & \textcolor{red}{-0.256} \\
      MA & \textcolor{red}{-0.011} & \textcolor{red}{-0.046} & \textcolor{red}{-0.039} & \textcolor{red}{-0.043} & \textcolor{red}{-0.052} & \textcolor{red}{-0.059} & \textcolor{red}{-0.057} & \textcolor{red}{-0.068} & \textcolor{red}{-0.075} & \textcolor{red}{-0.130} & \textcolor{red}{-0.115} & \textcolor{red}{-0.105} \\
      MG & \textcolor{red}{-0.021} & \textcolor{red}{-0.040} & \textcolor{red}{-0.020} & \textcolor{red}{-0.026} & \textcolor{red}{-0.061} & \textcolor{red}{-0.092} & \textcolor{red}{-0.096} & \textcolor{red}{-0.097} & \textcolor{red}{-0.080} & \textcolor{red}{-0.106} & \textcolor{red}{-0.115} & \textcolor{red}{-0.137} \\
      MS & \textcolor{red}{-0.045} & \textcolor{red}{-0.090} & \textcolor{red}{-0.058} & \textcolor{red}{-0.072} & \textcolor{red}{-0.082} & \textcolor{red}{-0.121} & \textcolor{red}{-0.120} & \textcolor{red}{-0.123} & \textcolor{red}{-0.132} & \textcolor{red}{-0.178} & \textcolor{red}{-0.201} & \textcolor{red}{-0.206} \\
      MT & 0.005 & \textcolor{red}{-0.016} & \textcolor{red}{-0.013} & \textcolor{red}{-0.024} & \textcolor{red}{-0.033} & \textcolor{red}{-0.028} & \textcolor{red}{-0.017} & \textcolor{red}{-0.018} & \textcolor{red}{-0.035} & \textcolor{red}{-0.129} & \textcolor{red}{-0.109} & \textcolor{red}{-0.092} \\
      PA & 0.001 & \textcolor{red}{-0.037} & \textcolor{red}{-0.030} & \textcolor{red}{-0.073} & \textcolor{red}{-0.136} & \textcolor{red}{-0.174} & \textcolor{red}{-0.159} & \textcolor{red}{-0.123} & \textcolor{red}{-0.056} & \textcolor{red}{-0.051} & \textcolor{red}{-0.067} & \textcolor{red}{-0.107} \\
      PB & \textcolor{red}{-0.028} & \textcolor{red}{-0.007} & \textcolor{red}{-0.020} & \textcolor{red}{-0.023} & \textcolor{red}{-0.012} & \textcolor{red}{-0.046} & \textcolor{red}{-0.030} & 0.017 & \textcolor{red}{-0.020} & \textcolor{red}{-0.002} & 0.001 & 0.026 \\
      PE & 0.019 & 0.011 & 0.009 & \textcolor{red}{-0.016} & \textcolor{red}{-0.042} & \textcolor{red}{-0.065} & \textcolor{red}{-0.069} & \textcolor{red}{-0.036} & \textcolor{red}{-0.011} & \textcolor{red}{-0.005} & \textcolor{red}{-0.017} & \textcolor{red}{-0.053} \\
      PI & \textcolor{red}{-0.021} & \textcolor{red}{-0.040} & \textcolor{red}{-0.042} & \textcolor{red}{-0.048} & \textcolor{red}{-0.051} & \textcolor{red}{-0.055} & \textcolor{red}{-0.043} & \textcolor{red}{-0.032} & \textcolor{red}{-0.035} & \textcolor{red}{-0.056} & \textcolor{red}{-0.079} & \textcolor{red}{-0.080} \\
      PR & \textcolor{red}{-0.036} & \textcolor{red}{-0.059} & \textcolor{red}{-0.012} & \textcolor{red}{-0.003} & \textcolor{red}{-0.015} & \textcolor{red}{-0.062} & \textcolor{red}{-0.047} & \textcolor{red}{-0.033} & \textcolor{red}{-0.052} & \textcolor{red}{-0.111} & \textcolor{red}{-0.233} & \textcolor{red}{-0.362} \\
      RJ & 0.015 & \textcolor{red}{-0.033} & \textcolor{red}{-0.028} & \textcolor{red}{-0.047} & \textcolor{red}{-0.066} & \textcolor{red}{-0.091} & \textcolor{red}{-0.087} & \textcolor{red}{-0.070} & \textcolor{red}{-0.079} & \textcolor{red}{-0.089} & \textcolor{red}{-0.107} & \textcolor{red}{-0.122} \\
      RN & 0.004 & \textcolor{red}{-0.000} & \textcolor{red}{-0.018} & \textcolor{red}{-0.022} & \textcolor{red}{-0.028} & \textcolor{red}{-0.032} & \textcolor{red}{-0.024} & \textcolor{red}{-0.013} & \textcolor{red}{-0.024} & \textcolor{red}{-0.031} & \textcolor{red}{-0.038} & \textcolor{red}{-0.034} \\
      RO & \textcolor{red}{-0.017} & \textcolor{red}{-0.017} & \textcolor{red}{-0.016} & \textcolor{red}{-0.037} & \textcolor{red}{-0.065} & \textcolor{red}{-0.069} & \textcolor{red}{-0.058} & \textcolor{red}{-0.044} & \textcolor{red}{-0.044} & \textcolor{red}{-0.058} & \textcolor{red}{-0.071} & \textcolor{red}{-0.083} \\
      RR & 0.013 & \textcolor{red}{-0.040} & \textcolor{red}{-0.022} & \textcolor{red}{-0.056} & \textcolor{red}{-0.065} & \textcolor{red}{-0.078} & \textcolor{red}{-0.064} & \textcolor{red}{-0.037} & \textcolor{red}{-0.038} & \textcolor{red}{-0.024} & \textcolor{red}{-0.054} & \textcolor{red}{-0.048} \\
      RS & \textcolor{red}{-0.060} & \textcolor{red}{-0.109} & \textcolor{red}{-0.032} & \textcolor{red}{-0.013} & \textcolor{red}{-0.043} & \textcolor{red}{-0.127} & \textcolor{red}{-0.149} & \textcolor{red}{-0.148} & \textcolor{red}{-0.115} & \textcolor{red}{-0.217} & \textcolor{red}{-0.281} & \textcolor{red}{-0.286} \\
      SC & \textcolor{red}{-0.024} & \textcolor{red}{-0.022} & 0.040 & 0.052 & 0.035 & 0.006 & 0.027 & 0.036 & 0.030 & \textcolor{red}{-0.011} & \textcolor{red}{-0.032} & \textcolor{red}{-0.131} \\
      SE & \textcolor{red}{-0.011} & \textcolor{red}{-0.018} & \textcolor{red}{-0.004} & \textcolor{red}{-0.048} & \textcolor{red}{-0.060} & \textcolor{red}{-0.113} & \textcolor{red}{-0.097} & \textcolor{red}{-0.065} & \textcolor{red}{-0.035} & \textcolor{red}{-0.051} & \textcolor{red}{-0.091} & \textcolor{red}{-0.161} \\
      SP & \textcolor{red}{-0.043} & \textcolor{red}{-0.116} & \textcolor{red}{-0.053} & \textcolor{red}{-0.038} & \textcolor{red}{-0.028} & \textcolor{red}{-0.074} & \textcolor{red}{-0.068} & \textcolor{red}{-0.034} & \textcolor{red}{-0.037} & \textcolor{red}{-0.072} & \textcolor{red}{-0.142} & \textcolor{red}{-0.386} \\
      TO & \textcolor{red}{-0.008} & \textcolor{red}{-0.030} & \textcolor{red}{-0.015} & \textcolor{red}{-0.034} & \textcolor{red}{-0.058} & \textcolor{red}{-0.070} & \textcolor{red}{-0.055} & \textcolor{red}{-0.035} & \textcolor{red}{-0.026} & \textcolor{red}{-0.058} & \textcolor{red}{-0.080} & \textcolor{red}{-0.119} \\
      \hline
    \end{tabular}%
  }
  \label{tab:adm_var_sh}
\end{table}

\begin{table}[!htb]
  \caption{$\RelRMSE^{\Base}$ of dismissal series. VAR model for base forecasts and the shrinkage approach to estimate $\bm{W}$. Values in red indicate a $\RelRMSE^{\Base}$ less than 0.}
  \centering
  \resizebox{1\textwidth}{!}{%
    \begin{tabular}{lrrrrrrrrrrrr}
      \hline
      \multirow{2}{*}{Series} & \multicolumn{12}{c}{Forecast horizons ($h$)} \\ \cline{2-13}
      & 1 & 2 & 3 & 4 & 5 & 6 & 7 & 8 & 9 & 10 & 11 & 12 \\ \hline
      Total & 0.064 & 0.064 & 0.071 & 0.047 & 0.041 & 0.045 & 0.045 & 0.047 & 0.044 & 0.029 & 0.018 & 0.011 \\
      Midwest & \textcolor{red}{-0.023} & \textcolor{red}{-0.031} & \textcolor{red}{-0.041} & \textcolor{red}{-0.019} & \textcolor{red}{-0.013} & \textcolor{red}{-0.009} & \textcolor{red}{-0.012} & \textcolor{red}{-0.014} & \textcolor{red}{-0.009} & \textcolor{red}{-0.001} & 0.003 & 0.009 \\
      Northeast & \textcolor{red}{-0.007} & \textcolor{red}{-0.007} & \textcolor{red}{-0.010} & \textcolor{red}{-0.005} & \textcolor{red}{-0.003} & \textcolor{red}{-0.003} & \textcolor{red}{-0.004} & \textcolor{red}{-0.006} & \textcolor{red}{-0.002} & \textcolor{red}{-0.000} & 0.002 & 0.002 \\
      North & \textcolor{red}{-0.251} & \textcolor{red}{-0.228} & \textcolor{red}{-0.334} & \textcolor{red}{-0.192} & \textcolor{red}{-0.227} & \textcolor{red}{-0.129} & \textcolor{red}{-0.079} & \textcolor{red}{-0.059} & 0.036 & 0.100 & 0.066 & 0.059 \\
      Southeast & \textcolor{red}{-0.000} & \textcolor{red}{-0.001} & \textcolor{red}{-0.001} & 0.001 & \textcolor{red}{-0.000} & 0.000 & 0.000 & \textcolor{red}{-0.002} & \textcolor{red}{-0.001} & \textcolor{red}{-0.003} & \textcolor{red}{-0.003} & \textcolor{red}{-0.005} \\
      South & \textcolor{red}{-0.009} & \textcolor{red}{-0.013} & \textcolor{red}{-0.013} & \textcolor{red}{-0.005} & \textcolor{red}{-0.007} & \textcolor{red}{-0.009} & \textcolor{red}{-0.008} & \textcolor{red}{-0.012} & \textcolor{red}{-0.012} & \textcolor{red}{-0.016} & \textcolor{red}{-0.014} & \textcolor{red}{-0.030} \\
      AC & \textcolor{red}{-0.149} & \textcolor{red}{-0.068} & \textcolor{red}{-0.050} & 0.023 & 0.017 & \textcolor{red}{-0.064} & \textcolor{red}{-0.186} & \textcolor{red}{-0.276} & \textcolor{red}{-0.210} & \textcolor{red}{-0.083} & \textcolor{red}{-0.041} & 0.049 \\
      AL & 0.021 & 0.009 & \textcolor{red}{-0.005} & 0.013 & 0.014 & 0.021 & \textcolor{red}{-0.003} & \textcolor{red}{-0.011} & \textcolor{red}{-0.006} & \textcolor{red}{-0.003} & 0.002 & \textcolor{red}{-0.002} \\
      AM & \textcolor{red}{-0.142} & \textcolor{red}{-0.086} & \textcolor{red}{-0.103} & 0.010 & 0.063 & \textcolor{red}{-0.031} & \textcolor{red}{-0.028} & \textcolor{red}{-0.008} & 0.032 & 0.085 & 0.036 & 0.055 \\
      AP & \textcolor{red}{-3.329} & \textcolor{red}{-2.950} & \textcolor{red}{-2.216} & \textcolor{red}{-1.857} & \textcolor{red}{-3.051} & \textcolor{red}{-3.120} & \textcolor{red}{-3.847} & \textcolor{red}{-4.045} & \textcolor{red}{-4.269} & \textcolor{red}{-5.488} & \textcolor{red}{-5.674} & \textcolor{red}{-8.273} \\
      BA & \textcolor{red}{-0.052} & \textcolor{red}{-0.002} & 0.069 & 0.045 & 0.045 & 0.033 & 0.028 & 0.027 & 0.009 & 0.005 & 0.007 & \textcolor{red}{-0.001} \\
      CE & \textcolor{red}{-0.007} & \textcolor{red}{-0.024} & \textcolor{red}{-0.021} & \textcolor{red}{-0.025} & \textcolor{red}{-0.014} & \textcolor{red}{-0.012} & \textcolor{red}{-0.013} & \textcolor{red}{-0.023} & \textcolor{red}{-0.029} & \textcolor{red}{-0.044} & \textcolor{red}{-0.044} & \textcolor{red}{-0.034} \\
      DF & \textcolor{red}{-0.021} & \textcolor{red}{-0.031} & \textcolor{red}{-0.039} & \textcolor{red}{-0.008} & \textcolor{red}{-0.005} & \textcolor{red}{-0.013} & \textcolor{red}{-0.009} & \textcolor{red}{-0.011} & 0.001 & 0.009 & 0.009 & 0.003 \\
      ES & \textcolor{red}{-0.378} & \textcolor{red}{-0.365} & \textcolor{red}{-0.062} & \textcolor{red}{-0.011} & 0.064 & 0.066 & 0.061 & 0.126 & 0.171 & 0.297 & 0.373 & 0.413 \\
      GO & \textcolor{red}{-0.114} & \textcolor{red}{-0.084} & \textcolor{red}{-0.053} & \textcolor{red}{-0.101} & 0.049 & \textcolor{red}{-0.006} & 0.067 & 0.048 & 0.003 & 0.068 & 0.030 & 0.120 \\
      MA & \textcolor{red}{-0.156} & \textcolor{red}{-0.240} & \textcolor{red}{-0.244} & \textcolor{red}{-0.042} & \textcolor{red}{-0.154} & 0.011 & 0.038 & \textcolor{red}{-0.068} & \textcolor{red}{-0.094} & \textcolor{red}{-0.178} & \textcolor{red}{-0.162} & \textcolor{red}{-0.276} \\
      MG & \textcolor{red}{-0.092} & \textcolor{red}{-0.011} & 0.046 & 0.063 & 0.123 & 0.099 & 0.090 & 0.093 & 0.080 & 0.114 & 0.117 & 0.138 \\
      MS & 0.121 & 0.069 & 0.082 & 0.051 & 0.008 & 0.002 & \textcolor{red}{-0.002} & \textcolor{red}{-0.007} & \textcolor{red}{-0.024} & \textcolor{red}{-0.063} & \textcolor{red}{-0.082} & \textcolor{red}{-0.112} \\
      MT & \textcolor{red}{-0.022} & \textcolor{red}{-0.003} & \textcolor{red}{-0.001} & \textcolor{red}{-0.014} & 0.018 & 0.010 & 0.034 & 0.022 & 0.004 & 0.022 & 0.032 & 0.072 \\
      PA & \textcolor{red}{-0.021} & \textcolor{red}{-0.007} & \textcolor{red}{-0.012} & \textcolor{red}{-0.010} & \textcolor{red}{-0.013} & \textcolor{red}{-0.016} & \textcolor{red}{-0.009} & 0.003 & 0.009 & 0.011 & 0.006 & 0.015 \\
      PB & \textcolor{red}{-0.001} & 0.009 & 0.010 & 0.012 & 0.015 & 0.005 & \textcolor{red}{-0.003} & \textcolor{red}{-0.008} & \textcolor{red}{-0.002} & \textcolor{red}{-0.018} & \textcolor{red}{-0.027} & \textcolor{red}{-0.042} \\
      PE & 0.001 & 0.000 & \textcolor{red}{-0.000} & 0.003 & 0.002 & \textcolor{red}{-0.001} & \textcolor{red}{-0.006} & \textcolor{red}{-0.008} & \textcolor{red}{-0.005} & \textcolor{red}{-0.007} & \textcolor{red}{-0.008} & \textcolor{red}{-0.011} \\
      PI & \textcolor{red}{-0.005} & \textcolor{red}{-0.007} & \textcolor{red}{-0.014} & \textcolor{red}{-0.012} & \textcolor{red}{-0.005} & \textcolor{red}{-0.005} & 0.003 & 0.011 & 0.024 & 0.020 & 0.013 & 0.009 \\
      PR & 0.089 & 0.092 & 0.111 & 0.059 & 0.030 & 0.066 & 0.061 & 0.089 & 0.107 & 0.121 & 0.117 & 0.094 \\
      RJ & \textcolor{red}{-0.014} & \textcolor{red}{-0.033} & \textcolor{red}{-0.005} & \textcolor{red}{-0.029} & 0.034 & 0.057 & 0.050 & 0.093 & 0.139 & 0.226 & 0.267 & 0.292 \\
      RN & \textcolor{red}{-1.064} & \textcolor{red}{-0.841} & \textcolor{red}{-0.316} & \textcolor{red}{-0.450} & \textcolor{red}{-0.430} & \textcolor{red}{-0.352} & \textcolor{red}{-0.574} & \textcolor{red}{-0.397} & \textcolor{red}{-0.423} & \textcolor{red}{-0.144} & 0.059 & \textcolor{red}{-0.060} \\
      RO & \textcolor{red}{-2.480} & \textcolor{red}{-2.298} & \textcolor{red}{-2.047} & \textcolor{red}{-1.828} & \textcolor{red}{-1.606} & \textcolor{red}{-1.433} & \textcolor{red}{-1.274} & \textcolor{red}{-1.313} & \textcolor{red}{-1.267} & \textcolor{red}{-0.959} & \textcolor{red}{-0.746} & \textcolor{red}{-0.449} \\
      RR & \textcolor{red}{-0.145} & \textcolor{red}{-0.158} & \textcolor{red}{-0.090} & \textcolor{red}{-0.072} & \textcolor{red}{-0.230} & \textcolor{red}{-0.245} & \textcolor{red}{-0.568} & \textcolor{red}{-1.269} & \textcolor{red}{-1.951} & \textcolor{red}{-1.756} & \textcolor{red}{-1.387} & \textcolor{red}{-1.214} \\
      RS & \textcolor{red}{-0.009} & \textcolor{red}{-0.012} & \textcolor{red}{-0.012} & \textcolor{red}{-0.005} & \textcolor{red}{-0.007} & \textcolor{red}{-0.008} & \textcolor{red}{-0.008} & \textcolor{red}{-0.010} & \textcolor{red}{-0.008} & \textcolor{red}{-0.003} & 0.001 & 0.006 \\
      SC & \textcolor{red}{-0.210} & \textcolor{red}{-0.289} & \textcolor{red}{-0.216} & \textcolor{red}{-0.086} & \textcolor{red}{-0.084} & \textcolor{red}{-0.115} & \textcolor{red}{-0.150} & \textcolor{red}{-0.170} & \textcolor{red}{-0.182} & \textcolor{red}{-0.248} & \textcolor{red}{-0.194} & \textcolor{red}{-0.226} \\
      SE & \textcolor{red}{-1.712} & \textcolor{red}{-3.290} & \textcolor{red}{-6.043} & \textcolor{red}{-6.216} & \textcolor{red}{-5.338} & \textcolor{red}{-7.419} & \textcolor{red}{-8.615} & \textcolor{red}{-8.109} & \textcolor{red}{-7.832} & \textcolor{red}{-7.931} & \textcolor{red}{-6.458} & \textcolor{red}{-6.773} \\
      SP & \textcolor{red}{-0.004} & \textcolor{red}{-0.003} & \textcolor{red}{-0.003} & \textcolor{red}{-0.002} & \textcolor{red}{-0.002} & \textcolor{red}{-0.002} & \textcolor{red}{-0.003} & \textcolor{red}{-0.002} & \textcolor{red}{-0.002} & \textcolor{red}{-0.002} & \textcolor{red}{-0.001} & \textcolor{red}{-0.002} \\
      TO & 0.002 & \textcolor{red}{-0.004} & 0.009 & 0.038 & \textcolor{red}{-0.008} & 0.010 & 0.008 & \textcolor{red}{-0.022} & \textcolor{red}{-0.022} & \textcolor{red}{-0.046} & \textcolor{red}{-0.036} & \textcolor{red}{-0.078} \\
      \hline
    \end{tabular}%
  }
  \label{tab:dem_var_sh}
\end{table}


\end{document}
